% concatenated from main.tex
\pdfoutput=1 % ensure arXiv direct PDF output
\documentclass[11pt]{amsart}
\usepackage{amssymb,amsthm,enumitem,colonequals,tikz-cd,microtype,rotating}
\usepackage[normalem]{ulem} 

%% Font style: old versions mlmodern, \usepackage[cal=euler,bfcal,bb=px,bfbb]{mathalpha}
\usepackage[osf]{Baskervaldx}
\usepackage[baskervaldx]{newtxmath}
\usepackage[cal=boondoxo]{mathalfa} % mathcal from STIX, unslanted a bit

%% Geometry with custom margins
\usepackage[top=3.75cm, bottom=3cm, left=3.5cm, right=3.5cm]{geometry}

%% High penalty to avoid footnote breaks
\interfootnotelinepenalty=10000 

%% Color definitions for links and references
\usepackage{xcolor}
\colorlet{darkblue}{blue!55!black}
\colorlet{darkcyan}{cyan!50!black}
\colorlet{darkgreen}{green!60!black}

%% URL hyphenation option must come before hyperref
\PassOptionsToPackage{hyphens}{url}
\usepackage{hyperref}
\hypersetup{
  colorlinks=true,
  linkcolor=darkblue,
  urlcolor=darkcyan,
  citecolor=darkgreen,
}

%% Colored eqref for emphasis (optional)
\def\eqref#1{\textcolor{darkblue}{(\ref{#1})}}

%% Cleveref utility
\usepackage[nameinlink]{cleveref} 
\Crefformat{section}{#2\S#1#3}
\Crefmultiformat{section}{#2\S\S#1#3}{ and~#2#1#3}{, #2#1#3}{, and~#2#1#3}
\Crefname{equation}{Diagram}{Diagrams}

%% Overfill notice
\usepackage[pagewise]{lineno}
\overfullrule=100pt % for debugging, comment or set to 0 for final

\let\oldequation\equation
\let\oldendequation\endequation
\renewenvironment{equation}{\linenomathNonumbers\oldequation}{\oldendequation\endlinenomath}
\expandafter\let\expandafter\oldequationstar\csname equation*\endcsname
\expandafter\let\expandafter\oldendequationstar\csname endequation*\endcsname
\renewenvironment{equation*}{\linenomathNonumbers\oldequationstar}{\oldendequationstar\endlinenomath}
\let\oldalign\align
\let\oldendalign\endalign
\renewenvironment{align}{\linenomathNonumbers\oldalign}{\oldendalign\endlinenomath}
\expandafter\let\expandafter\oldalignstar\csname align*\endcsname
\expandafter\let\expandafter\oldendalignstar\csname endalign*\endcsname
\renewenvironment{align*}{\linenomathNonumbers\oldalignstar}{\oldendalignstar\endlinenomath}

%% Theorem and counter setup
\theoremstyle{plain}
\newtheorem{theorem}{Theorem}[section]
\newtheorem{lemma}[theorem]{Lemma}
\newtheorem{corollary}[theorem]{Corollary}
\newtheorem{proposition}[theorem]{Proposition}

\theoremstyle{definition}
\newtheorem{definition}[theorem]{Definition}
\newtheorem{example}[theorem]{Example}
\newtheorem{remark}[theorem]{Remark}
\newtheorem{setup}[theorem]{Setup}
\newtheorem{terminology}[theorem]{Terminology}

\newtheorem*{disclaimer}{Disclaimer}
\newtheorem*{ack}{Acknowledgments}

% declare alias for cleveref, otherwise it doesn't do what it should; this probably isn't best fix but found it on the web
% \AddToHook{env/introthm/begin}{\crefalias{intro}{theorem}}
\AddToHook{env/conjecture/begin}{\crefalias{theorem}{conjecture}}
\AddToHook{env/lemma/begin}{\crefalias{theorem}{lemma}}
\AddToHook{env/corollary/begin}{\crefalias{theorem}{corollary}}
\AddToHook{env/proposition/begin}{\crefalias{theorem}{proposition}}
\AddToHook{env/definition/begin}{\crefalias{theorem}{definition}}
\AddToHook{env/remark/begin}{\crefalias{theorem}{remark}}
\AddToHook{env/setup/begin}{\crefalias{theorem}{setup}}
\AddToHook{env/example/begin}{\crefalias{theorem}{example}}
\AddToHook{env/terminology/begin}{\crefalias{theorem}{terminology}}

%% toc and numbering depths
\setcounter{tocdepth}{1}
\setcounter{secnumdepth}{2}
\numberwithin{equation}{section}
\numberwithin{theorem}{section}

%% Comment out collaborator notes before final version
\newcommand{\Pat}[1]{{\color{blue} $\clubsuit$ Pat: [#1]}}
\newcommand{\Kabeer}[1]{{\color{purple} $\clubsuit$ Kabeer: [#1]}}
\newcommand{\Neb}[1]{{\color{red} $\clubsuit$ Neb: [#1]}}
\newcommand{\Elias}[1]{{\color{cyan} $\clubsuit$ El\'{i}as: [#1]}}

%% Document information
\title[Fiberwise criteria for Fourier--Mukai equivalences]{Fiberwise criteria for Fourier--Mukai equivalences}

\author[E.~Guisado Villalgordo]{El\'{i}as Guisado Villalgordo}
\address{E.~Guisado Villalgordo,
BCAM -- Basque Center for Applied Mathematics, Mazarredo 14, 48009
Bilbao, Basque Country -- Spain}
\email{eguisado@bcamath.org}

\author[P.~Lank]{Pat Lank}
\address{P.~Lank,
Dipartimento di Matematica “F. Enriques”, Universit\`{a} degli Studi di Milano, Via Cesare Saldini 50, 20133 Milano, Italy}
\email{plankmathematics@gmail.com}

\author[K.~Manali Rahul]{Kabeer Manali Rahul}
\address{K.~Manali Rahul,
Max Planck Institute for Mathematics,
Bonn, Germany}
\email{kabeermr.maths@gmail.com}

\author[N.~Pavic]{Nebojsa Pavic}
\address{N.~Pavic,
Institut f\"{u}r Mathematik und Wissenschaftliches Rechnen, 
Universit\"{a}t Graz,
Graz, Austria}
\email{nebojsa.pavic@uni-graz.at}

\date{\today}

\keywords{Fourier--Mukai partners, integral transforms, singular varieties, derived categories, fully faithfulness}

\subjclass[2020]{14A30 (primary), 14F08, 18G80, 14B05}

%14A30 Fundamental constructions in algebraic geometry involving higher and derived categories (homotopical algebraic geometry, derived algebraic geometry, etc.)

%14F08 Derived categories of sheaves, dg categories, and related constructions in algebraic geometry 

%18G80 Derived categories, triangulated categories

%14B05 Singularities in algebraic geometry

\begin{document}
    
\begin{abstract}
    We study the behavior of integral transforms under base change. In particular, we establish a yoga of local algebra and fibers to test for derived equivalences or fully faithfulness via integral transforms. This generalizes a result of Orlov to singular varieties and strengthens several results in the literature by allowing arbitrary base fields. Additionally, it provides new insight into fibrations and their singularities in arithmetic settings (e.g.\ projective and flat schemes over a DVR).
\end{abstract}

\maketitle

\tableofcontents

%%%%%%%%%%%%%%%%%%%%%%%%%%%%%%%%%%%%%
\section{Introduction}
\label{sec:intro}
%%%%%%%%%%%%%%%%%%%%%%%%%%%%%%%%%%%%%

%%%%%%%%%%%%%%%%%%%%%%%%%%%%%%%%%%%%% 
\subsection{What is known}
\label{sec:intro_what_is_known}
%%%%%%%%%%%%%%%%%%%%%%%%%%%%%%%%%%%%%

The context for our work begins with Mukai \cite{Mukai:1981}. In that work, integral transforms were introduced to study triangulated equivalences between derived categories of abelian varieties. Since then, integral transforms have played a central role in the interplay between derived categories and algebraic geometry (e.g.\ \cite{Bridgeland:2002, Kawamata:2002a}). 

Much of this progress concerns smooth varieties, where geometric results can be obtained using categorical methods. However, extending these methods to singular varieties remains challenging, largely because bounded coherent complexes need not be perfect. Thus, an important task is to develop techniques for studying integral transforms in singular settings.

Recent work has begun to address these difficulties. Criteria determining when an integral transform is fully faithful or an equivalence have been established for varieties over algebraically closed fields, generalizing results from the smooth case \cite{Ruiperez/Hernandez/Martin/SanchodeSalas:2007, Ruiperez/Hernandez/Martin/SanchodeSalas:2009} and building upon foundational work of Bondal--Orlov \cite{Bondal/Orlov:1995}. Likewise, Bondal--Orlov's reconstruction theorem for smooth projective varieties has been extended to projective curves \cite{Spence:2023} and (anti-)Fano Gorenstein varieties \cite{Ballard:2011}.

There is also growing interest in studying integral transforms in the relative setting. Specifically, let $Y_1$ and $Y_2$ be proper $S$-schemes, where $S$ is Noetherian. An \textit{integral transform} is an exact functor 
\begin{displaymath}
    \Phi_K := \mathbf{R}(p_2)_\ast(\mathbf{L}p_1^\ast(-)\otimes^{\mathbf{L}} K) \colon D_{\operatorname{qc}}(Y_1) \to D_{\operatorname{qc}}(Y_2)
\end{displaymath}
where $K \in D_{\operatorname{qc}}(Y_1\times_S Y_2)$ and $p_i \colon Y_1\times_S Y_2 \to Y_i$ are the projection morphisms.
If there is $K$ such that $\Phi_K$ is an equivalence, we say $Y_1$ and $Y_2$ are \textit{Fourier--Mukai partners (over $S$)}.
The complex $K$ is called the \emph{kernel} of the integral transform $\Phi_K$.

Such transforms naturally arise in the study of \textit{fibrations}; that is, in situations involving proper and flat morphisms. See e.g.\ \cite{Ruiperez/Hernandez/Martin/SanchodeSalas:2007, Ruiperez/Hernandez/Martin/SanchodeSalas:2009}, which dealt with families whose special fibers have mild singularities. In particular, loc.\ cit.\ studied when an integral transform between \textit{locally projective} fibrations $Y_i \to S$ is an equivalence for $S$ of finite type over an algebraically closed field. It was shown that the induced integral transform on special fibers (i.e.\ closed points) over $S$ controls being an equivalence or fully faithful.

%%%%%%%%%%%%%%%%%%%%%%%%%%%%%%%%%%%%%
\subsection{What we do}
\label{sec:intro_what_we_do}
%%%%%%%%%%%%%%%%%%%%%%%%%%%%%%%%%%%%%

In this work, we deepen this perspective by systematically exploiting special fibers. Specifically, we study the behavior of integral transforms under base change using techniques from local algebra.

%%%%%%%%%%%%%%%%%%%%%%%%%%%%%%%%%%%%%
\subsubsection{Preservation of adjointness}
\label{sec:intro_what_we_do_adjointness}
%%%%%%%%%%%%%%%%%%%%%%%%%%%%%%%%%%%%%

Neeman has developed a new framework called approximable triangulated categories (see e.g.\ \cite{Canonaco/Neeman/Stellari:2025,Neeman:2023}), which has resolved open conjectures in algebraic geometry \cite{Neeman:2021a,Neeman:2022}. As a consequence, any integral transform $\Phi_K$ inducing an equivalence on $D_{\operatorname{qc}}$ restricts to equivalences on $D^b_{\operatorname{coh}}$ and $\operatorname{Perf}$, see \Cref{lem:equivalences_induced}.

Consequently, questions about equivalences may be reduced to the behavior of integral transforms on these smaller subcategories. Whether an integral transform restricts to these smaller categories is now understood \cite{Rizzardo:2017,Ballard:2009,Dutta/Lank/ManaliRahul:2025}. These results shed light on the existence of left or right adjoints to integral transforms and the conditions under which such adjoints also restrict to these subcategories.

However, the behavior of adjointness under base change is not yet fully understood. Our focus is on base changes induced by certain morphisms. This leads us to the following.

\begin{theorem}
    \label{thm:intro_preservation_adjoint}
    Let $S$ be a Noetherian scheme. Suppose $Y_1$ and $Y_2$ are proper and flat $S$-schemes. Consider $K\in D^b_{\operatorname{coh}}(Y_1\times_S Y_2)$. Then the following are equivalent:
    \begin{enumerate}
        \item $K$ is relatively perfect over $Y_1$ (resp.\ $Y_2$)
        \item $\mathbf{L}(t^\prime)^\ast K$ is relatively perfect over $Y_1\times_S T$ (resp.\ $Y_2\times_S T$) for any affine morphism $t$ from a Noetherian scheme
        \item $\mathbf{L}(t^\prime)^\ast K$ is relatively perfect over $Y_1\times_S \operatorname{Spec}(k)$ (resp.\ $Y_2\times_S \operatorname{Spec}(k)$) for all morphisms $t\colon \operatorname{Spec}(k)\to S$ from a field
        \item $\mathbf{L}(t^\prime)^\ast K$ is relatively perfect over $Y_1\times_S \operatorname{Spec}(\kappa(p))$ (resp.\ $Y_2\times_S \operatorname{Spec}(\kappa(p))$) for every closed point $p\in S$ with associated closed immersion $t\colon \operatorname{Spec}(\kappa(p))\to S$
        \item $\mathbf{L}(t^\prime)^\ast K$ is relatively perfect over $Y_1\times_S T$ (resp.\ $Y_2\times_S T$) for some affine surjective morphism $t\colon T\to S$.
    \end{enumerate}
    Here, $t^\prime\colon Y_1\times_S Y_2 \times_S T\to Y_1 \times_S Y_2$ is the natural morphism. In fact, in either case the corresponding left or right adjoints are preserved.
\end{theorem}

This follows from \Cref{thm:bounded_pseudocoherence_perfectness_faithfully_flat_affine} together with \Cref{prop:right_adjoint_pullback,prop:left_adjoint_pullback}. In short, \Cref{thm:intro_preservation_adjoint} describes how integral transforms behave under base change when restricted to $\operatorname{Perf}$ or $D^b_{\operatorname{coh}}$. Importantly, it addresses the preservation of adjoints under base change. To the best of our knowledge, this result is new. In particular, it applies in the presence of singularities on the $Y_i$. An important instance where the result is applicable is $t$ arising from field extensions.

%%%%%%%%%%%%%%%%%%%%%%%%%%%%%%%%%%%%%
\subsubsection{Fully faithfulness \& equivalences}
\label{sec:intro_what_we_dofully_faithfulness_equivalence}
%%%%%%%%%%%%%%%%%%%%%%%%%%%%%%%%%%%%%

Next, we shift gears to studying fully faithfulness and equivalences under base change. Specifically, the ascent and descent of these properties. Before doing so, we briefly review the relevant literature.

As mentioned earlier, \cite{Ruiperez/Hernandez/Martin/SanchodeSalas:2007,Ruiperez/Hernandez/Martin/SanchodeSalas:2009} studied these questions for fibrations whose base is of finite type over an algebraically closed field. A key ingredient in their work was to look at the behavior on special fibers. More recently, such questions have also been studied for smooth proper morphisms over a Noetherian base scheme (see \cite[Proposition/Definition 2.3]{Kurama:2024}). However, smoothness assumptions restrict the range of examples one can consider, e.g.\ fibrations admitting singular fibers.

A result of Orlov addresses similar questions in the case of smooth varieties over a field. Specifically, for smooth projective varieties over a field, an integral transform is fully faithful or an equivalence if, and only if, the induced integral transform satisfies the same property after any base field extension \cite[Lemma 2.12]{Orlov:2002}. The strategy of loc.\ cit.\ uses `canonical morphisms' between the kernels of integral transforms that describe the (co)units of adjunction. However, it remains unclear whether such canonical morphisms exist in the singular setting. This uncertainty provides further motivation for our work: we bypass the absence of such tools in the singular setting by developing different techniques.

This brings attention to the following.

\begin{theorem}
    \label{introthm:fully_faithfulness_equivalence}
    Let $S$ be a Noetherian scheme. Suppose $Y_1$ and $Y_2$ are proper and flat $S$-schemes. Consider $K\in D^b_{\operatorname{coh}}(Y_1\times_S Y_2)$ which is relatively perfect over each $Y_i$. Then the following are equivalent:
    \begin{enumerate}
        \item $\Phi_{K}$ is fully faithful (resp.\ an equivalence) on $D^b_{\operatorname{coh}}$
        \item $\Phi_{\mathbf{L}(t^\prime)^\ast K}$ is fully faithful (resp.\ an equivalence) on $D^b_{\operatorname{coh}}$ for any affine morphism $t$
        \item $\Phi_{\mathbf{L}(t^\prime)^\ast K}$ is fully faithful (resp.\ an equivalence) on $D^b_{\operatorname{coh}}$ for every morphism $t\colon \operatorname{Spec}(k)\to S$ from a field
        \item $\Phi_{\mathbf{L}(t^\prime)^\ast K}$ is fully faithful (resp.\ an equivalence) on $D^b_{\operatorname{coh}}$ for every closed point $p\in S$ with associated closed immersion $t\colon \operatorname{Spec}(\kappa(p))\to S$
        \item $\Phi_{\mathbf{L}(t^\prime)^\ast K}$ is fully faithful (resp.\ an equivalence) on $D^b_{\operatorname{coh}}$ for some affine surjection $t$.
    \end{enumerate}
    Here, $t^\prime\colon Y_1\times_S Y_2 \times_S T\to Y_1 \times_S Y_2$ is the natural morphism induced by $t\colon T \to S$. 
\end{theorem}

This appears later as \Cref{thm:descent_ascent}. Notably, it generalizes \cite[Proposition 2.15]{Ruiperez/Hernandez/Martin/SanchodeSalas:2009}, which only proved the statement for closed points. As a consequence, this allows us to generalize Orlov's result to the singular setting (cf.\ \cite[Lemma 2.12]{Orlov:2002}).

\begin{corollary}
    \label{cor:orlov_varieties}
    Consider proper schemes $Y_1$ and $Y_2$ over a field $k$.
    Let $K\in D^b_{\operatorname{coh}}(Y_1\times_k Y_2)$ be relatively perfect over each $Y_i$.
    Let $L/k$ be a field extension.
    Then $\Phi_{K_L}$ is fully faithful (resp.\ an equivalence) on $D^b_{\operatorname{coh}}$
    if and only if $\Phi_{K}$ is fully faithful (resp.\ an equivalence) on $D^b_{\operatorname{coh}}$.
\end{corollary}

Our proofs require careful bookkeeping of how the (co)units of adjunction for integral transforms behave under base change. See \Cref{sec:fully_faithful_equivalence} for details. Notably, the techniques we use are independent of those appearing in \cite[Lemma 2.12]{Orlov:2002}.

We compare the methodology to that of \cite{Bergh/Schnurer:2020}. 
Our approach avoids representing (co)units by morphisms between perfect kernels.
Although mates appear in \cite{Bergh/Schnurer:2020}, their argument crucially requires perfectness. 
In particular, the proof of \cite[Proposition 4.9]{Bergh/Schnurer:2020} does not hold for nonperfect kernels, and so the base change machinery we develop develop is necessary. 

In fact, the proof of \cite[Proposition 2.15]{Ruiperez/Hernandez/Martin/SanchodeSalas:2009} led us to approach this problem with mates. 
Loc.\ cit.\ had used that units are functorial with respect to derived pushforward along a closed immersion.
To us, this seemed like a very nontrivial thing, nor could we find a reference for the fact. 
To address this concern, we have rigorously filled all of the details, and more, with mates. 
See e.g.\ \Cref{prop:pullforwards_is_a_morphism_of_adjoint_integral_transforms}.
However, we argue in a way that is fairly self-contained, and does not mimic the approach of \cite{Ruiperez/Hernandez/Martin/SanchodeSalas:2009} via spectral sequences (see proof of \cite[Proposition 2.15]{Ruiperez/Hernandez/Martin/SanchodeSalas:2009}).

%%%%%%%%%%%%%%%%%%%%%%%%%%%%%%%%%%%%%
\subsection{What is next}
\label{sec:intro_what_next}
%%%%%%%%%%%%%%%%%%%%%%%%%%%%%%%%%%%%%

We conclude by discussing several consequences of the results above. Our results extend the literature on schemes of finite type over algebraically closed fields to the more general setting of arbitrary fields. They reveal a new interplay between the intuition of \cite{Ruiperez/Hernandez/Martin/SanchodeSalas:2007,Ruiperez/Hernandez/Martin/SanchodeSalas:2009} concerning special fibers and faithfully flat descent along affine morphisms. In particular, certain singularities are shown to be derived invariants for fibrations over Noetherian schemes. This broadens the scope to include arithmetically flavored situations of potential interest for future research.

For a pair of geometrically integral projective curves over a field $k$ of characteristic zero, we show Fourier--Mukai partnership implies they are $k$-forms. Particularly: 

\begin{proposition}
    \label{prop:curves_isomorphic_after_base_change}
    Let $k$ be a field of characteristic zero. Suppose $C$ and $C^\prime$ are geometrically integral projective curves over $k$. If $C$ and $C^\prime$ are Fourier--Mukai partners over $k$, then $\operatorname{Spec}(L)\times_k C \cong \operatorname{Spec}(L)\times_k C^\prime$ for some finite field extension $L/k$ (i.e.\ are $k$-forms).
\end{proposition}

Here, a $k$-form means there exists a finite extension $L/k$ such that their base changes along $L/k$ become isomorphic. The result holds for singular curves and extends the algebraically closed case of \cite{Spence:2023}. For fibrations of relative dimension one over a Noetherian $\mathbb{Q}$-scheme, Fourier--Mukai partnership forces the genera of special fibers to coincide (see \Cref{cor:special_fibers_genii_invariance}). It would be interesting to find a prime or mixed characteristic analog of these results.

Next, we show that smoothness for proper schemes over an arbitrary field is likewise invariant under Fourier--Mukai partnership (see \Cref{lem:smoothness_derived_invariance}). This subtlety arises especially over imperfect fields. Since regularity is characterized by $D^b_{\operatorname{coh}} = \operatorname{Perf}$, one might expect this condition to persist after a change of field extensions, but there exist regular projective curves that become singular after such a base change (see e.g.\ \cite[Remark 16]{Kollar:2011}).

In the setting of fibrations, smoothness of the total space is also a derived invariant:

\begin{proposition}
    \label{prop:smooth_morphism_derived_invariance}
    Let $f_i \colon Y_i \to S$ be proper flat morphisms to a Noetherian scheme. Assume that $Y_1$ and $Y_2$ are Fourier--Mukai partners over $S$. Then $f_1$ is smooth if, and only if, $f_2$ is smooth.
\end{proposition}

Finally, generalizing \cite{Ruiperez/Hernandez/Martin/SanchodeSalas:2009}, we show that the Cohen--Macaulay and Gorenstein properties of a morphism remain invariant under Fourier--Mukai partnership for Noetherian base schemes. Specifically:

\begin{proposition}
    \label{prop:CM_or_Gorensteinness_morphism_derived_invariance}
    Let $f_i \colon Y_i \to S$ be projective flat morphisms with geometrically integral fibers to a Noetherian scheme. Assume that $Y_1$ and $Y_2$ are Fourier--Mukai partners over $S$. Then $f_1$ is Cohen--Macaulay (resp.\ Gorenstein) if, and only if, $f_2$ satisfies the same condition.
\end{proposition}

%%%%%%%%%%%%%%%%%%%%%%%%%%%%%%%%
\subsection*{Notation}
\label{sec:intro_notation}
%%%%%%%%%%%%%%%%%%%%%%%%%%%%%%%%

Let $X$ be a scheme. We work with the following triangulated categories:
\begin{enumerate}
    \item $D(X):=D(\operatorname{Mod}(X))$ is the derived category of $\mathcal{O}_X$-modules.
    \item $D_{\operatorname{qc}}(X)$ is the (strictly full) subcategory of $D(X)$ consisting of complexes with quasi-coherent cohomology.
    \item If $X$ is Noetherian, then $D_{\operatorname{coh}}^b(X)$ (resp.\ $D_{\operatorname{qc}}^b(X)$) is the (strictly full) subcategory of $D(X)$ consisting of complexes having bounded and coherent (resp.\ quasi-coherent) cohomology.
    \item $\operatorname{Perf}(X)$ is the (strictly full) subcategory of $D_{\operatorname{qc}}(X)$ consisting of the perfect complexes on $X$.
\end{enumerate}
At times if $X$ is affine, we might abuse notation and write $D^b_{\operatorname{coh}}(R):=D^b_{\operatorname{coh}}(X)$ where $R:=H^0(X,\mathcal{O}_X)$ is the ring of global sections; similar conventions will occur for the other categories mentioned here and later.

\begin{ack}
    Lank was supported under the ERC Advanced Grant 101095900-TriCatApp. Part of this work was completed while Lank visited the Basque Center for Applied Mathematics (BCAM) and thanks them for their hospitality. Guisado Villalgordo was supported under the grant PRE2022-104796 by the Spanish Ministry of Science and Innovation. Manali Rahul is grateful to Max Planck Institute for Mathematics in
    Bonn for its hospitality and financial support. The authors thank Dylan Spence, Ana Cristina L\'{o}pez Mart\'{i}n, and Fernando Sancho de Salas for discussions.
\end{ack}

%%%%%%%%%%%%%%%%%%%%%%%%%%%%%%%%%%%%%
\section{Preliminaries}
\label{sec:prelim}
%%%%%%%%%%%%%%%%%%%%%%%%%%%%%%%%%%%%%

Let $X$ be a Noetherian scheme.

%%%%%%%%%%%%%%%%%%%%%%%%%%%%%%%%
\subsection{Perfect \& pseudocoherent complexes}
\label{sec:prelim_perfectness_pseudocoherence}
%%%%%%%%%%%%%%%%%%%%%%%%%%%%%%%%

We say $P\in D_{\operatorname{qc}}(X)$ is \textbf{perfect} if it is locally isomorphic to a bounded complex of locally free sheaves of finite rank. Denote by $\operatorname{Perf}(X)$ the full subcategory of perfect complexes. Also, given an integer $N$, we say $E\in D_{\operatorname{qc}}(X)$ is $N$-\textbf{pseudocoherent} if locally there is a morphism $P \to E$ from some $P\in \operatorname{Perf}(X)$ such that the induced morphism on cohomology sheaves $\mathcal{H}^n (P) \to \mathcal{H}^n (E)$ is an isomorphism for $n>N$ and is surjective for $n=N$.
We say $E\in D_{\operatorname{qc}}(X)$ is \textbf{pseudocoherent} if it is $N$-pseudocoherent for every $N\in\mathbb{Z}$.

In general, $E\in D_{\operatorname{qc}}(X)$ is pseudocoherent if, and only if, locally on $X$ there is a quasi-isomorphism $P\to E$, with $P$ a bounded above complex of finite-rank free modules \cite[Corollary 22.47]{Gortz/Wedhorn:2023}.
As $X$ is Noetherian, $E\in D_{\operatorname{qc}}(X)$ is pseudocoherent if, and only if, $\mathcal{H}^n(E)$ is coherent for all $n$ and vanishes for $n\gg 0$ (see e.g.\ \cite[\href{https://stacks.math.columbia.edu/tag/08E8}{Tag 08E8}]{stacks-project}). The subcategory of pseudocoherent complexes of $D_{\operatorname{qc}}(X)$ coincides with $D^-_{\operatorname{coh}}(X)$ (i.e.\ the full subcategory of complexes with bounded above and coherent cohomology).

Consider a morphism $f\colon Y \to X$ of Noetherian schemes. 
Denote the right adjoint of $\mathbf{R}f_\ast$ on $D_{\operatorname{qc}}$ by $f^\times$ \cite[\href{https://stacks.math.columbia.edu/tag/0A9E}{Tag 0A9E}]{stacks-project}.
If $f$ is separated and of finite type, then $f^!\colon D^+_{\operatorname{qc}}(X) \to D^+_{\operatorname{qc}}(Y)$ denotes the so-called upper shriek functor \cite[\href{https://stacks.math.columbia.edu/tag/0A9Y}{Tag 0A9Y}]{stacks-project}. 
Then one has $\mathbf{L}f^\ast D^-_{\operatorname{coh}}(X)\subseteq D^-_{\operatorname{coh}}(Y)$ (see e.g.\ \cite[\href{https://stacks.math.columbia.edu/tag/08H4}{Tag 08H4}]{stacks-project}).
%%NOTE: See also 09U7 if needed

We record a few straightforward results.

\begin{lemma}
    \label{lem:pseudocoherence_affine_faithfully_flat_cover}
    Let $f\colon Y \to X$ be a faithfully flat morphism of Noetherian schemes. An object $E\in D_{\operatorname{qc}}(X)$ belongs to $D^-_{\operatorname{coh}}(X)$ (resp.\ $\operatorname{Perf}(X)$) if, and only if, $\mathbf{L}f^\ast E\in D^-_{\operatorname{coh}}(Y)$ (resp.\ $\mathbf{L}f^\ast E\in \operatorname{Perf}(Y)$).
\end{lemma}

\begin{proof}
    This is a special case of \cite[Proposition 22.52]{Gortz/Wedhorn:2023}. 
\end{proof}

Denote by $\mathbf{R}\operatorname{\mathcal{H}\! \mathit{om}}$ the internal Hom functor for $D_{\operatorname{qc}}(X)$, and by $\mathbb{R}\operatorname{\mathcal{H}\! \mathit{om}}$ the internal Hom for $D(X)$. 

\begin{lemma}
    \label{lem:gortz_wedhorn_internal_hom}
    Let $f\colon W \to Y$ be a morphism of schemes. Choose $E,G\in D(Y)$. Assume one of the following holds:
    \begin{enumerate}
        \item \label{lem:gortz_wedhorn_internal_hom1} $E$ is perfect
        \item \label{lem:gortz_wedhorn_internal_hom2} $E$ is pseudocoherent and both $G,\mathbf{L}f^\ast G$ are locally bounded below.
    \end{enumerate}
    Then there exists an isomorphism
    \begin{displaymath}
        \mathbf{L}f^\ast \mathbf{R}\operatorname{\mathcal{H}\! \mathit{om}} (E,G) \to  \mathbf{R}\operatorname{\mathcal{H}\! \mathit{om}} (\mathbf{L}f^\ast E , \mathbf{L}f^\ast G).
    \end{displaymath}
\end{lemma}

\begin{proof}
    In both cases, \cite[\href{https://stacks.math.columbia.edu/tag/0A6H}{Tag 0A6H}]{stacks-project} yields isomorphisms 
    \begin{displaymath}
        \mathbf{R}\operatorname{\mathcal{H}\! \mathit{om}}(E,G)
        \cong \mathbb{R}\operatorname{\mathcal{H}\! \mathit{om}}(E,G)
    \end{displaymath}
    and 
    \begin{displaymath}
        \mathbf{R}\operatorname{\mathcal{H}\! \mathit{om}}(\mathbf{L}f^\ast E, \mathbf{L}f^\ast G)
        \cong \mathbb{R}\operatorname{\mathcal{H}\! \mathit{om}}(\mathbf{L}f^\ast E,\mathbf{L}f^\ast G)
    \end{displaymath}
    By \cite[\href{https://stacks.math.columbia.edu/tag/08I3}{Tag 08I3}]{stacks-project}, there exists a natural morphism 
    \begin{displaymath}
        \mathbf{L}f^\ast \mathbb{R}\operatorname{\mathcal{H}\! \mathit{om}} (E,G) \to  \mathbb{R}\operatorname{\mathcal{H}\! \mathit{om}} (\mathbf{L}f^\ast E , \mathbf{L}f^\ast G).
    \end{displaymath}
    The claim follows is we can show this is an isomorphism. 
    Firstly, \eqref{lem:gortz_wedhorn_internal_hom1} is a special case of \cite[Proposition 22.70]{Gortz/Wedhorn:2023}. 
    Secondly, \eqref{lem:gortz_wedhorn_internal_hom2} is argued exactly as in the case of loc.\ cit.\, where `$f$ has finite tor-dimension' can be replaced with our hypothesis that $\mathbf{L}f^\ast G$ is locally bounded below. 
    %%NOTE: Indeed, this is the only part in the proof of loc.\ cit. which requires the assumption of finite tor-dimension.
\end{proof}

Next, we give a characterization of perfectness. To do so, we need notation. Given a scheme $X$ and a point $p\in X$, denote by $\sigma_p\colon \operatorname{Spec}(\mathcal{O}_{X,p})\to X$ and $i_p \colon \operatorname{Spec}(\kappa(p))\to X$ for the natural morphisms. Note that the set theoretic image of $\sigma_p$ is the intersection of all open subschemes $U$ of $X$ which contain $p$. 

\begin{lemma}
    \label{lem:perfect_around_a_point}
    Let $X$ be a scheme, $E\in D_{\operatorname{qc}}^b(X)$ be pseudocoherent, and $p\in X$. Then the following are equivalent:
    \begin{enumerate}
        \item \label{lem:perfect_around_a_point1} There is an open immersion $j\colon U \to X$ such that $p\in U$ and $\mathbf{L}j^\ast E$ is perfect
        \item \label{lem:perfect_around_a_point2} $\mathbf{L}\sigma_p^\ast E$ is perfect
        \item \label{lem:perfect_around_a_point3} $\mathbf{L}i_p^\ast E$ is bounded.
    \end{enumerate}
\end{lemma}

\begin{proof}
    To see that $\eqref{lem:perfect_around_a_point1}\implies \eqref{lem:perfect_around_a_point2}$ and $\eqref{lem:perfect_around_a_point2}\implies \eqref{lem:perfect_around_a_point3}$, use the natural factorizations of $\sigma_p$ and of $i_p$. Specifically,
    \begin{displaymath}
        \operatorname{Spec}(\kappa(p))
        \to\operatorname{Spec}(\mathcal{O}_{X,p})
        \to U\xrightarrow{j} X.
    \end{displaymath}
    Indeed, one uses the fact that derived pullback preserves perfect complexes (see e.g.\ \cite[\href{https://stacks.math.columbia.edu/tag/09UA}{Tag 09UA}]{stacks-project}).

    Lastly, we check $\eqref{lem:perfect_around_a_point3}\implies \eqref{lem:perfect_around_a_point1}$. In this case, we may assume $X$ is affine. However, the result follows from \cite[Theorem 2.3, (iii)$\implies$(i)]{AlonsoTarrio/JeremiasLopez/SanchodeSalas:2023}.
    %%NOTE: There is $E\in D(\mathcal{O}_{X,x})$ such that $\tilde{E}\cong \mathbf{L}\sigma_x^*\mathcal{E}$ \cite[Lemma 22.36]{Gortz/Wedhorn:2023}.
    % Thus, since tildification commutes with derived pullbacks \cite[Lemma 22.41]{Gortz/Wedhorn:2023}, by \cite[Lemma 22.44]{Gortz/Wedhorn:2023} it suffices to show: given a local ring $R$ with residue field $\kappa$, for a pseudo-coherent complex $K\in D(R)$, if $K\otimes^{\mathbf{L}}_R\kappa$ is bounded, then $K$ is perfect.
    % This is done in the proof of \cite[\href{https://stacks.math.columbia.edu/tag/068W}{Tag 068W}]{stacks-project}, second paragraph.
\end{proof}

The following is an extension of \cite[Theorem 2.3]{AlonsoTarrio/JeremiasLopez/SanchodeSalas:2023} to closed points.

\begin{proposition}
    \label{prop:perfection_is_stalk_local_property}
    Let $X$ be a quasi-separated quasi-compact scheme and $E\in D^b_{\operatorname{qc}}(X)$ be pseudocoherent. Then the following conditions are equivalent:
    \begin{enumerate}
        \item \label{prop:perfection_is_stalk_local_property1} $E$ is perfect
        \item \label{prop:perfection_is_stalk_local_property2} $\mathbf{L}i_p^\ast E$ is bounded for every $p \in X$
        \item \label{prop:perfection_is_stalk_local_property2'} $\mathbf{L}i_p^\ast E$ is bounded for every $p \in X$ which is closed
        \item \label{prop:perfection_is_stalk_local_property3} $\mathbf{L}\sigma_p^\ast E$ is perfect for every $p \in X$
        \item \label{prop:perfection_is_stalk_local_property3'} $\mathbf{L}\sigma_p^\ast E$ is perfect for every $p \in X$ which is closed.
    \end{enumerate}
\end{proposition}

\begin{proof}
    It is straightforward to see that $\eqref{prop:perfection_is_stalk_local_property2}\implies \eqref{prop:perfection_is_stalk_local_property2'}$. Also, by \cite[Theorem 2.3]{AlonsoTarrio/JeremiasLopez/SanchodeSalas:2023}, we have $\eqref{prop:perfection_is_stalk_local_property1}\iff \eqref{prop:perfection_is_stalk_local_property2}$. Moreover, from \Cref{lem:perfect_around_a_point}, we have $\eqref{prop:perfection_is_stalk_local_property3} \iff \eqref{prop:perfection_is_stalk_local_property2}$ and $\eqref{prop:perfection_is_stalk_local_property3'} \iff \eqref{prop:perfection_is_stalk_local_property2'}$. Let $p\in X$.
    Since $X$ is quasi-compact, there is a closed point $x\in X$ such that $x\in\overline{\{p\}}$ \cite[Exercise 3.13]{Gortz/Wedhorn:2020}. Hence, $\sigma_p$ factors through the natural morphism $\operatorname{Spec}(\mathcal{O}_{X,p})\to\operatorname{Spec}(\mathcal{O}_{X,x})$ with $\sigma_x$, and so $\eqref{prop:perfection_is_stalk_local_property3'} \iff \eqref{prop:perfection_is_stalk_local_property3}$ by \cite[\href{https://stacks.math.columbia.edu/tag/09UA}{Tag 09UA}]{stacks-project}.
\end{proof}

Let $X$ be a Noetherian scheme. 
Recall the \textbf{support} of $E\in D_{\operatorname{qc}}(X)$ is given by the set of $p\in X$ such that $\mathbf{L}\sigma_p^\ast E\not\cong 0$.
It is denoted by $\operatorname{Supp}(E)$.
Therefore,
\begin{equation*}
    \label{eq:support_of_complex}
    \operatorname{Supp}(E)
    =\bigcup_{i\in \mathbb{Z}}
    \operatorname{Supp}(\mathcal{H}^i(E)).
\end{equation*}
If $E\in D^-_{\operatorname{coh}}(X)$, then \cite[Lemma A.1]{Iyengar/Lipman/Neeman:2015} says $\operatorname{Supp}(E)$ coincides with $p\in X$ such that $\mathbf{L}i^\ast_p E\not\cong 0$.
If furthermore $E\in D_{\operatorname{coh}}^b(Y)$ then $\operatorname{Supp}(E)$ is closed \cite[\href{https://stacks.math.columbia.edu/tag/01BA}{Tag 01BA}]{stacks-project}.

The following generalizes \cite[Lemma 2.3]{Ruiperez/Hernandez/Martin/SanchodeSalas:2007}.

\begin{lemma}
    \label{lm:fiberwise_vanishing_implies_vanishing}
    Let $f:Y\to S$ be a morphism of Noetherian schemes which preserves closed points, and let $E\in D_{\operatorname{coh}}^b(Y)$ be an object.
    Suppose $\mathbf{L}j_s^*E\cong 0$ for every closed point $s\in S$, where $j_s:Y_s\to Y$ is the immersion of the fiber over $s$.
    Then $E\cong 0$.
\end{lemma}

\begin{proof}
    Condition $E=0$ is equivalent to $\operatorname{Supp}(E)$ being empty.
    Since the subspace $\operatorname{Supp}(E)$ is quasi-compact (any subset of a Noetherian topological space is quasi-compact) and also closed, it can be identified with a closed, reduced and quasi-compact subscheme of $Y$.
    Therefore if $\operatorname{Supp}(E)$ were nonempty it would contain a closed point \cite[Exercise 3.13]{Gortz/Wedhorn:2020}.
    It suffices to show $\operatorname{Supp}(E)$ contains no closed points, which we shall do now.

    Write $r_q\colon \operatorname{Spec}(\kappa(q))\to Y$ for the closed immersion associated to a closed point $q\in Y$, then $s:=(f\circ r_q)(q)$ is a closed point in $S$.
    By \cite[\href{https://stacks.math.columbia.edu/tag/01J9}{Tag 01J9}]{stacks-project}, there exist a field $\ell$ and a commutative diagram
    \begin{displaymath}
        % https://q.uiver.app/#q=WzAsNCxbMSwwLCJcXG9wZXJhdG9ybmFtZXtTcGVjfShcXGthcHBhKHMpKSJdLFsxLDEsIlMiXSxbMCwxLCJcXG9wZXJhdG9ybmFtZXtTcGVjfShcXGthcHBhKHEpKSJdLFswLDAsIlxcb3BlcmF0b3JuYW1le1NwZWN9KFxcZWxsKSJdLFswLDEsInRfcyJdLFsyLDEsImZcXGNpcmMgcl9xIiwyXSxbMywyLCJiXzEiLDJdLFszLDAsImJfMiJdXQ==
        \begin{tikzcd}
            {\operatorname{Spec}(\ell)} & {\operatorname{Spec}(\kappa(s))} \\
            {\operatorname{Spec}(\kappa(q))} & S
            \arrow["{b_2}", from=1-1, to=1-2]
            \arrow["{b_1}"', from=1-1, to=2-1]
            \arrow["{t_s}", from=1-2, to=2-2]
            \arrow["{f\circ r_q}"', from=2-1, to=2-2]
        \end{tikzcd}
    \end{displaymath}
    where $t_s$ is the associated closed immersion.
    Thus, there exists a unique morphism $b$ making the diagram commute,
    \begin{displaymath}
        % https://q.uiver.app/#q=WzAsNixbMiwxLCJcXG9wZXJhdG9ybmFtZXtTcGVjfShcXGthcHBhKHMpKSJdLFsyLDIsIlMuIl0sWzAsMSwiXFxvcGVyYXRvcm5hbWV7U3BlY30oXFxrYXBwYShxKSkiXSxbMCwwLCJcXG9wZXJhdG9ybmFtZXtTcGVjfShcXGVsbCkiXSxbMSwyLCJZXzEiXSxbMSwxLCJZXzFcXHRpbWVzX1MgXFxvcGVyYXRvcm5hbWV7U3BlY30oXFxrYXBwYShzKSkiXSxbMCwxLCJ0X3MiXSxbMywyLCJiXzEiLDJdLFszLDAsImJfMiJdLFs0LDEsImZfMSIsMl0sWzIsNCwicl9xIiwyXSxbNSw0LCJ0XzEiLDJdLFs1LDAsImdfMSIsMl0sWzMsNSwiYiIsMl1d
        \begin{tikzcd}
            {\operatorname{Spec}(\ell)} && \\
            {\operatorname{Spec}(\kappa(q))} & {Y_s} & {\operatorname{Spec}(\kappa(s))} \\
            & Y & {S.}
            \arrow["{b_1}"', from=1-1, to=2-1]
            \arrow["b"', from=1-1, to=2-2]
            \arrow["{b_2}", from=1-1, to=2-3]
            \arrow["{r_q}"', from=2-1, to=3-2]
            \arrow["{g_1}"', from=2-2, to=2-3]
            \arrow["{j_s}"', from=2-2, to=3-2]
            \arrow["{t_s}", from=2-3, to=3-3]
            \arrow["{f_1}"', from=3-2, to=3-3]
        \end{tikzcd}
    \end{displaymath}
    By hypothesis, $\mathbf{L}j_s^\ast E\cong 0$, and hence, $\mathbf{L}(j_s \circ b)^\ast E\cong 0$.
    This implies that $\mathbf{L}(r_q \circ b_1)^\ast E\cong 0$.
    Since $b_1$ is faithfully flat, it follows that $\mathbf{L}r_q^\ast E\cong 0$.
    Therefore, $q\not\in \operatorname{Supp}(E)$, as desired.
\end{proof}

%%%%%%%%%%%%%%%%%%%%%%%%%%%%%%%%
\subsection{Integral transforms}
\label{sec:prelim_integral_transform}
%%%%%%%%%%%%%%%%%%%%%%%%%%%%%%%%

Suppose $f_1 \colon Y_1 \to S$ and $f_2 \colon Y_2 \to S$ are morphisms of finite type to a Noetherian scheme. Consider the fibered square
\begin{displaymath}
    \begin{tikzcd}[ampersand replacement=\&]
        {Y_1\times_S Y_2} \& {Y_2} \\
        {Y_1} \& S.
        \arrow["{p_2}", from=1-1, to=1-2]
        \arrow["{p_1}"', from=1-1, to=2-1]
        \arrow["{f_2}", from=1-2, to=2-2]
        \arrow["{f_1}"', from=2-1, to=2-2]
    \end{tikzcd}
\end{displaymath} 
For any $K\in D_{\operatorname{qc}}(Y_1\times_S Y_2)$, its associated \textbf{integral $S$-transform} is the exact functor $\Phi_K\colon D_{\operatorname{qc}}(Y_1)\to D_{\operatorname{qc}}(Y_2)$ given by $E \mapsto \mathbf{R}p_{2,\ast} (\mathbf{L}p_1^\ast E \otimes^\mathbf{L} K)$. 
We say $Y_1$ and $Y_2$ are \textbf{Fourier--Mukai $S$-partners} if there is such a $K$ for which $\Phi_K$ yields an equivalence $D_{\operatorname{qc}}(Y_1) \to D_{\operatorname{qc}}(Y_2)$. When clear from context, the hyphen `$S$-' is dropped. 

Now, suppose each $f_i$ be proper and flat. An object $K\in D^-_{\operatorname{coh}}(Y_1\times_S Y_2)$ is called \textbf{relatively perfect} over $Y_1$ if 
\begin{displaymath}
    \mathbf{R}(p_1)_\ast (K\otimes^{\mathbf{L}} \operatorname{Perf}(Y_1\times_S Y_2))\subseteq \operatorname{Perf}(Y_1).
\end{displaymath} 
%%OLD: This is equivalent to the condition that
%\begin{displaymath}
%    K\otimes^{\mathbf{L}} \mathbf{L}p_1^\ast D^b_{\operatorname{coh}}%(Y_1)\subseteq D^b_{\operatorname{coh}}(Y_1\times_S Y_2).
%\end{displaymath}
%To see, observe this can be argued affine locally on the target, which case the claim follows from \cite[Lemma 3.4]{Dutta/Lank/ManaliRahul:2025}. 
By symmetry, the notion makes sense over $Y_2$. The following two facts are \cite[Proposition's 3.6 \& 3.2]{Dutta/Lank/ManaliRahul:2025}: $K$ being relatively perfect over $Y_1$ is equivalent to $\Phi_K (D^b_{\operatorname{coh}}(Y_1))\subseteq D^b_{\operatorname{coh}}(Y_2)$; $K$ being relatively perfect over $Y_2$ is the same as $\Phi_K (\operatorname{Perf}(Y_1))\subseteq \operatorname{Perf}(Y_2)$.

More generally, the definition of relatively perfect objects make sense for any proper morphism of Noetherian schemes $f\colon Y \to X$. Indeed, an object $K\in D^-_{\operatorname{coh}}(Y)$ is called \textbf{$f$-perfect} if $\mathbf{R}f_\ast (K\otimes^{\mathbf{L}} \operatorname{Perf}(Y))\subseteq \operatorname{Perf}(X)$. So, in our setting, $K\in D^-_{\operatorname{coh}}(Y_1\times_S Y_2)$ is relatively perfect over $Y_1$ precisely when it is $p_1$-perfect. This notion has appeared in the literature but we record a lemma for convenience to avoid any potential confusion.

\begin{lemma}
    \label{lem:relatively_perfect_pseudo_coherent}
    Let $f\colon Y \to X$ be a proper morphism of Noetherian schemes. Then the following are equivalent for any $E\in D^-_{\operatorname{coh}} (Y)$:
    \begin{enumerate}
        \item \label{lem:relatively_perfect_pseudo_coherent1} $E$ is $f$-perfect
        \item \label{lem:relatively_perfect_pseudo_coherent2} $E \otimes^{\mathbf{L}} \mathbf{L}f^\ast D^b_{\operatorname{coh}}(X)\subseteq D^b_{\operatorname{coh}}(Y)$
        \item \label{lem:relatively_perfect_pseudo_coherent3} $E \otimes^{\mathbf{L}} \mathbf{L}f^\ast D^b_{\operatorname{qc}}(X)\subseteq D^b_{\operatorname{qc}}(Y)$.
    \end{enumerate}
\end{lemma}

\begin{proof}
    $\eqref{lem:relatively_perfect_pseudo_coherent1} \iff \eqref{lem:relatively_perfect_pseudo_coherent2}$ can be argued affine locally on the target by using \cite[Lemma 3.4]{Dutta/Lank/ManaliRahul:2025}; $\eqref{lem:relatively_perfect_pseudo_coherent1} \iff \eqref{lem:relatively_perfect_pseudo_coherent3}$ is \cite[Theorem 4.3]{AlonsoTarrio/JeremiasLopez/SanchodeSalas:2023}.
\end{proof}

\begin{remark}
    A useful observation from \Cref{lem:relatively_perfect_pseudo_coherent} is that any object $E\in D^-_{\operatorname{coh}} (Y)$ being $f$-perfect implies it belongs to $D^b_{\operatorname{coh}}(Y)$.
\end{remark}

%%%%%%%%%%%%%%%%%%%%%%%%%%%%%%%%
\subsection{Generation for triangulated categories}
\label{sec:prelim_generation}
%%%%%%%%%%%%%%%%%%%%%%%%%%%%%%%%

Let $\mathcal{T}$ be a triangulated category with shift functor $[1]\colon \mathcal{T} \to \mathcal{T}$. Fix a subcategory $\mathcal{S}\subseteq \mathcal{T}$. Denote by $\operatorname{add}(\mathcal{S})$ the strictly full (that is, closed under isomorphisms) subcategory of $\mathcal{T}$ consisting of all direct summands of objects of the form $\oplus_{n\in \mathbb{Z}} S_n^{\oplus r_n}[n]$ where $S_n \in \mathcal{S}$ for all $n$ and finitely many $r_n$ are nonzero. Inductively, let 
\begin{displaymath}
    \langle \mathcal{S}\rangle_n :=
    \begin{cases}
        \operatorname{add}(\varnothing) & \text{if }n=0 \\
        \operatorname{add}(\mathcal{S} )& \text{if }n=1 \\
        \operatorname{add}(\{ \operatorname{cone}\phi \mid \phi \in \operatorname{Hom}(\langle \mathcal{S}  \rangle_{n-1}, \langle \mathcal{S}  \rangle_1) \}) & \text{if }n>1.
    \end{cases}
\end{displaymath}
Additionally, let $\langle\mathcal{S} \rangle:=\cup_{n\geq 0} \langle\mathcal{S} \rangle_n$. It may be checked that $\langle\mathcal{S} \rangle$ is the smallest \textbf{thick} (that is, triangulated subcategory closed under direct summands) subcategory containing $\mathcal{S}$. In the case $\mathcal{S}$ consists of a single object $G$, we write the above as $\langle G \rangle_n$ and $\langle G \rangle$. An object $G\in \mathcal{T}$ is called a \textbf{classical generator} for $\mathcal{T}$ if $\langle G \rangle = \mathcal{T}$.

Suppose that $\mathcal{T}$ is a compactly generated triangulated category (to clarify, we require that $\mathcal{T}$ admits all small coproducts in the definition of being compactly generated). Denote the collection of compact objects of $\mathcal{T}$ by $\mathcal{T}^c$. We say $G\in \mathcal{T}^c$ is a \textbf{compact generator} for $\mathcal{T}$ if for any $E\in\mathcal{T}$, one has $\operatorname{Hom}(G,E[n]) = 0$ for all $n\in \mathbb{Z}$ if, and only if, $E=0$. A well-known fact is that classical generators for $\mathcal{T}^c$ coincide with compact generators for $\mathcal{T}$ (see e.g.\ \cite[\href{https://stacks.math.columbia.edu/tag/09SR}{Tag 09SR}]{stacks-project}).

The following is an important case.

\begin{example}
    [{\cite[Theorem 3.1.1]{Bondal/VandenBergh:2003}}]
    \label{ex:perfect_compacts_classical_generator}
    Let $X$ be a quasi-compact quasi-separated scheme. The subcategory of compact objects in $D_{\operatorname{qc}}(X)$ coincides with $\operatorname{Perf}(X)$. Moreover, $D_{\operatorname{qc}}(X)$ is compactly generated by a single object from $\operatorname{Perf}(X)$, see \cite[Theorem 3.1.1]{Bondal/VandenBergh:2003} for details.
\end{example}

\begin{terminology}
    Let $\Phi\colon D_{\operatorname{qc}}(Y_1) \to D_{\operatorname{qc}}(Y_2)$ be an exact functor where each $Y_i$ is Noetherian. Suppose $\mathcal{S}\in \{\operatorname{Perf}, D^b_{\operatorname{coh}}, D^-_{\operatorname{coh}} \}$. If $\Phi_K (\mathcal{S}(Y_1)) \subseteq \mathcal{S}(Y_2)$, then we say $\Phi_K$ \textbf{restricts} to an exact functor on $\mathcal{S}$. The same goes for properties such as being fully faithful or an equivalence.
\end{terminology}

Recall that a triangulated subcategory of $\mathcal{T}$ is called \textbf{localizing} if it is closed under small coproducts.
Given a set $\mathcal{S}$ of objects in $\mathcal{T}$, we denote by $\overline{\langle \mathcal{S} \rangle}$ for the smallest localizing subcategory of $\mathcal{T}$ which contains $\mathcal{S}$.

\begin{lemma}
    \label{lem:localizing_iff_cpt_gen}
    Assume $\mathcal{T}$ is compactly generated. 
    Then $G\in \mathcal{T}^c$ compactly generates $\mathcal{T}$ if, and only if, $\overline{\langle G \rangle} = \mathcal{T}$. 
\end{lemma}

\begin{proof}
    This is well-known but we highlight why this is true,
    Indeed, \cite[Remark 1.16, Theorem 1.17, \& Proposition 9.1.19]{Neeman:2001} says the Verdier localization $\pi \colon \mathcal{T}\to \mathcal{T}/\overline{\langle G \rangle}$ admits a right adjoint $Q\colon \mathcal{T}/\overline{\langle G \rangle}\to \mathcal{T}$.
    Then for each $E\in \mathcal{T}$, there exists a distinguished triangle 
    \begin{displaymath}
        S_E \to E \to (Q\circ \pi)(E) \to S_E [1]
    \end{displaymath}
    where $S_E \in \overline{\langle G \rangle}$ and $(Q\circ \pi)(E)\in \mathcal{T}/\overline{\langle G \rangle}^\perp$ (e.g.\ $A\in \mathcal{T}$ such that $\operatorname{Hom}(G,A)\cong 0$ for all $G\in \mathcal{T}/\overline{\langle G \rangle}$).
    Thus, the claim follows.
\end{proof}

\begin{lemma}
    \label{lem:equivalence_or_fully_faithful_via_compacts}
    Let $F\colon \mathcal{T} \leftrightarrows \mathcal{S} \colon G$ be a pair of exact adjoint functors between triangulated categories which are singly compactly generated. 
    Assume $G$ preserves small coproducts. 
    Then $F$ is an equivalence if, and only if, $F$ restricts to an equivalence $\mathcal{T}^c \to \mathcal{S}^c$.
\end{lemma}

\begin{proof}
    Let $G$ be a compact generator for $\mathcal{T}$.
    By \cite[Theorem 5.1]{Neeman:1996}, $F$ preserves compacts.
    Any equivalence $\mathcal{T}\to\mathcal{S}$ induces an equivalence on compact objects $\mathcal{T}^c \to \mathcal{S}^c$. 
    Conversely, assume $F$ restricts to an equivalence $F:\mathcal{T}^c \to \mathcal{S}^c$.
    Therefore $F:\mathcal{T}^c \to \mathcal{S}^c$ has a right adjoint.
    Thus \cite[Lemma 2.6]{Balmer/DellAmbrogio/Sanders:2016} implies $G$ preserves compact objects and the adjunction of $F$ and $G$ on ambient categories restricts to an adjunction on categories of compact objects.
    To show $F$ is an equivalence we show that $F$ and $G$ are both fully faithful. 
    Denote the counit and unit of the adjunction respectively by $\epsilon$ and $\eta$. The strictly full subcategories $\mathcal{T}^\prime\subseteq\mathcal{T}$ and $\mathcal{S}^\prime\subseteq\mathcal{S}$ consisting respectively of objects where $\eta$ and $\epsilon$ are isomorphisms contain $\mathcal{T}^c$ and $\mathcal{S}^c$.
    One may verify that $\mathcal{T}^\prime$ and $\mathcal{S}^\prime$ are closed under shifts, iterated cones, and small coproducts (hence, homotopy colimits). 
    However, \cite[\href{https://stacks.math.columbia.edu/tag/09SN}{Tag 09SN}]{stacks-project} tells us each object in $\mathcal{T}$ and $\mathcal{S}$ can be obtained respectively by $\mathcal{T}^c$ and $\mathcal{S}^c$ using these operations. Hence, $\mathcal{T} \subseteq \mathcal{T}^\prime$ and $\mathcal{S} \subseteq \mathcal{S}^\prime$, which completes the proof.
\end{proof}

%%%%%%%%%%%%%%%%%%%%%%%%%%%%%%%%%%%%%
\section{Preservation of adjointness}
\label{sec:preserve_adjoints}
%%%%%%%%%%%%%%%%%%%%%%%%%%%%%%%%%%%%%

This section studies the behavior of adjoints for integral transforms under a change of base schemes. To avoid repeatedly stating similar constraints, the following is used as a placeholder throughout our work.

\begin{setup}
    \label{setup:fm_cube_pullback}
    Let $t\colon T \to S$ be a morphism of Noetherian schemes. Suppose $f_i\colon Y_i \to S$ are proper flat morphisms with $i\in \{1,2\}$. Consider the commutative diagram
    \begin{equation}
        \label{diag:fm_cube_pullback}
        % https://q.uiver.app/#q=WzAsOCxbNCwxLCJUIl0sWzQsMywiUyJdLFsxLDMsIllfMSJdLFsyLDIsIllfMiJdLFsxLDEsIllfMVxcdGltZXNfUyBUIl0sWzIsMCwiWV8yXFx0aW1lc19TIFQiXSxbMCwwLCJZXzFcXHRpbWVzX1MgWV8yXFx0aW1lc19TIFQiXSxbMCwyLCJZXzFcXHRpbWVzX1MgWV8yIl0sWzAsMSwidCJdLFsyLDEsImZfMSIsMl0sWzMsMSwiZl8yIl0sWzQsMCwiZ18xIl0sWzUsMCwiZ18yIl0sWzQsMiwidF8xIiwwLHsibGFiZWxfcG9zaXRpb24iOjcwfV0sWzUsMywidF8yIiwyLHsibGFiZWxfcG9zaXRpb24iOjgwLCJzdHlsZSI6eyJib2R5Ijp7Im5hbWUiOiJkYXNoZWQifX19XSxbNiw0LCJnXzJeXFxwcmltZSIsMl0sWzYsNSwiZ18xXlxccHJpbWUiLDJdLFs3LDMsImZeXFxwcmltZV8xIiwwLHsibGFiZWxfcG9zaXRpb24iOjMwLCJzdHlsZSI6eyJib2R5Ijp7Im5hbWUiOiJkYXNoZWQifX19XSxbNywyLCJmXlxccHJpbWVfMiIsMl0sWzYsNywidF5cXHByaW1lIiwyXV0=
        \begin{tikzcd}
            {Y_1\times_S Y_2\times_S T} && {Y_2\times_S T} \\
            & {Y_1\times_S T} &&& T \\
            {Y_1\times_S Y_2} && {Y_2} \\
            & {Y_1} &&& S
            \arrow["{g_1^\prime}"', from=1-1, to=1-3]
            \arrow["{g_2^\prime}"', from=1-1, to=2-2]
            \arrow["{t^\prime}"', from=1-1, to=3-1]
            \arrow["{g_2}", from=1-3, to=2-5]
            \arrow["{t_2}"'{pos=0.8}, dashed, from=1-3, to=3-3]
            \arrow["{g_1}", from=2-2, to=2-5]
            \arrow["{t_1}"{pos=0.7}, from=2-2, to=4-2]
            \arrow["t", from=2-5, to=4-5]
            \arrow["{f^\prime_1}"{pos=0.3}, dashed, from=3-1, to=3-3]
            \arrow["{f^\prime_2}"', from=3-1, to=4-2]
            \arrow["{f_2}", from=3-3, to=4-5]
            \arrow["{f_1}"', from=4-2, to=4-5]
        \end{tikzcd}
    \end{equation}
    obtained by base change along $t$.
\end{setup}

\begin{lemma}
    \label{lem:fm_cube_locally_finite_tor_dimension}
    Consider \Cref{setup:fm_cube_pullback}. Then each face of \Cref{diag:fm_cube_pullback} is a tor-independent square (in the sense of \cite[Definition 3.10.2]{Lipman:2009}).
\end{lemma}

\begin{proof}
    As each $f_i$, $f^\prime_i$, $g_i$, and $g^\prime_i$ are flat, the desired claim follows from \cite[Remark 22.94]{Gortz/Wedhorn:2023}.
\end{proof}

\begin{lemma}
    \label{lem:induced_dbcoh_perf_preservation_upon_pullback_with_t_affine}
    Consider \Cref{setup:fm_cube_pullback} with $t$ an affine morphism. If $K\in D^b_{\operatorname{coh}}(Y_1\times_S Y_2)$ is relatively perfect over $Y_1$ (resp.\ over $Y_2$), then $\Phi_{\mathbf{L} (t^\prime)^\ast K}$ restricts to an exact functor on $D^b_{\operatorname{coh}}$ (resp.\ on $\operatorname{Perf}$).
\end{lemma}

\begin{proof}
    We only prove the case of $K$ being relatively perfect over $Y_2$ 
    %%NOTE: That is, it implies $\Phi_{\mathbf{L} (t^\prime) K}$ restricts to an exact functor on $\operatorname{Perf}$ 
    because a similar argument applies for $D^b_{\operatorname{coh}}$. As $\Phi_K (\operatorname{Perf}(Y_1))\subseteq \operatorname{Perf}(Y_2)$, \cite[Proposition 3.2]{Dutta/Lank/ManaliRahul:2025} implies 
    \begin{displaymath}
        \mathbf{R}(f_1^\prime)_\ast (K \otimes^{\mathbf{L} }\operatorname{Perf}(Y_1\times_S Y_2)) \subseteq \operatorname{Perf}(Y_2).
    \end{displaymath}
    By base change, we know that $t_1,t_2,t^\prime$ are affine morphisms. Choose $G\in \operatorname{Perf}(Y_1\times_S Y_2)$ such that $\operatorname{Perf}(Y_1\times_S Y_2) = \langle G \rangle$. Now, $\mathbf{L}(t^\prime)^\ast G$ satisfies 
    \begin{displaymath}
        \operatorname{Perf}(Y_1\times_S Y_2\times_S T) = \langle \mathbf{L} (t^\prime)^\ast G \rangle,
    \end{displaymath}
    see e.g.\ \cite[\href{https://stacks.math.columbia.edu/tag/0BQT}{Tag 0BQT}]{stacks-project}. 
    Also, via base change, we know that $g_1,g_2, f^\prime_2,f^\prime_1,g_2^\prime,g_1^\prime$ are proper flat morphisms.
    Using flat base change (see e.g.\ \cite[Remark 22.94 \& Theorem 22.99]{Gortz/Wedhorn:2023}), we can tie things together and see that 
    \begin{displaymath}
        \begin{aligned}
            \mathbf{R} (g^\prime_1)_\ast & (\mathbf{L}(t^\prime)^\ast K \otimes^\mathbf{L} \operatorname{Perf}(Y_1\times_S Y_2 \times_S T)) 
            \\&\subseteq \mathbf{R} (g^\prime_1)_\ast \langle \mathbf{L}(t^\prime)^\ast (K \otimes^\mathbf{L} G) \rangle
            \\&\subseteq \langle \mathbf{R} (g^\prime_1)_\ast \mathbf{L}(t^\prime)^\ast (K \otimes^\mathbf{L} G) \rangle
            \\&\subseteq \langle \mathbf{L}t_2^\ast \mathbf{R}(f_1^\prime)_\ast (K \otimes^\mathbf{L} G) \rangle
            \\&\subseteq \operatorname{Perf}(Y_2\times_S T).
        \end{aligned}
    \end{displaymath}
    Therefore, by \cite[Proposition 3.2]{Dutta/Lank/ManaliRahul:2025}, $\Phi_{\mathbf{L}(t^\prime)^\ast K}$ restricts to an exact functor on $\operatorname{Perf}$. 
\end{proof}

\begin{lemma}
	\label{lem:fm_cube_flat_base_change_natural_isomorphism}
	Consider \Cref{setup:fm_cube_pullback}.
	Let $K\in D_{\operatorname{qc}}(Y_1\times_S Y_2)$.
	Then on $D_{\operatorname{qc}}$ there is a natural isomorphism
	\begin{displaymath}
        \beta^K \colon \mathbf{L} t_2^\ast \circ \Phi_{K} \to \Phi_{\mathbf{L}(t^\prime)^\ast K} \circ \mathbf{L} t_1^\ast.
    \end{displaymath}
\end{lemma}

\begin{proof}
	This follows from the string of isomorphisms for each $E\in D_{\operatorname{qc}}(Y_1)$:
    \begin{align*}
        \mathbf{L} t_2^\ast \Phi_{K} (E)
        &= \mathbf{L} t_2^\ast \mathbf{R}(f^\prime_1)_\ast (\mathbf{L} (f^\prime_2)^\ast E \otimes^{\mathbf{L}} K)
        \\&\cong \mathbf{R}(g^\prime_1)_\ast \mathbf{L}(t^\prime)^\ast (\mathbf{L} (f^\prime_2)^\ast E \otimes^{\mathbf{L}} K)
        &\text{(flat base change)}
        \\&\cong \mathbf{R}(g^\prime_1)_\ast ( \mathbf{L}(t^\prime)^\ast  \mathbf{L} (f^\prime_2)^\ast E \otimes^{\mathbf{L}}  \mathbf{L}(t^\prime)^\ast K)
        &\text{(monoidality of $\mathbf{L}(-)^\ast$)}
        \\&\cong \mathbf{R}(g^\prime_1)_\ast ( \mathbf{L} (g^\prime_2)^\ast  \mathbf{L} (t_1)^\ast E \otimes^{\mathbf{L}}  \mathbf{L}(t^\prime)^\ast K)
        &\text{(pseudofunctoriality of $\mathbf{L}(-)^\ast$)}
        \\&= \Phi_{\mathbf{L}(t^\prime)^\ast K}( \mathbf{L} t_1^\ast E). && \qedhere
    \end{align*}
\end{proof}

\begin{lemma}
    \label[lemma]{lem:reflecting_bounded_pseudocoherence_perfectness}
    Consider \Cref{setup:fm_cube_pullback}. Let $K\in D^b_{\operatorname{coh}}(Y_1\times_S Y_2)$. 
    \begin{enumerate}[label=(\arabic*), ref=\theremark(\arabic*)]
        \item \label[lemma]{lem:reflecting_bounded_pseudocoherence_perfectness1} If $\mathbf{L}(t^\prime)^\ast K$ is relatively perfect over $Y_1\times_S T$ (resp.\ $Y_2\times_S T$) where $t$ is affine and faithfully flat, then $K$ is relatively perfect over $Y_1$ (resp.\ $Y_2$). 
        \item \label[lemma]{lem:reflecting_bounded_pseudocoherence_perfectness2} If $\mathbf{L}(t^\prime)^\ast K$ is relatively perfect over $Y_1\times_S \operatorname{Spec}(k)$ (resp.\ $Y_2\times_S \operatorname{Spec}(k)$) for all morphisms $\operatorname{Spec}(k)\to S$ from a field, then $K$ is relatively perfect over $Y_1$ (resp.\ $Y_2$). 
        \item \label[lemma]{lem:reflecting_bounded_pseudocoherence_perfectness3} If $\mathbf{L}(t^\prime)^\ast K$ is relatively perfect over $Y_1\times_S \operatorname{Spec}(\kappa(p))$ (resp.\ $Y_2\times_S \operatorname{Spec}(\kappa(p))$) for all closed immersions $\operatorname{Spec}(\kappa(p))\to S$ associated to a closed point $p\in S$, then $K$ is relatively perfect over $Y_1$ (resp.\ $Y_2$). 
    \end{enumerate}
\end{lemma}

\begin{proof}
    We only show the case for $\mathbf{L}(t^\prime)^\ast K$ being relatively perfect over $Y_1\times_S T$. 
    %%NOTE: That is, it implies $K$ is relatively perfect over $Y_1$. 
    To show the other case,
    %%NOTE: That is, of perfect complexes being preserved by integral transforms, 
    one can argue like below and use \Cref{lem:pseudocoherence_affine_faithfully_flat_cover}. 
    %%OLD: Let us make an observation before proving things. Choose $P_i\in\operatorname{Perf}(Y_i)$ such that $\langle P_i \rangle = \operatorname{Perf}(Y_i)$ for each $i$. From \cite[Lemma 3.4.1]{Bondal/VandenBergh:2003}, we know that $\langle \mathbf{L}f_2^\ast P_1 \otimes^{\mathbf{L}} \mathbf{L} f_1^\ast P_2 \rangle = \operatorname{Perf}(Y_1\times_S Y_2)$. Then, by \cite[\href{https://stacks.math.columbia.edu/tag/0BQT}{Tag 0BQT}]{stacks-project}, it follows that 
    % \begin{displaymath}
    %     \langle \mathbf{L}(t^\prime)^\ast (\mathbf{L}f_2^\ast P_1 \otimes^{\mathbf{L}} \mathbf{L} f_1^\ast P_2) \rangle = \operatorname{Perf}(Y_1\times_S Y_2\times_S T).
    % \end{displaymath}
    %However, we have an isomorphism in \Cref{diag:fm_cube_pullback},
    % \begin{displaymath}
    %     \begin{aligned}
    %         \mathbf{L}(t^\prime)^\ast (\mathbf{L}f_2^\ast P_1 \otimes^{\mathbf{L}} \mathbf{L} f_1^\ast P_2) 
    %         &\cong  \mathbf{L}(t^\prime)^\ast \mathbf{L}f_2^\ast P_1 \otimes^{\mathbf{L}}  \mathbf{L}(t^\prime)^\ast \mathbf{L} f_1^\ast P_2
    %         \\&\cong  \mathbf{L}(g_2^\prime)^\ast \mathbf{L}t_1^\ast P_1 \otimes^{\mathbf{L}} \mathbf{L}(g_1^\prime)^\ast \mathbf{L} t_2^\ast P_2.
    %     \end{aligned}
    % \end{displaymath}

    To start, we prove the first claim. Let $E\in D^b_{\operatorname{coh}}(Y_1)$. As $t^\prime$ is faithfully flat, it follows that $\mathbf{L} t_1^\ast E\in D^b_{\operatorname{coh}}(Y_1\times_S T)$. By \Cref{lem:fm_cube_flat_base_change_natural_isomorphism}, we know that $\mathbf{L} t_2^\ast \Phi_K (E)\cong \Phi_{\mathbf{L}(t^\prime)^\ast K} (\mathbf{L} t^\ast_1 E)$. From \Cref{lem:pseudocoherence_affine_faithfully_flat_cover}, the hypothesis implies $\Phi_K (E) \in D_{\operatorname{coh}}^-(Y_2)$. 
    %%NOTE: This is because $\mathbf{L}(t^\prime)^\ast K$ is relatively perfect over $Y_1\times_S T$, so must be pseudocoherent by definition
    So, it suffices to show that $\mathcal{H}^j (\Phi_K (E))=0$ for all but finitely many $j\in \mathbb{Z}$. Our hypothesis ensures that $\Phi_{\mathbf{L}(t^\prime)^\ast K} (\mathbf{L} t_1^\ast E) \in D^b_{\operatorname{coh}}(Y_2\times_S T)$. Then, as $t_2$ is flat, for each $j\in \mathbb{Z}$,
    \begin{displaymath}
        \begin{aligned}
            \operatorname{Supp}(\mathcal{H}^j (\Phi_K (E)) )
            &= t_2 (\operatorname{Supp}(t_2^\ast \mathcal{H}^j (\Phi_K (E)) )) && \textrm{(\cite[\href{https://stacks.math.columbia.edu/tag/056J}{Tag 056J}]{stacks-project})}
            \\&= t_2 (\operatorname{Supp}(\mathbf{L}t_2^\ast \mathcal{H}^j (\Phi_K (E)) )) && \textrm{(Flatness of $t_2$)}
            \\&= t_2 (\operatorname{Supp}(\mathcal{H}^j (\mathbf{L}t_2^\ast \Phi_K (E)) ))  && \textrm{(Flatness of $t_2$)}.
        \end{aligned}
    \end{displaymath}
    Hence, the desired claim regarding bounded cohomology follows. Consequently, we have $\Phi_K (D^b_{\operatorname{coh}}(Y_1)) \subseteq D^b_{\operatorname{coh}}(Y_2)$, which implies $K$ is relatively perfect over $Y_1$. 

    Next, we check the second claim. Choose any $p\in Y_1$. Denote by $t\colon \operatorname{Spec}(\kappa(f_1(p)))\to S$ for the natural morphism. There is a commutative diagram
    %%NOTE: See bottom of Tag 01J5 and use universal property of fiber products.
    \begin{displaymath}
        % https://q.uiver.app/#q=WzAsNSxbMiwxLCJcXG9wZXJhdG9ybmFtZXtTcGVjfShcXGthcHBhKGZfMShwKSkpIl0sWzIsMiwiUyJdLFsxLDIsIllfMSJdLFsxLDEsIllfMVxcdGltZXNfUyBcXG9wZXJhdG9ybmFtZXtTcGVjfShcXGthcHBhKGZfMShwKSkpIl0sWzAsMCwiXFxvcGVyYXRvcm5hbWV7U3BlY30oXFxrYXBwYShwKSkiXSxbMCwxLCJ0Il0sWzIsMSwiZl8xIl0sWzMsMCwiZ18xIl0sWzMsMiwidF8xIiwyXSxbNCwzLCJoIiwxXSxbNCwyLCJoXzEiLDJdLFs0LDAsImhfMiJdXQ==
        \begin{tikzcd}
            {\operatorname{Spec}(\kappa(p))} && \\
            & {Y_1\times_S \operatorname{Spec}(\kappa(f_1(p)))} & {\operatorname{Spec}(\kappa(f_1(p)))} \\
            & {Y_1} & S
            \arrow["h"{description}, from=1-1, to=2-2]
            \arrow["{h_2}", bend right = -12pt, from=1-1, to=2-3]
            \arrow["{h_1}"', bend right = 12pt, from=1-1, to=3-2]
            \arrow["{g_1}", from=2-2, to=2-3]
            \arrow["{t_1}"', from=2-2, to=3-2]
            \arrow["t", from=2-3, to=3-3]
            \arrow["{f_1}", from=3-2, to=3-3]
        \end{tikzcd}
    \end{displaymath}
    where $h_1$ is the natural morphism. Using the hypothesis, it follows that 
    \begin{displaymath}
        \begin{aligned}
            \mathbf{R}(g^\prime_2)_\ast 
            & (\mathbf{L}(t^\prime)^\ast \big( K \otimes^{\mathbf{L}} \operatorname{Perf}(Y_1\times_S Y_2) \big) )
            \\&\subseteq \mathbf{R}(g^\prime_2)_\ast (\mathbf{L}(t^\prime)^\ast K \otimes^{\mathbf{L}} \operatorname{Perf}(Y_1\times_S Y_2 \times_S \operatorname{Spec}(\kappa(f_1(p)))))
            \\&\subseteq \operatorname{Perf}(Y_1\times_S \operatorname{Spec}(\kappa(f_1(p)))).
        \end{aligned}
    \end{displaymath}
    There is a fibered square 
    \begin{displaymath}
        % https://q.uiver.app/#q=WzAsNCxbMSwxLCJZXzEuIl0sWzEsMCwiWV8xXFx0aW1lc19TIFxcb3BlcmF0b3JuYW1le1NwZWN9KFxca2FwcGEoZl8xKHApKSkiXSxbMCwwLCJZXzFcXHRpbWVzX1MgWV8yXFx0aW1lc19TIFxcb3BlcmF0b3JuYW1le1NwZWN9KFxca2FwcGEoZl8xKHApKSkiXSxbMCwxLCJZXzFcXHRpbWVzX1MgWV8yIl0sWzEsMCwidF8xIl0sWzIsMSwiZ18yXlxccHJpbWUiLDJdLFszLDAsImZeXFxwcmltZV8yIiwyXSxbMiwzLCJ0XlxccHJpbWUiLDJdXQ==
        \begin{tikzcd}
            {Y_1\times_S Y_2\times_S \operatorname{Spec}(\kappa(f_1(p)))} & {Y_1\times_S \operatorname{Spec}(\kappa(f_1(p)))} \\
            {Y_1\times_S Y_2} & {Y_1.}
            \arrow["{g_2^\prime}"', from=1-1, to=1-2]
            \arrow["{t^\prime}"', from=1-1, to=2-1]
            \arrow["{t_1}", from=1-2, to=2-2]
            \arrow["{f^\prime_2}"', from=2-1, to=2-2]
        \end{tikzcd}
    \end{displaymath}
    So, from flat base change, we have for all $P\in \operatorname{Perf}(Y_1\times_S Y_2)$,
    \begin{displaymath}
        \begin{aligned}
            \mathbf{L}t_1^\ast \mathbf{R}(f^\prime_2)_\ast ( K \otimes^{\mathbf{L}} P) 
            &\cong \mathbf{R}(g^\prime_2)_\ast \mathbf{L}(t^\prime)^\ast ( K \otimes^{\mathbf{L}} P) 
            \\&\in \operatorname{Perf}(Y_1\times_S \operatorname{Spec}(\kappa(f_1(p)))).
        \end{aligned}
    \end{displaymath}
    This implies that
    \begin{displaymath}
        \mathbf{L}h_1^\ast \mathbf{R}(f^\prime_2)_\ast ( K \otimes^{\mathbf{L}} P) 
        \cong \mathbf{L}h^\ast \mathbf{L}t_1^\ast \mathbf{R}(f^\prime_2)_\ast ( K \otimes^{\mathbf{L}} P) 
        \in D^b_{\operatorname{coh}}(\kappa(p)),
    \end{displaymath}
    because the derived pullback of perfect complexes remain perfect (see e.g.\ \cite[\href{https://stacks.math.columbia.edu/tag/09UA}{Tag 09UA}]{stacks-project}). Using \cite[Theorem 2.3(iv)]{AlonsoTarrio/JeremiasLopez/SanchodeSalas:2023}, we deduce that $\mathbf{R}(f^\prime_2)_\ast ( K \otimes^{\mathbf{L}} P) \in \operatorname{Perf}(Y_1)$; note that $\mathbf{R}(f^\prime_2)_\ast ( K \otimes^{\mathbf{L}} P)$ is bounded because $K \otimes^{\mathbf{L}} P$ is a bounded pseudocoherent complex and $f^\prime_2$ is proper. Indeed, we have shown that $\mathbf{L}b^\ast \mathbf{R}(f^\prime_2)_\ast ( K \otimes^{\mathbf{L}} P)$ is bounded for all $q\in Y_1$ with natural morphism $b\colon \operatorname{Spec}(\kappa(q)) \to Y_1$. Thus, $K$ is relatively perfect over $Y_1$. 

    Lastly, we check the third claim. However, we can argue essentially like above. In particular, the argument above shows that the stalks of $\mathbf{R}(f^\prime_2)_\ast ( K \otimes^{\mathbf{L}} P)$ at all closed points of $Y_1$ are perfect. 
    Indeed, $f_1$ is proper, and hence a closed morphism on the underlying topological spaces.
    Consequently, the desired claim follows from \Cref{prop:perfection_is_stalk_local_property}.
\end{proof}

\begin{theorem}
    \label[theorem]{thm:bounded_pseudocoherence_perfectness_faithfully_flat_affine}
    Consider \Cref{setup:fm_cube_pullback}. Let $K\in D^b_{\operatorname{coh}}(Y_1\times_S Y_2)$. Then the following are equivalent:
    \begin{enumerate}
        \item \label[theorem]{thm:bounded_pseudocoherence_perfectness_faithfully_flat_affine1} $K$ is relatively perfect over $Y_1$ (resp.\ $Y_2$)
        \item \label[theorem]{thm:bounded_pseudocoherence_perfectness_faithfully_flat_affine2} $\mathbf{L}(t^\prime)^\ast K$ is relatively perfect over $Y_1\times_S T$ (resp.\ $Y_2\times_S T$) for every affine morphism $t$ from a Noetherian scheme
        \item \label[theorem]{thm:bounded_pseudocoherence_perfectness_faithfully_flat_affine3} $\mathbf{L}(t^\prime)^\ast K$ is relatively perfect over $Y_1\times_S \operatorname{Spec}(k)$ (resp.\ $Y_2\times_S \operatorname{Spec}(k)$) for all morphisms $t\colon \operatorname{Spec}(k)\to S$ from a field
        \item \label[theorem]{thm:bounded_pseudocoherence_perfectness_faithfully_flat_affine3_closed} $\mathbf{L}(t^\prime)^\ast K$ is relatively perfect over $Y_1\times_S \operatorname{Spec}(\kappa(p))$ (resp.\ $Y_2\times_S \operatorname{Spec}(\kappa(p))$) for all closed immersions $t\colon \operatorname{Spec}(k(p))\to S$ associated to a closed point $p\in S$
        \item \label[theorem]{thm:bounded_pseudocoherence_perfectness_faithfully_flat_affine4} $\mathbf{L}(t^\prime)^\ast K$ is relatively perfect over $Y_1\times_S T$ (resp.\ $Y_2\times_S T$) for some affine surjective morphism $t\colon T\to S$ (note \Cref{setup:fm_cube_pullback} requires $T$ to be Noetherian).
    \end{enumerate}
\end{theorem}

\begin{proof}
    We only prove the case of being relatively perfect over $Y_1$ because the other can be shown analogously. By \Cref{lem:induced_dbcoh_perf_preservation_upon_pullback_with_t_affine}, we see that \eqref{thm:bounded_pseudocoherence_perfectness_faithfully_flat_affine1} implies \eqref{thm:bounded_pseudocoherence_perfectness_faithfully_flat_affine2}. Moreover, we know that \eqref{thm:bounded_pseudocoherence_perfectness_faithfully_flat_affine2} implies \eqref{thm:bounded_pseudocoherence_perfectness_faithfully_flat_affine3} by \cite[\href{https://stacks.math.columbia.edu/tag/01SI}{Tag 01SI}]{stacks-project}. Clearly, \eqref{thm:bounded_pseudocoherence_perfectness_faithfully_flat_affine3} implies \eqref{thm:bounded_pseudocoherence_perfectness_faithfully_flat_affine3_closed}. Furthermore, to see \eqref{thm:bounded_pseudocoherence_perfectness_faithfully_flat_affine3_closed} implies \eqref{thm:bounded_pseudocoherence_perfectness_faithfully_flat_affine4}, use that \Cref{lem:reflecting_bounded_pseudocoherence_perfectness3} tells us $K$ is relatively perfect over $Y_1$, i.e.\ take $t= 1_S$. 
    
    So, we need to check \eqref{thm:bounded_pseudocoherence_perfectness_faithfully_flat_affine4} implies \eqref{thm:bounded_pseudocoherence_perfectness_faithfully_flat_affine1}. Let $t\colon T \to S$ be an affine surjective morphism such that $\mathbf{L}(t^\prime)^\ast K$ is relatively perfect over $Y_1\times_S T$. Choose any morphism $s\colon \operatorname{Spec}(k)\to S$ from a field. 
    Consider the fiber product
    \begin{displaymath}
        % https://q.uiver.app/#q=WzAsNCxbMSwxLCJTLiJdLFswLDEsIlQiXSxbMSwwLCJcXG9wZXJhdG9ybmFtZXtTcGVjfShrKSJdLFswLDAsIlRcXHRpbWVzX1MgXFxvcGVyYXRvcm5hbWV7U3BlY30oaykiXSxbMSwwLCJ0IiwyXSxbMiwwLCJzIl0sWzMsMiwidF5cXHByaW1lIl0sWzMsMSwic15cXHByaW1lIiwyXV0=
        \begin{tikzcd}
            {T\times_S \operatorname{Spec}(k)} & {\operatorname{Spec}(k)} \\
            T & {S.}
            \arrow["{t^\prime}", from=1-1, to=1-2]
            \arrow["{s^\prime}"', from=1-1, to=2-1]
            \arrow["s", from=1-2, to=2-2]
            \arrow["t"', from=2-1, to=2-2]
        \end{tikzcd}
    \end{displaymath}
    By base change, $s^\prime$ is affine and $t^\prime$ is an affine surjection. To avoid complicated diagrams involving base changes, we need some notation. Specifically, let $K_{\#}$ be the derived pullback of $K$ along the natural morphism $Y_1 \times_S Y_2 \times_S \# \to Y_1\times_S Y_2$ for $\#$ any scheme appearing in the diagram above. Now, the hypothesis tells us that $K_T$ is relatively perfect over $Y_1\times_S T$. Hence, \Cref{lem:induced_dbcoh_perf_preservation_upon_pullback_with_t_affine} implies $K_{T\times_S \operatorname{Spec}(k)}$ is relatively perfect over $Y_1\times_S T\times_S \operatorname{Spec}(k)$ because $s^\prime$ is affine. However, $t^\prime$ is affine and faithfully flat, so \Cref{lem:reflecting_bounded_pseudocoherence_perfectness1} ensures that $K_{\operatorname{Spec}(k)}$ is relatively perfect over $Y_1\times_S \operatorname{Spec}(k)$.
    %%NOTE: In some sense, use derived pullbacks along the square above, so composition of one is that of the other.
    However, $s$ was an arbitrary morphism from a field, so \Cref{lem:reflecting_bounded_pseudocoherence_perfectness2} tells us $K$ is relatively perfect over $Y_1$ as desired.
\end{proof}

\begin{lemma}
    [Base change for relative dualizing complex]
    \label{lem:base_change_relative_dualizing_complex}
    Let $S$ be a quasi-compact quasi-separated scheme. Consider a fibered square
    \begin{displaymath}
        % https://q.uiver.app/#q=WzAsNCxbMSwwLCJTXlxccHJpbWUiXSxbMSwxLCJTIl0sWzAsMSwiWCJdLFswLDAsIlheXFxwcmltZSJdLFswLDEsImciXSxbMiwxLCJmIiwyXSxbMywwLCJmXlxccHJpbWUiXSxbMywyLCJnXlxccHJpbWUiLDJdXQ==
        \begin{tikzcd}
            {X^\prime} & {S^\prime} \\
            X & S
            \arrow["{f^\prime}", from=1-1, to=1-2]
            \arrow["{g^\prime}"', from=1-1, to=2-1]
            \arrow["g", from=1-2, to=2-2]
            \arrow["f"', from=2-1, to=2-2]
        \end{tikzcd}
    \end{displaymath}
    where $f$ is proper, flat, and of finite presentation. Then $\mathbf{L}(g^\prime)^\ast f^! \mathcal{O}_S$ is a relative dualizing complex for $f^\prime$. In fact, there is an isomorphism $\mathbf{L}(g^\prime)^\ast f^! \mathcal{O}_S \to (f^\prime)^! \mathcal{O}_{S^\prime}$.
\end{lemma}

\begin{proof}
    This is known but we include it for convenience. By \cite[\href{https://stacks.math.columbia.edu/tag/0B6S}{Tag 0B6S} \& \href{https://stacks.math.columbia.edu/tag/0E2Z}{Tag 0E2Z}]{stacks-project}, we know that $(f^\prime)^! \mathcal{O}_{S^\prime}$ and $f^! \mathcal{O}_S$ are respectively relative dualizing complexes for $f^\prime$ and $f$. Moreover, \cite[\href{https://stacks.math.columbia.edu/tag/0E2Y}{Tag 0E2Y}]{stacks-project} tells us that $\mathbf{L}(g^\prime)^\ast f^! \mathcal{O}_S$ is a relative dualizing complex for $f^\prime$. However, \cite[\href{https://stacks.math.columbia.edu/tag/0E2W}{Tag 0E2W}]{stacks-project} implies relative dualizing complexes are unique up to isomorphism, and so $\mathbf{L}(g^\prime)^\ast f^! \mathcal{O}_S \cong (f^\prime)^! \mathcal{O}_{S^\prime}$.
\end{proof}

\begin{proposition}
    \label{prop:right_adjoint_pullback}
    Consider \Cref{setup:fm_cube_pullback} with $t$ affine. Let $K\in D^b_{\operatorname{coh}}(Y_1\times_S Y_2)$ be relatively perfect over $Y_2$. Then $\Phi_{\mathbf{L}(t^\prime)^\ast K^\prime}$ is right adjoint to $\Phi_{\mathbf{L} (t^\prime)^\ast K}$ on $D_{\operatorname{qc}}$ where $K^\prime:= \mathbf{R}\operatorname{\mathcal{H}\! \mathit{om}} (  K,  (f^\prime_1)^! \mathcal{O}_{Y_2})$. In both cases, $K^\prime\in D^b_{\operatorname{coh}}(Y_1\times_S Y_2)$.
\end{proposition}

\begin{proof}
    That the last claim holds follows from \cite[Definition 2.11 \& Lemma 2.12]{Rizzardo:2017}. So, we check the first claim. For any $E\in D_{\operatorname{qc}}(Y_1\times_S T)$ and $A\in D_{\operatorname{qc}}(Y_2\times_S T)$, there is a string of natural isomorphisms obtained by adjunctions:
    \begin{displaymath}
        \begin{aligned}
            \operatorname{Hom} & ( \Phi_{\mathbf{L} (t^\prime)^\ast K} (E) , A )
            \\&= \operatorname{Hom} ( \mathbf{R} (g^\prime_1)_\ast (\mathbf{L} (g^\prime_2)^\ast E \otimes^{\mathbf{L}} \mathbf{L} (t^\prime)^\ast K), A ) && \textrm{(Definition)}
            \\&\cong \operatorname{Hom} ( \mathbf{L} (g^\prime_2)^\ast E \otimes^{\mathbf{L}} \mathbf{L} (t^\prime)^\ast K, (g^\prime_1)^! A ) && \textrm{(see e.g.\ \cite[\href{https://stacks.math.columbia.edu/tag/0A9E}{Tag 0A9E}]{stacks-project})}
            \\&\cong \operatorname{Hom} ( \mathbf{L} (g^\prime_2)^\ast  E , \mathbf{R}\operatorname{\mathcal{H}\! \mathit{om}} ( \mathbf{L} (t^\prime)^\ast K, (g^\prime_1)^! A)) && \textrm{(see e.g.\ \cite[\href{https://stacks.math.columbia.edu/tag/08DH}{Tag 08DH}]{stacks-project})}
            \\&\cong \operatorname{Hom} ( E , \mathbf{R} (g^\prime_2)_\ast \mathbf{R}\operatorname{\mathcal{H}\! \mathit{om}} ( \mathbf{L} (t^\prime)^\ast K, (g^\prime_1)^! A)).
        \end{aligned}
    \end{displaymath}
    The desired claim follows if we can find an isomorphism
    \begin{displaymath}
        \begin{aligned}
            \mathbf{R} & (g^\prime_2)_\ast \mathbf{R}\operatorname{\mathcal{H}\! \mathit{om}} ( \mathbf{L} (t^\prime)^\ast K, (g^\prime_1)^! A) \\&\cong \mathbf{R} (g^\prime_2)_\ast \bigg( \mathbf{L} (t^\prime)^\ast \big( \mathbf{R}\operatorname{\mathcal{H}\! \mathit{om}} (  K,  (f^\prime_1)^! \mathcal{O}_{Y_2})\big) \otimes^\mathbf{L} \mathbf{L} (g^\prime_1)^\ast A \bigg) =: \Phi_{\mathbf{L}(t^\prime)^\ast K^\prime} (A).
        \end{aligned}
    \end{displaymath}
    Moreover, $(f^\prime_1)^! \mathcal{O}_{Y_2} \in D^+_{\operatorname{coh}}(Y_1\times_S Y_2)$ as upper shriek functors for proper morphisms preserve complexes with bounded below and coherent cohomology (see e.g.\ \cite[\href{https://stacks.math.columbia.edu/tag/0AU1}{Tag 0AU1}]{stacks-project}). So, we have a string of isomorphisms
    \begin{displaymath}
        \begin{aligned}
            \mathbf{L}(t^\prime)^\ast & \mathbf{R}\operatorname{\mathcal{H}\! \mathit{om}}( K, (f^\prime_1)^! \mathcal{O}_{Y_2})
            \\&\cong \mathbf{R}\operatorname{\mathcal{H}\! \mathit{om}}(\mathbf{L}(t^\prime)^\ast K, \mathbf{L}(t^\prime)^\ast (f^\prime_1)^! \mathcal{O}_{Y_2}) && \textrm{(\Cref{lem:gortz_wedhorn_internal_hom})}
            \\&\cong \mathbf{R}\operatorname{\mathcal{H}\! \mathit{om}}(\mathbf{L}(t^\prime)^\ast K, (g^\prime_1)^! \mathcal{O}_{Y_2\times_S T}) && \textrm{(\Cref{lem:base_change_relative_dualizing_complex})}.
        \end{aligned}
    \end{displaymath}
    By \Cref{lem:induced_dbcoh_perf_preservation_upon_pullback_with_t_affine}, we see that $\Phi_{\mathbf{L}(t^\prime)^\ast K}$ restricts to an exact functor $\operatorname{Perf}(Y_1 \times_S T)\to \operatorname{Perf}(Y_2 \times_S T)$. Then \cite[Proposition 3.2]{Dutta/Lank/ManaliRahul:2025} implies 
    \begin{displaymath}
        \mathbf{R}(g_1^\prime)_\ast (\mathbf{L}(t^\prime)^\ast K \otimes^{\mathbf{L}} \operatorname{Perf}(Y_1\times_S Y_2 \times_S T))\subseteq \operatorname{Perf}(Y_2\times_S T).
    \end{displaymath}
    There is another string of isomorphisms for all $A \in D_{\operatorname{qc}}$:
    \begin{displaymath}
        \begin{aligned}
            \mathbf{R} & (g^\prime_2)_\ast \mathbf{R}\operatorname{\mathcal{H}\! \mathit{om}} ( \mathbf{L} (t^\prime)^\ast K, (g^\prime_1)^! A) 
            \\&\cong \mathbf{R} (g^\prime_2)_\ast \mathbf{R}\operatorname{\mathcal{H}\! \mathit{om}} \bigg( \mathbf{L} (t^\prime)^\ast K, \mathbf{L} (g^\prime_1)^\ast A \otimes^\mathbf{L} (g^\prime_1)^! \mathcal{O}_{Y_2 \times_S T} \bigg) && \textrm{(\cite[\href{https://stacks.math.columbia.edu/tag/0B6S}{Tag 0B6S}]{stacks-project})}
            \\&\cong \mathbf{R} (g^\prime_2)_\ast \bigg( \mathbf{R}\operatorname{\mathcal{H}\! \mathit{om}} \big( \mathbf{L} (t^\prime)^\ast K, (g^\prime_1)^! \mathcal{O}_{Y_2 \times_S T} \big) \otimes^\mathbf{L} \mathbf{L} (g^\prime_1)^\ast A \bigg)
            \\&\cong \mathbf{R} (g^\prime_2)_\ast \bigg( \mathbf{L} (t^\prime)^\ast \mathbf{R}\operatorname{\mathcal{H}\! \mathit{om}} (  K,  (f^\prime_1)^! \mathcal{O}_{Y_2}) \otimes^\mathbf{L} \mathbf{L} (g^\prime_1)^\ast A \bigg).
        \end{aligned}
    \end{displaymath}
    The last isomorphism comes from the work above, whereas the second isomorphism is \cite[Lemma 2.13]{Rizzardo:2017} together with the fact that $\mathbf{L} (t^\prime)^\ast K$ is $g_1^\prime$-perfect via \Cref{thm:bounded_pseudocoherence_perfectness_faithfully_flat_affine},
    which completes the proof.
\end{proof}

\begin{corollary}
    \label{cor:relatively_perfect_y2_implies_right_adjoint_restrict_to_dbcoh}
    Consider \Cref{setup:fm_cube_pullback}. Let $K\in D^b_{\operatorname{coh}}(Y_1\times_S Y_2)$ be relatively perfect over each $Y_i$. Denote by $K^\prime$ for the kernel of the integral transform obtained in \Cref{prop:right_adjoint_pullback} which is right adjoint to $\Phi_K$ on $D_{\operatorname{qc}}$. Then $\Phi_{K^\prime}$ restricts to $D^b_{\operatorname{coh}}$.
\end{corollary}

\begin{proof}
    By \Cref{thm:bounded_pseudocoherence_perfectness_faithfully_flat_affine}, it suffices to check that $\Phi_{\mathbf{L}(t^\prime)^\ast K^\prime}$ restricts to $D^b_{\operatorname{coh}}$ for every $t\colon \operatorname{Spec}(\kappa(p))\to S$ for $p\in S$ a closed point. So, we can reduce to the case $S=\operatorname{Spec}(k)$ for $k$ a field. In this setting, the desired claim follows from \cite[Lemma 3.1]{Dutta/Lank/ManaliRahul:2025}. Specifically, $\Phi_{K^\prime}\colon D^b_{\operatorname{coh}} (Y_2)\to D_{\operatorname{qc}}(Y_1)$ agrees with the unique functor $\Phi^\prime\colon D^b_{\operatorname{coh}}(Y_2)\to D^b_{\operatorname{coh}}(Y_1)$ of loc.\ cit.\ (indeed, \cite[Example 0.7]{Neeman:2021b} gives a triangulated equivalence between the category of finite cohomological functors on perfect complexes and the bounded derived category of bounded pseudocoherent complexes).
\end{proof}

\begin{lemma}
    \label{lem:left_adjoint}
    Let $f_i \colon Y_i \to S$ be proper flat morphisms to a Noetherian scheme where $i\in \{ 1,2\}$. Denote by $f^\prime_i \colon Y_1\times_S Y_2 \to Y_i$ the natural morphisms. Suppose $K\in D^b_{\operatorname{coh}}(Y_1\times_S Y_2)$ is relatively perfect over $Y_1$. Then $\Phi_K$ admits a left adjoint on $D_{\operatorname{qc}}$. In particular, the left adjoint is of the form $\Phi_{K^\prime}$ where $K^\prime := \mathbf{R} \operatorname{\mathcal{H}\! \mathit{om}} (K, (f^\prime_2)^! \mathcal{O}_{Y_1})$.
\end{lemma}

\begin{proof}
    This is baked into the proof of \cite[Proposition 4.3]{Dutta/Lank/ManaliRahul:2025} but we spell it out for convenience. In short, the desired claim follows from the string of natural isomorphisms for all $E\in D_{\operatorname{qc}}(Y_2)$ and $G\in D_{\operatorname{qc}}(Y_1)$:
    \begin{displaymath}
        \begin{aligned}
            \operatorname{Hom}& (E,\Phi_K (G))
            \cong \operatorname{Hom}( E ,\mathbf{R}(f_1^\prime)_\ast (K \otimes^{\mathbf{L}} \mathbf{L}(f^\prime_2)^\ast G)) && \textrm{(Definition)}
            \\&\cong \operatorname{Hom}( \mathbf{L}(f_1^\prime)^\ast E , K \otimes^{\mathbf{L}} \mathbf{L}(f_2^\prime)^\ast G) && \textrm{(Adjunction)}
            \\&\cong \operatorname{Hom}( \mathbf{R}(f_2^\prime)_\ast ( \mathbf{R} \operatorname{\mathcal{H}\! \mathit{om}} (K, (f^\prime_2)^! \mathcal{O}_{\mathcal{Y}_1}) \otimes^{\mathbf{L}} \mathbf{L}(f_1^\prime)^\ast E ) , G) && (\textrm{\cite[Lemma 3.10]{Ballard:2009}}).
        \end{aligned}
    \end{displaymath}
\end{proof}

\begin{proposition}
    \label{prop:left_adjoint_pullback}
    Consider \Cref{setup:fm_cube_pullback} with $t$ affine. Let $K\in D^b_{\operatorname{coh}}(Y_1\times_S Y_2)$ be relatively perfect over $Y_1$. Then $\Phi_{\mathbf{L}(t^\prime)^\ast K^\prime}$ is left adjoint to $\Phi_{\mathbf{L} (t^\prime)^\ast K}$ on $D_{\operatorname{qc}}$ where $K^\prime:= \mathbf{R}\operatorname{\mathcal{H}\! \mathit{om}} (  K,  (f^\prime_2)^! \mathcal{O}_{Y_1})$. In such a case, $K^\prime\in D^b_{\operatorname{coh}}(Y_1\times_S Y_2)$.
\end{proposition}

\begin{proof}
    That the last claim holds follows from \cite[Definition 2.11 \& Lemma 2.12]{Rizzardo:2017}. So, we check the first claim. By \Cref{lem:induced_dbcoh_perf_preservation_upon_pullback_with_t_affine}, $\mathbf{L} (t^\prime)^\ast K$ is relatively perfect over $Y_1 \times_S T$. Moreover, from \Cref{lem:left_adjoint}, the kernel of the integral transform which is left adjoint to $\Phi_{\mathbf{L} (t^\prime)^\ast K}$ is given by the object $\mathbf{R} \operatorname{\mathcal{H}\! \mathit{om}} ( \mathbf{L} (t^\prime)^\ast K, (g^\prime_2)^! \mathcal{O}_{Y_1\times_S T})$. So, it suffices to show there is an isomorphism
    \begin{displaymath}
        \begin{aligned}
            \mathbf{R}(g_2^\prime)_\ast & ( \mathbf{R} \operatorname{\mathcal{H}\! \mathit{om}} ( \mathbf{L} (t^\prime)^\ast K, (g^\prime_2)^! \mathcal{O}_{Y_1\times_S T}) \otimes^{\mathbf{L}} \mathbf{L}(g_1^\prime)^\ast E )
            \\&\to  \mathbf{R}(g_2^\prime)_\ast ( \mathbf{L} (t^\prime)^\ast  \mathbf{R} \operatorname{\mathcal{H}\! \mathit{om}} ( K, (f^\prime_2)^! \mathcal{O}_{Y_1}) \otimes^{\mathbf{L}} \mathbf{L}(g_1^\prime)^\ast E ).
        \end{aligned}
    \end{displaymath}
    This follows if we can find an isomorphism
    \begin{displaymath}
        \mathbf{R} \operatorname{\mathcal{H}\! \mathit{om}} ( \mathbf{L} (t^\prime)^\ast K, (g^\prime_2)^! \mathcal{O}_{Y_1\times_S T}) 
            \to  \mathbf{L} (t^\prime)^\ast ( \mathbf{R} \operatorname{\mathcal{H}\! \mathit{om}} ( K, (f^\prime_2)^! \mathcal{O}_{Y_1})).
    \end{displaymath}
    As in \Cref{lem:fm_cube_locally_finite_tor_dimension}, $(g^\prime_2)^! \mathcal{O}_{Y_1\times_S T} \in D^+_{\operatorname{coh}}(Y_1\times_S Y_2\times_S T)$ because upper shriek functors for proper morphisms preserve complexes with bounded below and coherent cohomology (see e.g.\ \cite[\href{https://stacks.math.columbia.edu/tag/0AU1}{Tag 0AU1}]{stacks-project}). Hence, we have a string of isomorphisms
    \begin{displaymath}
        \begin{aligned}
            \mathbf{L}(t^\prime)^\ast & \mathbf{R}\operatorname{\mathcal{H}\! \mathit{om}}( K, (f^\prime_2)^! \mathcal{O}_{Y_1})
            \\&\cong \mathbf{R}\operatorname{\mathcal{H}\! \mathit{om}}(\mathbf{L}(t^\prime)^\ast K, \mathbf{L}(t^\prime)^\ast (f^\prime_2)^! \mathcal{O}_{Y_1}) && \textrm{(\Cref{lem:gortz_wedhorn_internal_hom})}
            \\&\cong \mathbf{R}\operatorname{\mathcal{H}\! \mathit{om}}(\mathbf{L}(t^\prime)^\ast K, (g^\prime_2)^! \mathcal{O}_{Y_1\times_S T}) && \textrm{(\Cref{lem:base_change_relative_dualizing_complex})}.
        \end{aligned}
    \end{displaymath}
    This completes the proof.
\end{proof}

%%%%%%%%%%%%%%%%%%%%%%%%%%%%%%%%%%%%%
\section{Fully faithfulness \& equivalences}
\label{sec:fully_faithful_equivalence}
%%%%%%%%%%%%%%%%%%%%%%%%%%%%%%%%%%%%%

This section is concerned with the behavior of fully faithfulness or an equivalence for an integral transform under a change of base scheme.

\begin{lemma}
    [{{\cite[IV.7, Exercise 4]{MacLane:78}}}]
	\label[lemma]{lem:morph_adj}
	Consider a diagram of functors and adjoint functors
	\begin{equation}
        \label{diag:morph_adj}
        % https://q.uiver.app/#q=WzAsNCxbMCwwLCJcXG1hdGhjYWx7Q30iXSxbMCwxLCJcXG1hdGhjYWx7RH0iXSxbMSwwLCJcXG1hdGhjYWx7Q30nIl0sWzEsMSwiXFxtYXRoY2Fse0R9Jy4iXSxbMCwyLCJIIl0sWzEsMywiSyIsMl0sWzAsMSwiRiIsMix7Im9mZnNldCI6Mn1dLFsxLDAsIkciLDIseyJvZmZzZXQiOjJ9XSxbMiwzLCJGJyIsMix7Im9mZnNldCI6Mn1dLFszLDIsIkcnIiwyLHsib2Zmc2V0IjoyfV0sWzYsNywiIiwyLHsibGV2ZWwiOjEsInN0eWxlIjp7Im5hbWUiOiJhZGp1bmN0aW9uIn19XSxbOCw5LCIiLDIseyJsZXZlbCI6MSwic3R5bGUiOnsibmFtZSI6ImFkanVuY3Rpb24ifX1dXQ==
        \begin{tikzcd}
            {\mathcal{C}} & {\mathcal{C}'} \\
            {\mathcal{D}} & {\mathcal{D}'.}
            \arrow["H", from=1-1, to=1-2]
            \arrow[""{name=0, anchor=center, inner sep=0}, "F"', shift right=2, from=1-1, to=2-1]
            \arrow[""{name=1, anchor=center, inner sep=0}, "{F^\prime}"', shift right=2, from=1-2, to=2-2]
            \arrow[""{name=2, anchor=center, inner sep=0}, "G"', shift right=2, from=2-1, to=1-1]
            \arrow["K"', from=2-1, to=2-2]
            \arrow[""{name=3, anchor=center, inner sep=0}, "{G^\prime}"', shift right=2, from=2-2, to=1-2]
            \arrow["\dashv"{anchor=center}, draw=none, from=0, to=2]
            \arrow["\dashv"{anchor=center}, draw=none, from=1, to=3]
        \end{tikzcd}
	\end{equation}
	Let $\beta\colon F^\prime H \Rightarrow KF$ and $\alpha\colon  HG \Rightarrow G^\prime K$ be a pair of natural transformations.
	Then the following conditions are equivalent:
	\begin{enumerate}[label=(\arabic*), ref=\thelemma(\arabic*)]
		\item \label[lemma]{lem:morph_adj1} The following diagram commutes:
		\begin{displaymath}
            % https://q.uiver.app/#q=WzAsNCxbMCwwLCJGJ0hHIl0sWzAsMSwiRidHJ0siXSxbMSwxLCJLLiJdLFsxLDAsIktGRyJdLFswLDMsIlxcYWxwaGEgRyJdLFswLDEsIkYnXFxiZXRhIiwyXSxbMSwyLCJcXHZhcmVwc2lsb24nSyIsMl0sWzMsMiwiS1xcdmFyZXBzaWxvbiJdXQ==
			\begin{tikzcd}
				{F^\prime HG} & KFG \\
				{F^\prime G^\prime K} & {K.}
				\arrow["{\beta G}", from=1-1, to=1-2]
				\arrow["{F^\prime \alpha}"', from=1-1, to=2-1]
				\arrow["{K\varepsilon}", from=1-2, to=2-2]
				\arrow["{\varepsilon^\prime K}"', from=2-1, to=2-2]
			\end{tikzcd}
		\end{displaymath}
		\item \label[lemma]{lem:morph_adj2} The following diagram commutes:
		\begin{displaymath}
            % https://q.uiver.app/#q=WzAsNCxbMSwxLCJHXlxccHJpbWUgS0YuIl0sWzAsMCwiSCJdLFsxLDAsIkdeXFxwcmltZSBGXlxccHJpbWUgSCJdLFswLDEsIkhHRiJdLFsyLDAsIkdeXFxwcmltZSBcXGFscGhhIl0sWzMsMCwiXFxiZXRhIEYiLDJdLFsxLDMsIkhcXGV0YSIsMl0sWzEsMiwiXFxldGFeXFxwcmltZSBIIl1d
			\begin{tikzcd}
				H & {G^\prime F^\prime H} \\
				HGF & {G^\prime KF.}
				\arrow["{\eta^\prime H}", from=1-1, to=1-2]
				\arrow["{H\eta}"', from=1-1, to=2-1]
				\arrow["{G^\prime \beta}", from=1-2, to=2-2]
				\arrow["{\alpha F}"', from=2-1, to=2-2]
			\end{tikzcd}
		\end{displaymath}
		\item \label[lemma]{lem:morph_adj3} For all objects $c\in\mathcal{C}$, $d\in\mathcal{D}$,
		the following diagram commutes:
        \begin{displaymath}
            % https://q.uiver.app/#q=WzAsNixbMCwwLCJcXG1hdGhjYWx7RH0oRmMsZCkiXSxbMCwxLCJcXG1hdGhjYWx7RH1eXFxwcmltZShLRmMsS2QpIl0sWzAsMiwiXFxtYXRoY2Fse0R9XlxccHJpbWUgKEZeXFxwcmltZSBIYyxLZCkiXSxbMSwwLCJcXG1hdGhjYWx7Q30oYyxHZCkiXSxbMSwxLCJcXG1hdGhjYWx7Q31eXFxwcmltZSAoSGMsSEdkKSJdLFsxLDIsIlxcbWF0aGNhbHtDfV5cXHByaW1lIChIYyxHXlxccHJpbWUgS2QpLiJdLFswLDMsIlxcY29uZyJdLFsyLDUsIlxcY29uZyJdLFswLDEsIksiLDJdLFsxLDIsIlxcYWxwaGFfY14qIiwyXSxbMyw0LCJIIl0sWzQsNSwiXFxiZXRhX3tkLCp9Il1d
            \begin{tikzcd}
                {\mathcal{D}(Fc,d)} & {\mathcal{C}(c,Gd)} \\
                {\mathcal{D}^\prime(KFc,Kd)} & {\mathcal{C}^\prime (Hc,HGd)} \\
                {\mathcal{D}^\prime (F^\prime Hc,Kd)} & {\mathcal{C}^\prime (Hc,G^\prime Kd).}
                \arrow["\cong", from=1-1, to=1-2]
                \arrow["K"', from=1-1, to=2-1]
                \arrow["H", from=1-2, to=2-2]
                \arrow["{\alpha_c^\ast }"', from=2-1, to=3-1]
                \arrow["{\beta_{d,*}}", from=2-2, to=3-2]
                \arrow["\cong", from=3-1, to=3-2]
            \end{tikzcd}
        \end{displaymath}
	\end{enumerate}	
\end{lemma}

\begin{proof}
	First, we show $\eqref{lem:morph_adj2}\implies \eqref{lem:morph_adj3}$.
	Let $f\colon  Fc \to d$ be a morphism. In the diagram of \eqref{lem:morph_adj3}, $f$ is assigned along the top right composite as
	\begin{displaymath}
		f \mapsto Gf \cdot \eta_c \mapsto H(Gf \cdot \eta_c) \mapsto \alpha_d \cdot H(Gf \cdot \eta_c) = \alpha_d \cdot HGf \cdot H\eta_c.
	\end{displaymath}
	Moreover, $f$ is assigned along the left bottom composite as
	\begin{displaymath}
		f \mapsto Kf \mapsto Kf \cdot \beta_c \mapsto G^\prime (Kf \cdot \beta_c) \cdot \eta^\prime _{Hc} = G^\prime Kf \cdot G^\prime \beta_c \cdot \eta^\prime _{Hc}.
	\end{displaymath}
	So, by the naturality of $\alpha$ and \eqref{lem:morph_adj2}, we have
	\begin{displaymath}
        \begin{aligned}
            G^\prime Kf \cdot G^\prime \beta_c \cdot \eta^\prime _{Hc} 
            = G^\prime Kf \cdot \alpha_{Fc} \cdot H\eta_c 
            = \alpha_d \cdot HGf \cdot H\eta_c.
        \end{aligned}
    \end{displaymath}
	
	Next, we check that $\eqref{lem:morph_adj3}\implies \eqref{lem:morph_adj2}$. For $d=Fc$ and $f = 1_{Fc}$, we have equality between the last morphisms in the composites above. That is, it holds that $\alpha_{Fc} \cdot H\eta_c = G^\prime \beta_c \cdot \eta^\prime _{Hc}$, which gives the desired implication.
	
	To check that $\eqref{lem:morph_adj1}\iff \eqref{lem:morph_adj3}$, one may argue like above because it is essentially dual.
\end{proof}

\begin{definition}
	\label{def:morph_adj}
	A \textbf{(lax) morphism of adjunctions} from $F \dashv G$ to $F^\prime  \dashv G^\prime $ is comprised of a pair of functors
    \begin{displaymath}
        % https://q.uiver.app/#q=WzAsNCxbMCwwLCJcXG1hdGhjYWx7Q30iXSxbMCwxLCJcXG1hdGhjYWx7RH0iXSxbMSwwLCJcXG1hdGhjYWx7Q31eXFxwcmltZSJdLFsxLDEsIlxcbWF0aGNhbHtEfSdeXFxwcmltZSJdLFswLDIsIkgiXSxbMSwzLCJLIiwyXSxbMCwxLCJGIiwyLHsib2Zmc2V0IjoyfV0sWzEsMCwiRyIsMix7Im9mZnNldCI6Mn1dLFszLDIsIkdeXFxwcmltZSIsMix7Im9mZnNldCI6Mn1dLFsyLDMsIkZeXFxwcmltZSIsMix7Im9mZnNldCI6Mn1dLFs2LDcsIiIsMix7ImxldmVsIjoxLCJzdHlsZSI6eyJuYW1lIjoiYWRqdW5jdGlvbiJ9fV0sWzksOCwiIiwyLHsibGV2ZWwiOjEsInN0eWxlIjp7Im5hbWUiOiJhZGp1bmN0aW9uIn19XV0=
        \begin{tikzcd}
            {\mathcal{C}} & {\mathcal{C}^\prime} \\
            {\mathcal{D}} & {\mathcal{D}^\prime}
            \arrow["H", from=1-1, to=1-2]
            \arrow[""{name=0, anchor=center, inner sep=0}, "F"', shift right=2, from=1-1, to=2-1]
            \arrow[""{name=1, anchor=center, inner sep=0}, "{F^\prime}"', shift right=2, from=1-2, to=2-2]
            \arrow[""{name=2, anchor=center, inner sep=0}, "G"', shift right=2, from=2-1, to=1-1]
            \arrow["K"', from=2-1, to=2-2]
            \arrow[""{name=3, anchor=center, inner sep=0}, "{G^\prime}"', shift right=2, from=2-2, to=1-2]
            \arrow["\dashv"{anchor=center}, draw=none, from=0, to=2]
            \arrow["\dashv"{anchor=center}, draw=none, from=1, to=3]
        \end{tikzcd}
    \end{displaymath}
    plus a pair of natural transformations $\beta\colon  F^\prime H \Rightarrow KF$, $\alpha\colon  HG \Rightarrow G^\prime K$ such that the equivalent conditions in Lemma \ref{lem:morph_adj} hold.
	We call $\beta$ (resp.\ $\alpha$) the \textbf{left} (resp.\ \textbf{right}) \textbf{comparison transformation}. Furthermore, a \textbf{strict morphism of adjunctions} is just a lax morphism of adjunctions such that $F^\prime H = KF$, $HG = G^\prime K$, and $\beta$, $\alpha$ are the identity transformation.
\end{definition}

\begin{remark}
    In the category theory literature, a `morphism of adjunctions' is usually referred in its strict sense \cite[IV.7]{MacLane:78} and \cite[Exercise 4.2.v]{Riehl:17}, but a `lax' one is instead described by saying that $\beta$ and $\alpha$ are \textit{mates} under \Cref{diag:morph_adj}. See \cite[Sect. 1]{Cheng/Gurski/Riehl:2014} for the precise definition of mates. Indeed, the mate $\overline{\beta}$ of $\beta$ is by definition the natural transformation
    \begin{displaymath}
        \overline{\beta}\colon 
        HG\xrightarrow{\eta^\prime _{HG}}
        G^\prime F^\prime HG\xrightarrow{G^\prime \beta_G}
        G^\prime KFG\xrightarrow{G^\prime K\varepsilon}G^\prime K.
    \end{displaymath}
    Therefore, $\overline{\beta}$ will equal $\alpha$ if, and only if, the corresponding adjunct morphisms $F^\prime H G\to K$ (with respect to $F^\prime \dashv G^\prime $) are equal. This is exactly \Cref{lem:morph_adj1}. Dually, one also sees that \Cref{lem:morph_adj2} amounts to the mate $\overline{\alpha}$ of $\alpha$ being equal to $\beta$.
\end{remark}

\begin{lemma}
	\label{lem:fm_cube_pullback_natural_isomorphism}
	Consider \Cref{setup:fm_cube_pullback}. Let $K\in D_{\operatorname{qc}}(Y_1\times_S Y_2)$. Then on $D_{\operatorname{qc}}$ there is a natural isomorphism
    \begin{displaymath}
        \alpha^K \colon \Phi_K \circ \mathbf{R}(t_1)_\ast 
        \to \mathbf{R}(t_2)_\ast \circ \Phi_{\mathbf{L}(t^\prime)^\ast K}.
    \end{displaymath}
\end{lemma}

\begin{proof}
	This follows from the string of natural isomorphisms for each $E\in D_{\operatorname{qc}}(Y_1\times_S T)$:
    \begin{align*}
        \Phi_K \circ \mathbf{R} (t_1)_\ast (E)
        &= \mathbf{R} (f_1^\prime)_\ast 
        \big( \mathbf{L} (f^\prime_2)^\ast \mathbf{R} (t_1)_\ast E \otimes^{\mathbf{L}} K \big)
        \\
        &\cong 
        \mathbf{R} (f_1^\prime)_\ast 
        \big( 
        \mathbf{R} t^\prime_\ast \mathbf{L}(g^\prime_2)^\ast E 
        \otimes^{\mathbf{L}} K 
        \big)
        &
        \text{(flat base change)}
        \\
        &\cong 
        \mathbf{R} (f_1^\prime)_\ast 
        \mathbf{R} t^\prime_\ast 
        \big( 
        \mathbf{L}(g^\prime_2)^\ast E 
        \otimes^{\mathbf{L}} \mathbf{L}(t^\prime)^\ast K 
        \big) 
        &
        \text{(projection formula)}
        \\
        &\cong 
        \mathbf{R} (t_2)_\ast \mathbf{R} (g^\prime_1)_\ast
        \big( 
        \mathbf{L}(g^\prime_2)^\ast E 
        \otimes^{\mathbf{L}} \mathbf{L}(t^\prime)^\ast K 
        \big)
        &
        \text{(pseudofunctoriality of $\mathbf{R}(-)_\ast$)}
        \\
        &= 
        \mathbf{R} (t_2)_\ast 
        \circ 
        \Phi_{\mathbf{L}(t^\prime)^\ast K} (E). && \qedhere
    \end{align*}
\end{proof}

\begin{lemma}
	\label{lem:integral_transform_is_morphism_of_adjoints}
	Consider \Cref{setup:fm_cube_pullback}.
	Let $K\in D_{\operatorname{qc}}(Y_1\times_SY_2)$.
	Then we have a morphism of adjoint pairs as depicted in the diagram
	\begin{displaymath}
        % https://q.uiver.app/#q=WzAsNCxbMCwwLCJEX1xccWMoWV8xKSJdLFswLDEsIkRfXFxxYyhZXzFcXHRpbWVzX1NUKSJdLFsxLDAsIkRfXFxxYyhZXzIpIl0sWzEsMSwiRF9cXHFjKFlfMlxcdGltZXNfU1QpIl0sWzAsMiwiXFxQaGlfSyJdLFsxLDMsIlxcUGhpX3tcXG1hdGhiZntMfSh0JyleKkt9IiwyXSxbMCwxLCJcXG1hdGhiZntMfXRfMV4qIiwyLHsib2Zmc2V0IjoyfV0sWzEsMCwiXFxtYXRoYmZ7Un0odF8xKV8qIiwyLHsib2Zmc2V0IjoyfV0sWzIsMywiXFxtYXRoYmZ7TH10XzJeKiIsMix7Im9mZnNldCI6Mn1dLFszLDIsIlxcbWF0aGJme1J9KHRfMilfKiIsMix7Im9mZnNldCI6Mn1dLFs2LDcsIiIsMix7ImxldmVsIjoxLCJzdHlsZSI6eyJuYW1lIjoiYWRqdW5jdGlvbiJ9fV0sWzgsOSwiIiwyLHsibGV2ZWwiOjEsInN0eWxlIjp7Im5hbWUiOiJhZGp1bmN0aW9uIn19XV0=
        \begin{tikzcd}
            {D_{\operatorname{qc}}(Y_1)} & {D_{\operatorname{qc}}(Y_2)} \\
            {D_{\operatorname{qc}}(Y_1\times_ST)} & {D_{\operatorname{qc}}(Y_2\times_ST)}
            \arrow["{\Phi_K}", from=1-1, to=1-2]
            \arrow[""{name=0, anchor=center, inner sep=0}, "{\mathbf{L}t_1^\ast }"', shift right=2, from=1-1, to=2-1]
            \arrow[""{name=1, anchor=center, inner sep=0}, "{\mathbf{L}t_2^\ast }"', shift right=2, from=1-2, to=2-2]
            \arrow[""{name=2, anchor=center, inner sep=0}, "{\mathbf{R}(t_1)_\ast }"', shift right=2, from=2-1, to=1-1]
            \arrow["{\Phi_{\mathbf{L}(t^\prime)^\ast K}}"', from=2-1, to=2-2]
            \arrow[""{name=3, anchor=center, inner sep=0}, "{\mathbf{R}(t_2)_\ast }"', shift right=2, from=2-2, to=1-2]
            \arrow["\dashv"{anchor=center}, draw=none, from=0, to=2]
            \arrow["\dashv"{anchor=center}, draw=none, from=1, to=3]
        \end{tikzcd}
    \end{displaymath}
	with left and right comparison transformations given respectively by the isomorphisms $\beta^K$, $\alpha^K$ of \Cref{lem:fm_cube_pullback_natural_isomorphism,lem:fm_cube_flat_base_change_natural_isomorphism}.
\end{lemma}

\begin{proof}
	Denote by $\xi^i\colon \mathbf{L}t_i^\ast \mathbf{R}(t_i)_\ast \to\operatorname{id}$ the counit of the adjunction $\mathbf{L}t_i^\ast \dashv\mathbf{R}(t_i)_\ast $ with $i=1,2$. To prove the desired claim, it suffices to verify that \Cref{lem:morph_adj1} holds; that is, that the following diagram commutes:
    \begin{displaymath}
        % https://q.uiver.app/#q=WzAsNCxbMCwwLCJcXG1hdGhiZntMfXRfMl5cXGFzdFxcUGhpX0tcXG1hdGhiZntSfSh0XzEpX1xcYXN0Il0sWzEsMCwiXFxQaGlfe1xcbWF0aGJme0x9KHReXFxwcmltZSleXFxhc3QgS31cXG1hdGhiZntMfXRfMV5cXGFzdCBcXG1hdGhiZntSfSh0XzEpX1xcYXN0Il0sWzAsMSwiXFxtYXRoYmZ7TH10XzJeXFxhc3RcXG1hdGhiZntSfSh0XzIpX1xcYXN0XFxQaGlfe1xcbWF0aGJme0x9KHReXFxwcmltZSleXFxhc3QgS30iXSxbMSwxLCJcXFBoaV97XFxtYXRoYmZ7TH0odF5cXHByaW1lKV5cXGFzdCBLfS4iXSxbMCwyLCJcXG1hdGhiZntMfXRfMl5cXGFzdChcXGFscGhhXkspIiwyXSxbMCwxLCJcXGJldGFeS197XFxtYXRoYmZ7Un0odF8xKV9cXGFzdH0iXSxbMSwzLCJcXFBoaV97XFxtYXRoYmZ7TH0odF5cXHByaW1lKV5cXGFzdCBLfShcXHhpXjEpIl0sWzIsMywiXFx4aV4yX3tcXFBoaV97XFxtYXRoYmZ7TH0odF5cXHByaW1lKV5cXGFzdCBLfX0iLDJdXQ==
        \begin{tikzcd}
            {\mathbf{L}t_2^\ast\Phi_K\mathbf{R}(t_1)_\ast} & {\Phi_{\mathbf{L}(t^\prime)^\ast K}\mathbf{L}t_1^\ast \mathbf{R}(t_1)_\ast} \\
            {\mathbf{L}t_2^\ast\mathbf{R}(t_2)_\ast\Phi_{\mathbf{L}(t^\prime)^\ast K}} & {\Phi_{\mathbf{L}(t^\prime)^\ast K}.}
            \arrow["{\beta^K_{\mathbf{R}(t_1)_\ast}}", from=1-1, to=1-2]
            \arrow["{\mathbf{L}t_2^\ast(\alpha^K)}"', from=1-1, to=2-1]
            \arrow["{\Phi_{\mathbf{L}(t^\prime)^\ast K}(\xi^1)}", from=1-2, to=2-2]
            \arrow["{\xi^2_{\Phi_{\mathbf{L}(t^\prime)^\ast K}}}"', from=2-1, to=2-2]
        \end{tikzcd}
    \end{displaymath}
	For brevity, in the remainder of the proof, we drop the $\mathbf{L}$'s and $\mathbf{R}$'s in the notation for the derived functors (i.e.\
	all functors now are understood to be derived).	Hence, the diagram above is now
	\begin{equation}
		\label{eq:square_without_L_and_R}
        % https://q.uiver.app/#q=WzAsNCxbMCwxLCJ0XzJeXFxhc3QgXFxQaGlfSyAodF8xKV9cXGFzdCJdLFsxLDIsIlxcUGhpX3sodF5cXHByaW1lKV5cXGFzdCBLfXRfMV5cXGFzdCAodF8xKV5cXGFzdCJdLFsxLDAsInRfMl5cXGFzdCAodF8yKV9cXGFzdCBcXFBoaV97KHReXFxwcmltZSleXFxhc3QgS30iXSxbMiwxLCJcXFBoaV97KHReXFxwcmltZSleXFxhc3QgS30uIl0sWzAsMiwidF8yXlxcYXN0KFxcYWxwaGFeSykiXSxbMiwzLCJcXHhpXjJfe1xcUGhpX3sodF5cXHByaW1lKV5cXGFzdCBLfX0iXSxbMSwzLCJcXFBoaV97KHReXFxwcmltZSleXFxhc3QgS30oXFx4aV4xKSIsMl0sWzAsMSwiXFxiZXRhXktfeyh0XzEpX1xcYXN0fSIsMl1d
        \begin{tikzcd}
            & {t_2^\ast (t_2)_\ast \Phi_{(t^\prime)^\ast K}} & \\
            {t_2^\ast \Phi_K (t_1)_\ast} && {\Phi_{(t^\prime)^\ast K}.} \\
            & {\Phi_{(t^\prime)^\ast K}t_1^\ast (t_1)^\ast}
            \arrow["{\xi^2_{\Phi_{(t^\prime)^\ast K}}}", from=1-2, to=2-3]
            \arrow["{t_2^\ast(\alpha^K)}", from=2-1, to=1-2]
            \arrow["{\beta^K_{(t_1)_\ast}}"', from=2-1, to=3-2]
            \arrow["{\Phi_{(t^\prime)^\ast K}(\xi^1)}"', from=3-2, to=2-3]
        \end{tikzcd}
	\end{equation}
	Choose $E\in D_{\operatorname{qc}}(Y_1\times_ST)$.
	If we expand the definition of $\alpha^K$ and $\beta^K$ in \eqref{eq:square_without_L_and_R}
	(see the proofs of 
	\Cref{lem:fm_cube_pullback_natural_isomorphism,lem:fm_cube_flat_base_change_natural_isomorphism}),
	we obtain a large diagram consisting of various faces. In what follows, we make this a bit more explicit. Particularly, we spell out each face needed to understand \eqref{eq:square_without_L_and_R}. To describe said diagram more explicitly, we use the following abbreviations for the (natural) isomorphisms:
	\begin{itemize}
		\item $\operatorname{BC}_i$, $i\in\{1,2\}$ for the flat base change isomorphisms
		\item $\operatorname{PF}$ for the projection formula
		\item $F_{(-)_\ast }$ for functoriality of $(-)_\ast $
		\item $F_{(-)^\ast }$ for functoriality of $(-)^\ast $
		\item $M_{(-)^\ast }$ for monoidality of $(-)^\ast $
		\item $\xi^\prime $ for the counit of $(t^\prime)^\ast \dashv (t^\prime)_\ast $.
	\end{itemize}
    Now, consider the following diagrams:
    \begin{equation}
        \label{eq:big_boyA}
        % https://q.uiver.app/#q=WzAsNCxbMCwwLCJ0XzJeXFxhc3QgKGZfMSlfXFxhc3ReXFxwcmltZSAoKGZfMilfXFxhc3ReXFxwcmltZSAgKHRfMSlfXFxhc3QgRVxcb3RpbWVzIEspIl0sWzAsMiwiKGdfMSlfXFxhc3ReXFxwcmltZSAodF5cXHByaW1lKV5cXGFzdCAoKGZfMl5cXHByaW1lKV5cXGFzdCAgICh0XzEpX1xcYXN0IEVcXG90aW1lcyBLKSJdLFsyLDAsInRfMl5cXGFzdCAoZl8xKV9cXGFzdF5cXHByaW1lICgodF5cXHByaW1lKV9cXGFzdCAoZ18yXlxccHJpbWUpXlxcYXN0IEVcXG90aW1lcyBLKSJdLFsyLDIsIihnXzEpX1xcYXN0XlxccHJpbWUgKHReXFxwcmltZSleXFxhc3QgKHRfXFxhc3QgXlxccHJpbWUgKGdfMl5cXHByaW1lKV5cXGFzdCBFXFxvdGltZXMgSykiXSxbMCwxLCJcXG9wZXJhdG9ybmFtZXtCQ31fMiIsMV0sWzAsMiwiXFxvcGVyYXRvcm5hbWV7QkN9XzEiLDFdLFsyLDMsIlxcb3BlcmF0b3JuYW1le0JDfV8yIiwxXSxbMSwzLCJcXG9wZXJhdG9ybmFtZXtCQ31fMSIsMV1d
        \begin{tikzcd}
            {t_2^\ast (f_1^\prime)_\ast ((f_2^\prime)_\ast  (t_1)_\ast E\otimes K)} && {t_2^\ast (f_1^\prime)_\ast ((t^\prime)_\ast (g_2^\prime)^\ast E\otimes K)} \\
            \\
            {(g_1^\prime)_\ast (t^\prime)^\ast ((f_2^\prime)^\ast   (t_1)_\ast E\otimes K)} && {(g_1^\prime)_\ast (t^\prime)^\ast (t_\ast ^\prime (g_2^\prime)^\ast E\otimes K)}
            \arrow["{\operatorname{BC}_1}"{description}, from=1-1, to=1-3]
            \arrow["{\operatorname{BC}_2}"{description}, from=1-1, to=3-1]
            \arrow["{\operatorname{BC}_2}"{description}, from=1-3, to=3-3]
            \arrow["{\operatorname{BC}_1}"{description}, from=3-1, to=3-3]
        \end{tikzcd}
    \end{equation}

    \begin{equation}
        \label{eq:big_boyB}
        % https://q.uiver.app/#q=WzAsNCxbMCwwLCJ0XzJeXFxhc3QgKGZfMSlfXFxhc3ReXFxwcmltZSAoKHReXFxwcmltZSlfXFxhc3QgKGdfMl5cXHByaW1lKV5cXGFzdCBFXFxvdGltZXMgSykiXSxbMCwyLCIoZ18xKV9cXGFzdF5cXHByaW1lICh0XlxccHJpbWUpXlxcYXN0ICh0X1xcYXN0IF5cXHByaW1lIChnXzJeXFxwcmltZSleXFxhc3QgRVxcb3RpbWVzIEspIl0sWzIsMCwidF8yXlxcYXN0IChmXzEpX1xcYXN0XlxccHJpbWUgKHReXFxwcmltZSlfXFxhc3QgKChnXzJeXFxwcmltZSleXFxhc3QgRVxcb3RpbWVzICh0XlxccHJpbWUpXlxcYXN0IEspIl0sWzIsMiwiKGdfMSlfXFxhc3ReXFxwcmltZSAodF5cXHByaW1lKV5cXGFzdCB0X1xcYXN0IF5cXHByaW1lICgoZ18yXlxccHJpbWUpXlxcYXN0IEVcXG90aW1lcyAodF5cXHByaW1lKV5cXGFzdCBLKSJdLFswLDEsIlxcb3BlcmF0b3JuYW1le0JDfV8yIiwxXSxbMCwyLCJcXG9wZXJhdG9ybmFtZXtQRn0iLDFdLFsxLDMsIlxcb3BlcmF0b3JuYW1le1BGfSIsMV0sWzIsMywiXFxvcGVyYXRvcm5hbWV7QkN9XzIiLDFdXQ==
        \begin{tikzcd}
            {t_2^\ast (f_1^\prime)_\ast ((t^\prime)_\ast (g_2^\prime)^\ast E\otimes K)} && {t_2^\ast (f_1^\prime)_\ast (t^\prime)_\ast ((g_2^\prime)^\ast E\otimes (t^\prime)^\ast K)} \\
            \\
            {(g_1^\prime)_\ast (t^\prime)^\ast (t_\ast ^\prime (g_2^\prime)^\ast E\otimes K)} && {(g_1^\prime)_\ast (t^\prime)^\ast t_\ast ^\prime ((g_2^\prime)^\ast E\otimes (t^\prime)^\ast K)}
            \arrow["{\operatorname{PF}}"{description}, from=1-1, to=1-3]
            \arrow["{\operatorname{BC}_2}"{description}, from=1-1, to=3-1]
            \arrow["{\operatorname{BC}_2}"{description}, from=1-3, to=3-3]
            \arrow["{\operatorname{PF}}"{description}, from=3-1, to=3-3]
        \end{tikzcd}
    \end{equation}

    \begin{equation}
        \label{eq:big_boyC}
        % https://q.uiver.app/#q=WzAsNCxbMCwwLCIoZ18xKV9cXGFzdF5cXHByaW1lICh0XlxccHJpbWUpXlxcYXN0ICgoZl8yXlxccHJpbWUpXlxcYXN0ICAgKHRfMSlfXFxhc3QgRVxcb3RpbWVzIEspIl0sWzIsMCwiKGdfMSlfXFxhc3ReXFxwcmltZSAodF5cXHByaW1lKV5cXGFzdCAodF9cXGFzdCBeXFxwcmltZSAoZ18yXlxccHJpbWUpXlxcYXN0IEVcXG90aW1lcyBLKSJdLFswLDIsIihnXzEpX1xcYXN0XlxccHJpbWUgKCh0XlxccHJpbWUpXlxcYXN0IChmXzJeXFxwcmltZSleXFxhc3QgICAodF8xKV9cXGFzdCBFXFxvdGltZXMgKHReXFxwcmltZSleXFxhc3QgSykiXSxbMiwyLCIoZ18xKV9cXGFzdF5cXHByaW1lICgodF5cXHByaW1lKV5cXGFzdCB0X1xcYXN0IF5cXHByaW1lIChnXzJeXFxwcmltZSleXFxhc3QgRVxcb3RpbWVzICh0XlxccHJpbWUpXlxcYXN0IEspIl0sWzAsMSwiXFxvcGVyYXRvcm5hbWV7QkN9XzEiLDFdLFswLDIsIk1feygtKV5cXGFzdH0iLDFdLFsyLDMsIlxcb3BlcmF0b3JuYW1le0JDfV8xIiwxXSxbMSwzLCJNX3soLSl9XlxcYXN0IiwxXSxbMSwyLCIiLDEseyJzdHlsZSI6eyJib2R5Ijp7Im5hbWUiOiJub25lIn0sImhlYWQiOnsibmFtZSI6Im5vbmUifX19XV0=
        \begin{tikzcd}
            {(g_1^\prime)_\ast (t^\prime)^\ast ((f_2^\prime)^\ast   (t_1)_\ast E\otimes K)} && {(g_1^\prime)_\ast (t^\prime)^\ast (t_\ast ^\prime (g_2^\prime)^\ast E\otimes K)} \\
            \\
            {(g_1^\prime)_\ast ((t^\prime)^\ast (f_2^\prime)^\ast   (t_1)_\ast E\otimes (t^\prime)^\ast K)} && {(g_1^\prime)_\ast ((t^\prime)^\ast t_\ast ^\prime (g_2^\prime)^\ast E\otimes (t^\prime)^\ast K)}
            \arrow["{\operatorname{BC}_1}"{description}, from=1-1, to=1-3]
            \arrow["{M_{(-)^\ast}}"{description}, from=1-1, to=3-1]
            \arrow[draw=none, from=1-3, to=3-1]
            \arrow["{M_{(-)}^\ast}"{description}, from=1-3, to=3-3]
            \arrow["{\operatorname{BC}_1}"{description}, from=3-1, to=3-3]
        \end{tikzcd}
    \end{equation}

    \begin{equation}
        \label{eq:big_boyD}
        % https://q.uiver.app/#q=WzAsNCxbMiwwLCIoZ18xKV9cXGFzdF5cXHByaW1lICh0XlxccHJpbWUpXlxcYXN0ICh0X1xcYXN0IF5cXHByaW1lIChnXzJeXFxwcmltZSleXFxhc3QgRVxcb3RpbWVzIEspIl0sWzAsMCwiKGdfMSlfXFxhc3ReXFxwcmltZSAoKHReXFxwcmltZSleXFxhc3QgdF9cXGFzdCBeXFxwcmltZSAoZ18yXlxccHJpbWUpXlxcYXN0IEVcXG90aW1lcyAodF5cXHByaW1lKV5cXGFzdCBLKSJdLFsyLDIsIihnXzEpX1xcYXN0XlxccHJpbWUgKHReXFxwcmltZSleXFxhc3QgdF9cXGFzdCBeXFxwcmltZSAoKGdfMl5cXHByaW1lKV5cXGFzdCBFXFxvdGltZXMgKHReXFxwcmltZSleXFxhc3QgSykiXSxbMCwyLCIoZ18xKV9cXGFzdF5cXHByaW1lICgoZ18yXlxccHJpbWUpXlxcYXN0IEVcXG90aW1lcyAodF5cXHByaW1lKV5cXGFzdCBLKSJdLFswLDEsIk1feygtKX1eXFxhc3QiLDFdLFswLDIsIlxcb3BlcmF0b3JuYW1le1BGfSIsMV0sWzIsMywiXFx4aV5cXHByaW1lIiwxXSxbMSwzLCJcXHhpXlxccHJpbWUiLDFdXQ==
        \begin{tikzcd}
            {(g_1^\prime)_\ast ((t^\prime)^\ast t_\ast ^\prime (g_2^\prime)^\ast E\otimes (t^\prime)^\ast K)} && {(g_1^\prime)_\ast (t^\prime)^\ast (t_\ast ^\prime (g_2^\prime)^\ast E\otimes K)} \\
            \\
            {(g_1^\prime)_\ast ((g_2^\prime)^\ast E\otimes (t^\prime)^\ast K)} && {(g_1^\prime)_\ast (t^\prime)^\ast t_\ast ^\prime ((g_2^\prime)^\ast E\otimes (t^\prime)^\ast K)}
            \arrow["{\xi^\prime}"{description}, from=1-1, to=3-1]
            \arrow["{M_{(-)}^\ast}"{description}, from=1-3, to=1-1]
            \arrow["{\operatorname{PF}}"{description}, from=1-3, to=3-3]
            \arrow["{\xi^\prime}"{description}, from=3-3, to=3-1]
        \end{tikzcd}
    \end{equation}

    \begin{equation}
        \label{eq:big_boyE}
        % https://q.uiver.app/#q=WzAsNCxbMCwwLCJ0XzJeXFxhc3QgKGZfMSlfXFxhc3ReXFxwcmltZSAodF5cXHByaW1lKV9cXGFzdCAoKGdfMl5cXHByaW1lKV5cXGFzdCBFXFxvdGltZXMgKHReXFxwcmltZSleXFxhc3QgSykiXSxbMCwyLCIoZ18xKV9cXGFzdF5cXHByaW1lICh0XlxccHJpbWUpXlxcYXN0IHRfXFxhc3QgXlxccHJpbWUgKChnXzJeXFxwcmltZSleXFxhc3QgRVxcb3RpbWVzICh0XlxccHJpbWUpXlxcYXN0IEspIl0sWzIsMiwiKGdfMSlfXFxhc3ReXFxwcmltZSAoKGdfMl5cXHByaW1lKV5cXGFzdCBFXFxvdGltZXMgKHReXFxwcmltZSleXFxhc3QgSykiXSxbMiwwLCJ0XzJeXFxhc3QgKHRfMilfXFxhc3QgKGdfMSlfXFxhc3ReXFxwcmltZSAoKGdfMl5cXHByaW1lKV5cXGFzdCBFXFxvdGltZXMgKHReXFxwcmltZSleXFxhc3QgSykiXSxbMCwxLCJcXG9wZXJhdG9ybmFtZXtCQ31fMiIsMV0sWzEsMiwiXFx4aV5cXHByaW1lIiwxXSxbMCwzLCJGX3soLSleXFxhc3R9IiwxXSxbMywyLCJcXHhpXjIiLDFdXQ==
        \begin{tikzcd}
            {t_2^\ast (f_1^\prime)_\ast (t^\prime)_\ast ((g_2^\prime)^\ast E\otimes (t^\prime)^\ast K)} && {t_2^\ast (t_2)_\ast (g_1^\prime)_\ast ((g_2^\prime)^\ast E\otimes (t^\prime)^\ast K)} \\
            \\
            {(g_1^\prime)_\ast (t^\prime)^\ast t_\ast ^\prime ((g_2^\prime)^\ast E\otimes (t^\prime)^\ast K)} && {(g_1^\prime)_\ast ((g_2^\prime)^\ast E\otimes (t^\prime)^\ast K)}
            \arrow["{F_{(-)^\ast}}"{description}, from=1-1, to=1-3]
            \arrow["{\operatorname{BC}_2}"{description}, from=1-1, to=3-1]
            \arrow["{\xi^2}"{description}, from=1-3, to=3-3]
            \arrow["{\xi^\prime}"{description}, from=3-1, to=3-3]
        \end{tikzcd}
    \end{equation}

    \begin{equation}
        \label{eq:big_boyF}
        % https://q.uiver.app/#q=WzAsNCxbMCwwLCIoZ18xKV9cXGFzdF5cXHByaW1lICgodF5cXHByaW1lKV5cXGFzdCAoZl8yXlxccHJpbWUpXlxcYXN0ICAgKHRfMSlfXFxhc3QgRVxcb3RpbWVzICh0XlxccHJpbWUpXlxcYXN0IEspIl0sWzIsMCwiKGdfMSlfXFxhc3ReXFxwcmltZSAoKHReXFxwcmltZSleXFxhc3QgdF9cXGFzdCBeXFxwcmltZSAoZ18yXlxccHJpbWUpXlxcYXN0IEVcXG90aW1lcyAodF5cXHByaW1lKV5cXGFzdCBLKSJdLFsyLDIsIihnXzEpX1xcYXN0XlxccHJpbWUgKChnXzJeXFxwcmltZSleXFxhc3QgRVxcb3RpbWVzICh0XlxccHJpbWUpXlxcYXN0IEspIl0sWzAsMiwiKGdeXFxwcmltZV8xKV9cXGFzdCAoKGdfMl5cXHByaW1lKV5cXGFzdCB0XzFeXFxhc3QgICh0XzEpX1xcYXN0IEVcXG90aW1lcyAodF5cXHByaW1lKV5cXGFzdCBLKSJdLFswLDEsIlxcb3BlcmF0b3JuYW1le0JDfV8xIiwxXSxbMSwyLCJcXHhpXlxccHJpbWUiLDFdLFswLDMsIkZfeygtKV5cXGFzdH0iLDFdLFszLDIsIlxceGleMSIsMV1d
        \begin{tikzcd}
            {(g_1^\prime)_\ast ((t^\prime)^\ast (f_2^\prime)^\ast   (t_1)_\ast E\otimes (t^\prime)^\ast K)} && {(g_1^\prime)_\ast ((t^\prime)^\ast t_\ast ^\prime (g_2^\prime)^\ast E\otimes (t^\prime)^\ast K)} \\
            \\
            {(g^\prime_1)_\ast ((g_2^\prime)^\ast t_1^\ast  (t_1)_\ast E\otimes (t^\prime)^\ast K)} && {(g_1^\prime)_\ast ((g_2^\prime)^\ast E\otimes (t^\prime)^\ast K)}
            \arrow["{\operatorname{BC}_1}"{description}, from=1-1, to=1-3]
            \arrow["{F_{(-)^\ast}}"{description}, from=1-1, to=3-1]
            \arrow["{\xi^\prime}"{description}, from=1-3, to=3-3]
            \arrow["{\xi^1}"{description}, from=3-1, to=3-3]
        \end{tikzcd}
    \end{equation}
    Assume that we have shown \eqref{eq:big_boyA}, \eqref{eq:big_boyB}, \eqref{eq:big_boyC}, \eqref{eq:big_boyD}, \eqref{eq:big_boyE}, \eqref{eq:big_boyF} are commutative. Then \eqref{eq:square_without_L_and_R} (in the case of $E$) is the pasting of these diagrams, i.e.\ gives the desired morphism  
    \begin{displaymath}
        % https://q.uiver.app/#q=WzAsMixbMCwwLCJ0XzJeXFxhc3QgKGZfMSlfXFxhc3ReXFxwcmltZSAoKGZfMilfXFxhc3ReXFxwcmltZSAgKHRfMSlfXFxhc3QgRVxcb3RpbWVzIEspIl0sWzIsMCwiKGdfMSlfXFxhc3ReXFxwcmltZSAoKGdfMl5cXHByaW1lKV5cXGFzdCBFXFxvdGltZXMgKHReXFxwcmltZSleXFxhc3QgSykuIl0sWzAsMV1d
        \begin{tikzcd}
            {t_2^\ast (f_1^\prime)_\ast ((f_2^\prime)_\ast  (t_1)_\ast E\otimes K)} && {(g_1^\prime)_\ast ((g_2^\prime)^\ast E\otimes (t^\prime)^\ast K).}
            \arrow[from=1-1, to=1-3]
        \end{tikzcd}
    \end{displaymath}
    
    So, let us explain why \eqref{eq:big_boyA}, \eqref{eq:big_boyB}, \eqref{eq:big_boyC}, \eqref{eq:big_boyD}, \eqref{eq:big_boyE}, \eqref{eq:big_boyF} are commutative. As for \eqref{eq:big_boyA}, one may use naturality of $\operatorname{BC}_2$, whereas \eqref{eq:big_boyB} is due to the naturality of $\operatorname{BC}_2$. Also, for \eqref{eq:big_boyC}, one can use the naturality of $M_{(-)^\ast }$.

    Next, we explain \eqref{eq:big_boyD}. Observe that it is the functor $(g_1^\prime)_\ast $ applied to the following diagram evaluated at $((g_2^\prime)^\ast E,K)$:
    \begin{displaymath}
        % https://q.uiver.app/#q=WzAsNCxbMCwwLCIodF5cXHByaW1lKV5cXGFzdCAodF9cXGFzdF5cXHByaW1lICgtKVxcb3RpbWVzIC0pIl0sWzEsMCwiKHReXFxwcmltZSleXFxhc3QgdF9cXGFzdF5cXHByaW1lICgtXFxvdGltZXMgKHReXFxwcmltZSleXFxhc3QgKC0pKSJdLFswLDEsIih0XlxccHJpbWUpXlxcYXN0IHRfXFxhc3ReXFxwcmltZSAoLSlcXG90aW1lcyAodF5cXHByaW1lKV5cXGFzdCgtKSJdLFsxLDEsIigtKVxcb3RpbWVzICh0XlxccHJpbWUpXlxcYXN0KC0pIl0sWzAsMSwiXFx0ZXh0e1BGfSJdLFswLDIsIlxcbWF0aHJte019X3soLSleXFxhc3R9IiwyXSxbMiwzLCJcXHhpXlxccHJpbWUiLDJdLFsxLDMsIlxceGleXFxwcmltZSJdXQ==
        \begin{tikzcd}
            {(t^\prime)^\ast (t_\ast^\prime (-)\otimes -)} & {(t^\prime)^\ast t_\ast^\prime (-\otimes (t^\prime)^\ast (-))} \\
            {(t^\prime)^\ast t_\ast^\prime (-)\otimes (t^\prime)^\ast(-)} & {(-)\otimes (t^\prime)^\ast(-)}
            \arrow["{\text{PF}}", from=1-1, to=1-2]
            \arrow["{M_{(-)^\ast}}"', from=1-1, to=2-1]
            \arrow["{\xi^\prime}", from=1-2, to=2-2]
            \arrow["{\xi^\prime}"', from=2-1, to=2-2]
        \end{tikzcd}
    \end{displaymath}
    which commutes by definition of the projection formula, see e.g.\ proof of \cite[Proposition 22.81]{Gortz/Wedhorn:2023}.

    Now, for \eqref{eq:big_boyE}. Note that it is the following diagram evaluated at $(g_2^\prime)^\ast E\otimes (t^\prime)^\ast K$:
    \begin{displaymath}
        % https://q.uiver.app/#q=WzAsNCxbMCwwLCJ0XzJeXFxhc3QgKGZeXFxwcmltZV8xKV9cXGFzdCB0XlxccHJpbWVfXFxhc3QiXSxbMCwxLCIoZ15cXHByaW1lXzEpX1xcYXN0ICh0XlxccHJpbWUpXlxcYXN0IHReXFxwcmltZV9cXGFzdCJdLFsxLDEsIihnXzFeXFxwcmltZSlfXFxhc3QiXSxbMSwwLCJ0XzJeXFxhc3QgKHRfMilfXFxhc3QgKGdeXFxwcmltZV8xKV9cXGFzdCJdLFswLDMsIlxcbWF0aHJte0Z9X3soLSlfXFxhc3R9Il0sWzAsMSwiXFx0ZXh0e0JDfV8yIiwyXSxbMSwyLCJcXHhpXlxccHJpbWUiLDJdLFszLDIsIlxceGlfMiJdXQ==
        \begin{tikzcd}
            {t_2^\ast (f^\prime_1)_\ast t^\prime_\ast} & {t_2^\ast (t_2)_\ast (g^\prime_1)_\ast} \\
            {(g^\prime_1)_\ast (t^\prime)^\ast t^\prime_\ast} & {(g_1^\prime)_\ast}
            \arrow["{F_{(-)_\ast}}", from=1-1, to=1-2]
            \arrow["{\text{BC}_2}"', from=1-1, to=2-1]
            \arrow["{\xi_2}", from=1-2, to=2-2]
            \arrow["{\xi^\prime}"', from=2-1, to=2-2]
        \end{tikzcd}
    \end{displaymath}
    which commutes by definition of the base change morphism. See \cite[Definition and Remark 21.129]{Gortz/Wedhorn:2023} or \cite[(3.7.2)]{Lipman:2009}.
    
    Lastly, we check \eqref{eq:big_boyF}. However, it is the functor $(g_1^\prime)_\ast (-\otimes (t^\prime)^\ast K)$ applied to the following diagram evaluated at $E$:
    \begin{displaymath}
        % https://q.uiver.app/#q=WzAsNCxbMCwwLCIodF5cXHByaW1lKV5cXGFzdCB0XlxccHJpbWVfXFxhc3QgKGdfMl5cXHByaW1lKV5cXGFzdCJdLFswLDEsIih0XlxccHJpbWUpXlxcYXN0IChmXzJeXFxwcmltZSleXFxhc3QgKHRfMSlfXFxhc3QiXSxbMSwxLCIoZ18yXlxccHJpbWUpXlxcYXN0IHRfMV5cXGFzdCAodF8xKV9cXGFzdCJdLFsxLDAsIihnXzJeXFxwcmltZSleXFxhc3QiXSxbMCwzLCJcXHhpXlxccHJpbWUiXSxbMSwwLCJcXHRleHR7QkN9XzEiXSxbMSwyLCJcXG1hdGhybXtGfV97KC0pXlxcYXN0fSIsMl0sWzIsMywiXFx4aV8xIiwyXV0=
        \begin{tikzcd}
            {(t^\prime)^\ast t^\prime_\ast (g_2^\prime)^\ast} & {(g_2^\prime)^\ast} \\
            {(t^\prime)^\ast (f_2^\prime)^\ast (t_1)_\ast} & {(g_2^\prime)^\ast t_1^\ast (t_1)_\ast}
            \arrow["{\xi^\prime}", from=1-1, to=1-2]
            \arrow["{\text{BC}_1}", from=2-1, to=1-1]
            \arrow["{F_{(-)^\ast}}"', from=2-1, to=2-2]
            \arrow["{\xi_1}"', from=2-2, to=1-2]
        \end{tikzcd}
    \end{displaymath}
    which commutes by definition of the base change morphism.
\end{proof}

\begin{remark}
	Assume the hypothesis of \Cref{lem:integral_transform_is_morphism_of_adjoints}.
	Denote by $\zeta^i$ for the unit of $t_i^\ast \dashv (t_i)_\ast$ if $i\in\{1,2\}$.
	From \Cref{lem:morph_adj,lem:integral_transform_is_morphism_of_adjoints}, we have that the following diagram is commutative:
	\begin{equation}
		\label{eq:integral_transform_is_compatible_with_units_of_upper_and_lower_stars}
        % https://q.uiver.app/#q=WzAsNCxbMCwwLCJcXFBoaV9LIl0sWzAsMSwiXFxQaGlfSyAodF8xKV9cXGFzdCB0XzFeXFxhc3QiXSxbMSwwLCIodF8yKV9cXGFzdCB0XzJeXFxhc3QgXFxQaGlfSyJdLFsxLDEsIih0XzIpX1xcYXN0IFxcUGhpX3sodF5cXHByaW1lKV5cXGFzdCBLfXRfMV5cXGFzdCAuIl0sWzAsMSwiXFxQaGlfSyhcXHpldGFeMSkiLDJdLFswLDIsIlxcemV0YV4yX3tcXFBoaV9LfSJdLFsyLDMsIih0XzIpX1xcYXN0IChcXGJldGFeSykiXSxbMSwzLCJcXGFscGhhXktfe3RfMV5cXGFzdH0iXV0=
        \begin{tikzcd}
            {\Phi_K} & {(t_2)_\ast t_2^\ast \Phi_K} \\
            {\Phi_K (t_1)_\ast t_1^\ast} & {(t_2)_\ast \Phi_{(t^\prime)^\ast K}t_1^\ast .}
            \arrow["{\zeta^2_{\Phi_K}}", from=1-1, to=1-2]
            \arrow["{\Phi_K(\zeta^1)}"', from=1-1, to=2-1]
            \arrow["{(t_2)_\ast (\beta^K)}", from=1-2, to=2-2]
            \arrow["{\alpha^K_{t_1^\ast}}", from=2-1, to=2-2]
        \end{tikzcd}
	\end{equation}
\end{remark}

\begin{proposition}
	\label{prop:pullbacks_is_a_morphism_of_adjoint_integral_transforms}
	Consider \Cref{setup:fm_cube_pullback} where $t$ is affine. Let $K\in D^b_{\operatorname{coh}}(Y_1\times_S Y_2)$ be relatively perfect over $Y_2$. Denote by $K^\prime$ for the kernel of the integral transform obtained in \Cref{prop:right_adjoint_pullback}
	which is right adjoint to $\Phi_K$ on $D_{\operatorname{qc}}$.
	%%NOTE: by same result, $\Phi_{\mathbf{L}(t^\prime)^\ast K^\prime }$ is right adjoint to $\Phi_{\mathbf{L}(t^\prime)^\ast K}$
	Then, after possibly replacing $\Phi_{\mathbf{L}(t^\prime)^\ast K}$ by an isomorphic functor, we have a morphism of adjoint pairs:
    \begin{displaymath}
        % https://q.uiver.app/#q=WzAsNCxbMCwwLCJEX3tcXG9wZXJhdG9ybmFtZXtxY319KFlfMSkiXSxbMCwxLCJEX3tcXG9wZXJhdG9ybmFtZXtxY319KFlfMikiXSxbMSwwLCJEX3tcXG9wZXJhdG9ybmFtZXtxY319KFlfMVxcdGltZXNfU1QpIl0sWzEsMSwiRF97XFxvcGVyYXRvcm5hbWV7cWN9fShZXzJcXHRpbWVzX1MgVCkiXSxbMCwyLCJcXG1hdGhiZntMfXRfMV4qIl0sWzEsMywiXFxtYXRoYmZ7TH10XzJeXFxhc3QiLDJdLFswLDEsIlxcUGhpX0siLDIseyJvZmZzZXQiOjJ9XSxbMSwwLCJcXFBoaV97S15cXHByaW1lfSIsMix7Im9mZnNldCI6Mn1dLFsyLDMsIlxcUGhpX3tcXG1hdGhiZntMfSh0XlxccHJpbWUpXlxcYXN0IEt9IiwyLHsib2Zmc2V0IjoyfV0sWzMsMiwiXFxQaGlfe1xcbWF0aGJme0x9KHReXFxwcmltZSleXFxhc3QgS15cXHByaW1lfSIsMix7Im9mZnNldCI6Mn1dLFs2LDcsIiIsMix7ImxldmVsIjoxLCJzdHlsZSI6eyJuYW1lIjoiYWRqdW5jdGlvbiJ9fV0sWzgsOSwiIiwyLHsibGV2ZWwiOjEsInN0eWxlIjp7Im5hbWUiOiJhZGp1bmN0aW9uIn19XV0=
        \begin{tikzcd}
            {D_{\operatorname{qc}}(Y_1)} & {D_{\operatorname{qc}}(Y_1\times_ST)} \\
            {D_{\operatorname{qc}}(Y_2)} & {D_{\operatorname{qc}}(Y_2\times_S T)}
            \arrow["{\mathbf{L}t_1^\ast }", from=1-1, to=1-2]
            \arrow[""{name=0, anchor=center, inner sep=0}, "{\Phi_K}"', shift right=2, from=1-1, to=2-1]
            \arrow[""{name=1, anchor=center, inner sep=0}, "{\Phi_{\mathbf{L}(t^\prime)^\ast K}}"', shift right=2, from=1-2, to=2-2]
            \arrow[""{name=2, anchor=center, inner sep=0}, "{\Phi_{K^\prime}}"', shift right=2, from=2-1, to=1-1]
            \arrow["{\mathbf{L}t_2^\ast}"', from=2-1, to=2-2]
            \arrow[""{name=3, anchor=center, inner sep=0}, "{\Phi_{\mathbf{L}(t^\prime)^\ast K^\prime}}"', shift right=2, from=2-2, to=1-2]
            \arrow["\dashv"{anchor=center}, draw=none, from=0, to=2]
            \arrow["\dashv"{anchor=center}, draw=none, from=1, to=3]
        \end{tikzcd}
    \end{displaymath}
	with left and right comparison transformations given respectively by the isomorphisms
	\begin{displaymath}
		\begin{aligned}
            (\beta^K)^{-1}\colon \Phi_{\mathbf{L}(t^\prime)^\ast K}\circ\mathbf{L}t_1^\ast 
            \to\mathbf{L}t_2^\ast \circ\Phi_K,\\
            \beta^{K^\prime }\colon \mathbf{L} t_1^\ast \circ \Phi_{K^\prime }
            \to \Phi_{\mathbf{L}(t^\prime)^\ast K^\prime } \circ \mathbf{L} t_2^\ast.
        \end{aligned}
	\end{displaymath}	
\end{proposition}

\begin{proof}
    Write $\eta$ (resp.\ $\eta^\prime$) for the unit of the adjunction $\Phi_{K}\dashv\Phi_{K^\prime }$ (resp.\ of $\Phi_{(t^\prime)^\ast K}\dashv\Phi_{(t^\prime)^\ast K^\prime }$). By \Cref{lem:morph_adj2}, proving the Proposition amounts to showing the following diagram commutes:
	\begin{equation}
        \label{eq:unit_of_integral_transforms_downstairs_and_upstairs_are_compatible}% https://q.uiver.app/#q=WzAsNCxbMCwwLCJcXG1hdGhiZntMfSB0XzFeXFxhc3QiXSxbMCwxLCJcXG1hdGhiZntMfSB0XzFeXFxhc3RcXFBoaV97S15cXHByaW1lfVxcUGhpX0siXSxbMSwxLCJcXFBoaV97XFxtYXRoYmZ7TH0odF5cXHByaW1lKV5cXGFzdCBLXlxccHJpbWV9IFxcbWF0aGJme0x9dF8yXlxcYXN0IFxcUGhpX0suIl0sWzEsMCwiXFxQaGlfe1xcbWF0aGJme0x9KHReXFxwcmltZSleXFxhc3QgS15cXHByaW1lfVxcUGhpX3tcXG1hdGhiZntMfSh0XlxccHJpbWUpXlxcYXN0IEt9IFxcbWF0aGJme0x9dF8xXlxcYXN0Il0sWzAsMSwiXFxtYXRoYmZ7TH0gdF8xXlxcYXN0IChcXGV0YSkiLDJdLFsxLDIsIlxcYmV0YV57S15cXHByaW1lfV97XFxQaGlfS30iLDJdLFsyLDMsIlxcUGhpX3tcXG1hdGhiZntMfSh0XlxccHJpbWUpXlxcYXN0IEteXFxwcmltZX0oXFxiZXRhXkspIiwyXSxbMCwzLCJcXGV0YV5cXHByaW1lX3tcXG1hdGhiZntMfXRfMV5cXGFzdH0iXV0=
        \begin{tikzcd}
            {\mathbf{L} t_1^\ast} & {\Phi_{\mathbf{L}(t^\prime)^\ast K^\prime}\Phi_{\mathbf{L}(t^\prime)^\ast K} \mathbf{L}t_1^\ast} \\
            {\mathbf{L} t_1^\ast\Phi_{K^\prime}\Phi_K} & {\Phi_{\mathbf{L}(t^\prime)^\ast K^\prime} \mathbf{L}t_2^\ast \Phi_K.}
            \arrow["{\eta^\prime_{\mathbf{L}t_1^\ast}}", from=1-1, to=1-2]
            \arrow["{\mathbf{L} t_1^\ast (\eta)}"', from=1-1, to=2-1]
            \arrow["{\beta^{K^\prime}_{\Phi_K}}"', from=2-1, to=2-2]
            \arrow["{\Phi_{\mathbf{L}(t^\prime)^\ast K^\prime}(\beta^K)}"', from=2-2, to=1-2]
        \end{tikzcd}
	\end{equation}
    As in the proof of \Cref{lem:integral_transform_is_morphism_of_adjoints},
	we omit $\mathbf{L}$'s and $\mathbf{R}$'s in the notation for the derived functors to ease notation. Now, \eqref{eq:unit_of_integral_transforms_downstairs_and_upstairs_are_compatible} needs not commute. Although, it does commute if we replace $\Phi_{(t^\prime)^\ast K}$ by some appropriate isomorphic functor.
    Specifically, we claim there is an adjunction $\widetilde{\Phi}_{(t^\prime)^\ast K}\dashv\Phi_{(t^\prime)^\ast K^\prime}$ with unit $\tilde{\eta}^\prime$ making \eqref{eq:unit_of_integral_transforms_downstairs_and_upstairs_are_compatible} commute (i.e.\ inscribing $\widetilde{(-)}$ to $\Phi_{(t^\prime)^\ast K}$ and $\eta^\prime$) and that moreover there is an isomorphism
    $\theta:\Phi_{(t^\prime)^\ast K}\cong \widetilde{\Phi}_{(t^\prime)^\ast K}$ satisfying
    \begin{equation}
    \label{eq:iso_old_and_new_unit}
        \Phi_{(t^\prime)^\ast K^\prime}(\theta)\circ\eta^\prime =\tilde{\eta}^\prime.
    \end{equation}
	Assume that for every $E\in D_{\operatorname{qc}}(Y_1)$ we have shown:
    %% OLD tex:
    %    \begin{enumerate}[leftmargin=*,label={$\mathcal{P}(E)$.},ref=$\mathcal{P}(E)$]
    %        \item \label{ppty:reflection} The object $\Phi_{(t^\prime)^\ast K}t_1^\ast E$ with the morphism 
    % $\Phi_{(t^\prime)^\ast K}(\beta^K_E)\circ
    % \beta^{K^\prime }_{\Phi_KE}\circ
    % t_1^\ast (\eta_E)$ is a reflection of $t_1^\ast E$ along $\Phi_{(t^\prime)^\ast K^\prime }$ (in the sense of \cite[Definition 3.1.1]{Borceaux:1994}).
    %    \end{enumerate}
    \begin{equation}\tag{$\mathcal{P}(E)$}\label{ppty:reflection}
        \begin{minipage}{12cm}
            {The object $\Phi_{(t^\prime)^\ast K}t_1^\ast E$ with the morphism $\Phi_{(t^\prime)^\ast K}(\beta^K_E)\circ \beta^{K^\prime }_{\Phi_KE}\circ t_1^\ast (\eta_E)$ is a reflection of $t_1^\ast E$ along $\Phi_{(t^\prime)^\ast K^\prime }$ (in the sense of \cite[Definition 3.1.1]{Borceaux:1994}).}
        \end{minipage}
    \end{equation}
    Then we can define $\widetilde{\Phi}_{(t^\prime)^\ast K}$ by means of \cite[Proposition 3.1.3]{Borceaux:1994}. On the one hand, we have $\Phi_{(t^\prime)^\ast K}\dashv\Phi_{(t^\prime)^\ast K^\prime}$ with unit $\eta^\prime$; thus
    by \cite[Theorem 3.1.5]{Borceaux:1994} the pair $(\Phi_{(t^\prime)^\ast K}A,\eta^\prime_A)$ is a reflection of $A\in D_{\operatorname{qc}}(Y_1\times_ST)$ along $\Phi_{(t^\prime)^\ast K^\prime }$.
    Now define $\tilde{\eta}^\prime=\{\tilde{\eta}^\prime_A\}_{A\in D_{\operatorname{qc}}(Y_1\times_ST)}$ by setting $\tilde{\eta}^\prime_A:=\eta^\prime_A$ if $A$ is not in the image of $t_1^\ast:D_{\operatorname{qc}}(Y_1)\to D_{\operatorname{qc}}(Y_1\times_ST)$ and setting $\tilde{\eta}^\prime_{t_1^\ast E}$, $E\in D_{\operatorname{qc}}(Y_1)$, to be the unique top morphism in \eqref{eq:unit_of_integral_transforms_downstairs_and_upstairs_are_compatible} making it commute.
    By \cite[Proposition 3.1.3]{Borceaux:1994}, there is a unique functor $\widetilde{\Phi}_{(t^\prime)^\ast K}:D_{\operatorname{qc}}(Y_1\times_ST)\to D_{\operatorname{qc}}(Y_2\times_ST)$ whose action on objects coincides with $\Phi_{(t^\prime)^\ast K}$ and such that $\tilde{\eta}^\prime:\operatorname{id}\to\widetilde{\Phi}_{(t^\prime)^\ast K^\prime}\Phi_{(t^\prime)^\ast K}$ is a natural transformation.
    Moreover, by \cite[Theorem 3.1.5]{Borceaux:1994}, we have $\widetilde{\Phi}_{(t^\prime)^\ast K}\dashv\Phi_{(t^\prime)^\ast K^\prime}$ with unit $\tilde{\eta}^\prime$.
    Finally, by \cite[Proposition 4.4.1]{Riehl:17},
    there is a natural isomorphism $\theta:\Phi_{(t^\prime)^\ast K}\cong \widetilde{\Phi}_{(t^\prime)^\ast K}$ such that \eqref{eq:iso_old_and_new_unit} holds.

	It is left for us to show \ref{ppty:reflection}.
    So, let $A\in D_{\operatorname{qc}}(Y_2\times_T S)$ and
	consider some $\gamma\colon t_1^\ast E\to\Phi_{(t^\prime)^\ast K^\prime }A$. Our goal is to check that there is a unique morphism $\tilde{\gamma}\colon \Phi_{(t^\prime)^\ast K}t_1^\ast E\to A$ such that the diagram
	\begin{equation}
		\label{eq:pullback_unit_is_reflection}
        % https://q.uiver.app/#q=WzAsNSxbMCwwLCJ0XzFeXFxhc3QgRSJdLFswLDEsInRfMV5cXGFzdCBcXFBoaV97S15cXHByaW1lfVxcUGhpX0tFIl0sWzAsMywiXFxQaGlfeyh0XlxccHJpbWUpXlxcYXN0IEteXFxwcmltZSB9dF8yXlxcYXN0IFxcUGhpX0tFIl0sWzAsNCwiXFxQaGlfeyh0XlxccHJpbWUpXlxcYXN0IEteXFxwcmltZX1cXFBoaV97KHReXFxwcmltZSleXFxhc3QgS30gdF8xXlxcYXN0IEUiXSxbMSwyLCJcXFBoaV97KHReXFxwcmltZSleXFxhc3QgS15cXHByaW1lfSBBIl0sWzAsMSwidF8xXlxcYXN0IChcXGV0YV9FKSIsMl0sWzIsMywiXFxQaGlfeyh0XlxccHJpbWUpXlxcYXN0IEteXFxwcmltZX0oXFxiZXRhXktfRSkiLDJdLFsxLDIsIlxcYmV0YV57S15cXHByaW1lfV97XFxQaGlfSyhFKX0iLDJdLFswLDQsIlxcZ2FtbWEiXSxbMyw0LCJcXFBoaV97KHReXFxwcmltZSleXFxhc3QgS15cXHByaW1lfShcXHRpbGRle1xcZ2FtbWF9KSIsMl1d
        \begin{tikzcd}
            {t_1^\ast E} & \\
            {t_1^\ast \Phi_{K^\prime}\Phi_KE} \\[-1.6em]
            & {\Phi_{(t^\prime)^\ast K^\prime} A} \\[-1.6em]
            {\Phi_{(t^\prime)^\ast K^\prime }t_2^\ast \Phi_KE} \\
            {\Phi_{(t^\prime)^\ast K^\prime}\Phi_{(t^\prime)^\ast K} t_1^\ast E}
            \arrow["{t_1^\ast (\eta_E)}"', from=1-1, to=2-1]
            \arrow["\gamma", from=1-1, to=3-2]
            \arrow["{\beta^{K^\prime}_{\Phi_K(E)}}"', from=2-1, to=4-1]
            \arrow["{\Phi_{(t^\prime)^\ast K^\prime}(\beta^K_E)}"', from=4-1, to=5-1]
            \arrow["{\Phi_{(t^\prime)^\ast K^\prime}(\tilde{\gamma})}"', from=5-1, to=3-2]
        \end{tikzcd}
	\end{equation}
	commutes.
	Let $i\in\{1,2\}$.
	Denote $\zeta^i$ to the unit of $t_i^\ast \dashv (t_i)_\ast$.
	Given morphisms
	\begin{gather*}
		\phi\colon t^\ast_i M\to B,\quad
		\psi\colon C\to (t_i)_\ast D,\quad
		\rho\colon L\to\Phi_{K^\prime }P,
	\end{gather*}
	we will use notations
	\begin{gather*}
		\phi^{\flat_i}\colon M\to (t_i)_\ast B,\quad
		\psi^{\sharp_i}\colon t_i^\ast C\to D,\quad
		\rho^\sharp\colon \Phi_K L\to P
	\end{gather*}
	for the adjunct morphisms of $\phi$, $\psi$, $\rho$ with respect to the adjunctions $t_i^\ast \dashv (t_i)_\ast$ and $\Phi_K\dashv\Phi_{K^\prime }$.
	Suppose there is $\tilde{\gamma}$ turning \eqref{eq:pullback_unit_is_reflection} into a commutative diagram.
	Then $\gamma^{\flat_1}\colon E\to  (t_1)_\ast\Phi_{(t^\prime)^\ast K^\prime }A$ equals the outer clockwise composite in the following diagram:
    \begin{displaymath}
        % https://q.uiver.app/#q=WzAsMTAsWzEsMCwiKHRfMSlfXFxhc3QgdF8xXlxcYXN0IEUiXSxbMCwwLCJFIl0sWzEsMSwiKHRfMSlfXFxhc3QgdF8xXlxcYXN0IFxcUGhpX3tLXlxccHJpbWV9XFxQaGlfSyBFIl0sWzEsMiwiKHRfMSlfXFxhc3QgXFxQaGlfeyh0XlxccHJpbWUpXlxcYXN0IEteXFxwcmltZX10XzJeXFxhc3QgXFxQaGlfSyBFIl0sWzEsMywiKHRfMSlfXFxhc3QgXFxQaGlfeyh0XlxccHJpbWUpXlxcYXN0IEteXFxwcmltZX1cXFBoaV97KHReXFxwcmltZSleXFxhc3QgS310XzFeXFxhc3QgRSJdLFsxLDQsIih0XzEpX1xcYXN0IFxcUGhpX3sodF5cXHByaW1lKV5cXGFzdCBLXlxccHJpbWV9QS4iXSxbMCw0LCJcXFBoaV97S15cXHByaW1lfSh0XzIpX1xcYXN0IEEiXSxbMCwxLCJcXFBoaV97S15cXHByaW1lfVxcUGhpX0sgRSJdLFswLDIsIlxcUGhpX3tLXlxccHJpbWV9ICh0XzIpX1xcYXN0IHReXFxhc3RfMiBcXFBoaV9LIEUiXSxbMCwzLCJcXFBoaV97S15cXHByaW1lfSh0XzIpX1xcYXN0IFxcUGhpX3sodF5cXHByaW1lKV5cXGFzdCBLfXRfMV5cXGFzdCBFIl0sWzEsMCwiXFx6ZXRhX0VeMSJdLFswLDIsIih0XzEpX1xcYXN0IHRfMV5cXGFzdCAoXFxldGFfRSkiXSxbMiwzLCIodF8xKV9cXGFzdCAoXFxiZXRhXntLXlxccHJpbWV9X3tcXFBoaV9LKEUpfSkiXSxbMyw0LCIodF8xKV9cXGFzdCBcXFBoaV97KHReXFxwcmltZSleXFxhc3QgS15cXHByaW1lfShcXGJldGFfRV5LKSJdLFs0LDUsIih0XzEpX1xcYXN0IFxcUGhpX3sodF5cXHByaW1lKV5cXGFzdCBLXlxccHJpbWV9KFxcdGlsZGV7XFxnYW1tYX0pIl0sWzYsNSwiXFxhbHBoYV9BXntLXlxccHJpbWV9IiwyXSxbMSw3LCJcXGV0YV9FIl0sWzcsMiwiXFx6ZXRhXjFfe1xcUGhpX3tLXlxccHJpbWV9XFxQaGlfSyBFfSJdLFs3LDgsIlxcUGhpX3tLXlxccHJpbWV9KFxcemV0YV4yX3tcXFBoaV9LIEV9KSIsMV0sWzgsOSwiXFxQaGlfe0teXFxwcmltZX0gKHRfMilfXFxhc3QgKFxcYmV0YV9FXkspIl0sWzksNiwiXFxQaGlfe0teXFxwcmltZX0odF8yKV9cXGFzdChcXHRpbGRle1xcZ2FtbWF9KSJdLFs4LDMsIlxcYWxwaGFee0teXFxwcmltZX1fe3ReXFxhc3RfMlxcUGhpX0tFfSJdLFsxLDYsIlxcZGVsdGEiLDIseyJjdXJ2ZSI6NX1dXQ==
        \begin{tikzcd}
            E & {(t_1)_\ast t_1^\ast E} \\
        	{\Phi_{K^\prime}\Phi_K E} & {(t_1)_\ast t_1^\ast \Phi_{K^\prime}\Phi_K E} \\
        	{\Phi_{K^\prime} (t_2)_\ast t^\ast_2 \Phi_K E} & {(t_1)_\ast \Phi_{(t^\prime)^\ast K^\prime}t_2^\ast \Phi_K E} \\
        	{\Phi_{K^\prime}(t_2)_\ast \Phi_{(t^\prime)^\ast K}t_1^\ast E} & {(t_1)_\ast \Phi_{(t^\prime)^\ast K^\prime}\Phi_{(t^\prime)^\ast K}t_1^\ast E} \\
        	{\Phi_{K^\prime}(t_2)_\ast A} & {(t_1)_\ast \Phi_{(t^\prime)^\ast K^\prime}A.}
        	\arrow["{\zeta_E^1}", from=1-1, to=1-2]
        	\arrow["{\eta_E}"', from=1-1, to=2-1]
        	\arrow["\delta"', bend right = 80pt, from=1-1, to=5-1]
        	\arrow["{(t_1)_\ast t_1^\ast (\eta_E)}", from=1-2, to=2-2]
        	\arrow["{\zeta^1_{\Phi_{K^\prime}\Phi_K E}}", from=2-1, to=2-2]
        	\arrow["{\Phi_{K^\prime}(\zeta^2_{\Phi_K E})}"', from=2-1, to=3-1]
        	\arrow["{(t_1)_\ast (\beta^{K^\prime}_{\Phi_K(E)})}", from=2-2, to=3-2]
        	\arrow["{\alpha^{K^\prime}_{t^\ast_2\Phi_KE}}", from=3-1, to=3-2]
        	\arrow["{\Phi_{K^\prime} (t_2)_\ast (\beta_E^K)}", from=3-1, to=4-1]
        	\arrow["{(t_1)_\ast \Phi_{(t^\prime)^\ast K^\prime}(\beta_E^K)}", from=3-2, to=4-2]
        	\arrow["{\Phi_{K^\prime}(t_2)_\ast(\tilde{\gamma})}", from=4-1, to=5-1]
        	\arrow["{(t_1)_\ast \Phi_{(t^\prime)^\ast K^\prime}(\tilde{\gamma})}", from=4-2, to=5-2]
        	\arrow["{\alpha_A^{K^\prime}}"', from=5-1, to=5-2]
        \end{tikzcd}
    \end{displaymath}
	In this diagram, we note the upper square commutes by naturality of $\zeta^1$, the middle square commutes by \Cref{lem:integral_transform_is_morphism_of_adjoints}
	(it is \eqref{eq:integral_transform_is_compatible_with_units_of_upper_and_lower_stars} with $K$ replaced by $K^\prime $),
	and the bottom rectangle commutes by naturality of $\alpha^{K^\prime }$.
	Set $\delta$ for the left vertical composite in the last diagram, i.e.\
	$\delta=(\alpha^{K^\prime }_A)^{-1}\circ\gamma^{\flat_1}\colon E\to \Phi_{K^\prime }(t_2)_\ast A$. Then it follows that
	\begin{displaymath}
        \delta^\sharp\colon 
        \Phi_KE \xrightarrow{\zeta^2_{\Phi_KE}}
        (t_2)_\ast t_2^\ast \Phi_KE \xrightarrow{(t_2)_\ast(\beta^K_E)}
        (t_2)_\ast\Phi_{(t^\prime)^\ast K}t_1^\ast E \xrightarrow{(t_2)_\ast(\tilde{\gamma})}
        (t_2)_\ast A,
    \end{displaymath}
	which tells us 
    \begin{displaymath}
        (\delta^\sharp)^{\sharp_2}\colon 
        t_2^\ast \Phi_KE \xrightarrow{\beta_E^K}
        \Phi_{(t^\prime)^\ast K}t_1^\ast E \xrightarrow{\tilde{\gamma}}
        A.
    \end{displaymath}
    Consequently, we have
	\begin{displaymath}
        \begin{aligned}
		\tilde{\gamma}
		&=(\delta^\sharp)^{\sharp_2}\circ(\beta_E^K)^{-1}\\
		&=([(\alpha^{K^\prime }_A)^{-1}\circ\gamma^{\flat_1}]^\sharp)^{\sharp_2}
		\circ(\beta_E^K)^{-1},
	\end{aligned}
    \end{displaymath}
	which shows uniqueness of $\tilde{\gamma}$.
	Conversely, defining $\tilde{\gamma}$ by the last formula and reading the previous derivations in reverse order,
	it follows that \eqref{eq:pullback_unit_is_reflection} commutes.
\end{proof}

\begin{remark}
\label{rem:ff_descent}
    Consider the situation as in Proposition \ref{prop:pullbacks_is_a_morphism_of_adjoint_integral_transforms}.
	Denote $\varepsilon$ and $\varepsilon^\prime $ respectively the counits of the adjunctions $\Phi_K\dashv\Phi_{K^\prime }$ and $\Phi_{\mathbf{L}(t^\prime)^\ast K}\dashv\Phi_{\mathbf{L}(t^\prime)^\ast K^\prime }$.
	From \Cref{lem:morph_adj}, it follows that (after possibly replacing $\Phi_{\mathbf{L}(t^\prime)^\ast K}$ by an isomorphic functor) the following diagram commutes:
	\begin{equation}
		\label{eq:counit_of_integral_transforms_downstairs_and_upstairs_are_compatible}
        % https://q.uiver.app/#q=WzAsNCxbMCwwLCJcXFBoaV97XFxtYXRoYmZ7TH0odF5cXHByaW1lKV5cXGFzdCBLfVxcbWF0aGJme0x9dF8xXlxcYXN0IFxcUGhpX3tLXlxccHJpbWV9Il0sWzEsMCwiXFxtYXRoYmZ7TH10XzJeXFxhc3QgXFxQaGlfS1xcUGhpX3tLXlxccHJpbWV9Il0sWzEsMSwiXFxtYXRoYmZ7TH10XzJeXFxhc3QuIl0sWzAsMSwiXFxQaGlfe1xcbWF0aGJme0x9KHReXFxwcmltZSleXFxhc3QgS31cXFBoaV97XFxtYXRoYmZ7TH0odF5cXHByaW1lKV5cXGFzdCBLXlxccHJpbWV9XFxtYXRoYmZ7TH10XzJeXFxhc3QiXSxbMSwwLCJcXGJldGFeS197XFxQaGlfe0teXFxwcmltZX19IiwyXSxbMSwyLCJcXG1hdGhiZntMfXRfMl5cXGFzdCAoXFx2YXJlcHNpbG9uKSJdLFszLDIsIlxcdmFyZXBzaWxvbl5cXHByaW1lX3tcXG1hdGhiZntMfXRfMl5cXGFzdH0iXSxbMCwzLCJcXFBoaV97XFxtYXRoYmZ7TH0odF5cXHByaW1lKV5cXGFzdCBLfShcXGJldGFee0teXFxwcmltZX0pIiwyXV0=
        \begin{tikzcd}
            {\Phi_{\mathbf{L}(t^\prime)^\ast K}\mathbf{L}t_1^\ast \Phi_{K^\prime}} & {\mathbf{L}t_2^\ast \Phi_K\Phi_{K^\prime}} \\
            {\Phi_{\mathbf{L}(t^\prime)^\ast K}\Phi_{\mathbf{L}(t^\prime)^\ast K^\prime}\mathbf{L}t_2^\ast} & {\mathbf{L}t_2^\ast.}
            \arrow["{\Phi_{\mathbf{L}(t^\prime)^\ast K}(\beta^{K^\prime})}"', from=1-1, to=2-1]
            \arrow["{\beta^K_{\Phi_{K^\prime}}}"', from=1-2, to=1-1]
            \arrow["{\mathbf{L}t_2^\ast (\varepsilon)}", from=1-2, to=2-2]
            \arrow["{\varepsilon^\prime_{\mathbf{L}t_2^\ast}}", from=2-1, to=2-2]
        \end{tikzcd}
	\end{equation}
	Also, the proof of \ref{ppty:reflection} never used that $t$ is affine.
	Specifically, we showed that with \Cref{setup:fm_cube_pullback} and given $K\in D^b_{\operatorname{coh}}(Y_1\times_S Y_2)$ relatively perfect over $Y_2$, it is the case that \ref{ppty:reflection} holds for $E\in D_{\operatorname{qc}}(Y_1)$.\footnote{\label{ftnt:ff_descent_hypotheses_on_bc_morph}The proof of \Cref{prop:pullbacks_is_a_morphism_of_adjoint_integral_transforms} only needs $t$ affine to guarantee that $\Phi_{(t^\prime)^\ast K}\dashv\Phi_{(t^\prime)^\ast K^\prime }$, so that \ref{ppty:reflection} plus a replacement of $\Phi_{(t^\prime)^\ast K}$ by a certain isomorphic functor yields a morphism of adjunctions as in the statement of \Cref{prop:pullbacks_is_a_morphism_of_adjoint_integral_transforms}.}
	Thus, if it also happens that $t_1^\ast \colon D_{\operatorname{qc}}(Y_1)\to D_{\operatorname{qc}}(Y_1\times_ST)$ is essentially surjective, then $\Phi_{(t^\prime)^\ast K^\prime }\colon D_{\operatorname{qc}}(Y_2\times_ST)\to D_{\operatorname{qc}}(Y_1\times_ST)$ has a left adjoint $L$ (see e.g. \cite[Theorem 3.1.5]{Borceaux:1994}) such that there is an isomorphism
	\begin{displaymath}
		L(A)\cong\Phi_{(t^\prime)^\ast K}(A)
	\end{displaymath}
	for each $A\in D_{\operatorname{qc}}(Y_1\times_ST)$.
	However, one should be careful, for the isomorphism above might \textit{not} be natural in $A$.
\end{remark}

\begin{lemma}
    \label{lem:Ballard_relative_perf_implies_dual_is_such}
    Let $f\colon Y\to X$ be a proper morphism between Noetherian schemes. 
    Consider an $f$-perfect pseudocoherent complex $E$ on $Y$. 
    Then $\operatorname{\mathbf{R}\mathcal{H}\! \mathit{om}}(E, f^! P)$ is $f$-perfect for all $P\in \operatorname{Perf}(X)$.
\end{lemma}

\begin{proof}
    By \cite[\href{https://stacks.math.columbia.edu/tag/0A9I}{Tag 0A9I}]{stacks-project}, we know that $f^\times P \in D^+_{\operatorname{qc}}(Y)$. 
    Since $E$ is pseudocoherent, it follows that $\operatorname{\mathbb{R}\mathcal{H}\! \mathit{om}}(E, f^\times P)\in D_{\operatorname{qc}}(Y)$ \cite[\href{https://stacks.math.columbia.edu/tag/0A6H}{Tag 0A6H}]{stacks-project}. 
    In this situation, we obtain that
    \begin{displaymath}
        \operatorname{\mathbf{R}\mathcal{H}\! \mathit{om}}(E, f^\times P) \cong \operatorname{\mathbb{R}\mathcal{H}\! \mathit{om}}(E, f^\times P).
    \end{displaymath}
    Choose $Q\in \operatorname{Perf}(Y)$. 
    From \cite[\href{https://stacks.math.columbia.edu/tag/08DQ}{Tags 08DQ} \& \href{https://stacks.math.columbia.edu/tag/0A8A}{0A8A}]{stacks-project}, the double duality morphism $Q\to (Q^\vee)^\vee$ is an isomorphism where $Q^\vee := \operatorname{\mathbf{R}\mathcal{H}\! \mathit{om}}(Q, \mathcal{O}_{Y})$, and there exists a natural isomorphism of functors $\operatorname{\mathbf{R}\mathcal{H}\! \mathit{om}}(Q, -)\cong \operatorname{\mathbb{R}\mathcal{H}\! \mathit{om}}(Q, -)$. 
    Then
    \begin{displaymath}
        \begin{aligned}
            \operatorname{\mathbb{R}\mathcal{H}\! \mathit{om}}(E, f^\times P) \otimes^{\mathbf{L}} Q 
            &\cong \operatorname{\mathbb{R}\mathcal{H}\! \mathit{om}}(E, f^\times P) \otimes^{\mathbf{L}} (Q^\vee)^\vee
            \\&\cong \operatorname{\mathbb{R}\mathcal{H}\! \mathit{om}} (Q^\vee, \operatorname{\mathbb{R}\mathcal{H}\! \mathit{om}}(E, f^\times P) ) && (\textrm{\cite[\href{https://stacks.math.columbia.edu/tag/08DQ}{Tag 08DQ}]{stacks-project}})
            \\&\cong \operatorname{\mathbb{R}\mathcal{H}\! \mathit{om}} (Q^\vee \otimes^{\mathbf{L}} E ,f^\times P) && (\textrm{\cite[\href{https://stacks.math.columbia.edu/tag/08DJ}{Tag 08DJ}]{stacks-project}}).
        \end{aligned}
    \end{displaymath}
    Applying \cite[\href{https://stacks.math.columbia.edu/tag/0A6H}{Tag 0A6H}]{stacks-project}, 
    \begin{displaymath}
        \operatorname{\mathbb{R}\mathcal{H}\! \mathit{om}} (Q^\vee \otimes^{\mathbf{L}} E ,f^\times P) \cong \operatorname{\mathbf{R}\mathcal{H}\! \mathit{om}} (Q^\vee \otimes^{\mathbf{L}} E ,f^\times P).
    \end{displaymath}
    Moreover, \cite[\href{https://stacks.math.columbia.edu/tag/0AU3}{Tag 0AU3}]{stacks-project} says $f^!$ is naturally isomorphic to $f^\times$ on $D^+_{\operatorname{qc}}(X)$, and there is an isomorphism
    \begin{displaymath}
        \mathbf{R}f_\ast \operatorname{\mathbf{R}\mathcal{H}\! \mathit{om}} (Q^\vee \otimes^{\mathbf{L}} E ,f^! P) \cong \operatorname{\mathbf{R}\mathcal{H}\! \mathit{om}}( \mathbf{R}f_\ast (Q^\vee \otimes^{\mathbf{L}} E), P).
    \end{displaymath}
    Since $Q^\vee$ is perfect, the hypothesis implies that $\mathbf{R}f_\ast (Q^\vee \otimes^{\mathbf{L}} E)$ is perfect.
    As $P$ is perfect, we conclude that $\operatorname{\mathbf{R}\mathcal{H}\! \mathit{om}}( \mathbf{R}f_\ast (Q^\vee \otimes^{\mathbf{L}} E), P)$ is perfect, which completes the proof.
\end{proof}

\begin{lemma}
    \label{lem:equivalences_induced}
    Let $f_1\colon Y_1\to S$ and $f_2\colon Y_2\to S$ be proper flat morphisms of Noetherian schemes. 
    Suppose $K\in D^b_{\operatorname{coh}}(Y_1\times_S Y_2)$ is relatively perfect over each $Y_i$. 
    Then the following are equivalent:
    \begin{enumerate}
        \item \label{lem:equivalences_induced1} $\Phi_K$ is fully faithful (resp.\ an equivalence) on $D_{\operatorname{qc}}$
        \item \label{lem:equivalences_induced2} $\Phi_K$ restricts to a fully faithful functor (resp.\ an equivalence) on $D^b_{\operatorname{coh}}$
        \item \label{lem:equivalences_induced3} $\Phi_K$ restricts to a fully faithful functor (resp.\ an equivalence) on $\operatorname{Perf}$.
    \end{enumerate}
    In particular, these conditions for fully faithfulness or equivalences can be tested on a single compact generator (see proof for a precise statement).
\end{lemma}

\begin{proof}
    Denote by $K^\prime := \operatorname{\mathbf{R}\mathcal{H}\! \mathit{om}} (K,f_1^! \mathcal{O}_{Y_2})$ the kernel of the integral transform obtained in \Cref{prop:right_adjoint_pullback} which is right adjoint to $\Phi_K$ on $D_{\operatorname{qc}}$.
    Denote by $K^{\prime \prime}=\operatorname{\mathbf{R}\mathcal{H}\! \mathit{om}} (K,f_2^! \mathcal{O}_{Y_1})$ the kernel of the integral transform obtained in \Cref{prop:left_adjoint_pullback} which is left adjoint to $\Phi_K$ on $D_{\operatorname{qc}}$.
    We start with fully faithfulness. 
    Since $K$ is relatively perfect over both $Y_i$, \Cref{lem:induced_dbcoh_perf_preservation_upon_pullback_with_t_affine} implies $\Phi_K$ restricts to a functor on $D^b_{\operatorname{coh}}$ and $\operatorname{Perf}$.
    Hence, $\eqref{lem:equivalences_induced1} \implies \eqref{lem:equivalences_induced2} \implies \eqref{lem:equivalences_induced3}$. 
    Assume \eqref{lem:equivalences_induced3} holds; that is, $\Phi_K$ restricts to a fully faithful functor on $\operatorname{Perf}$.
    Denote by $e\colon \Phi_{K^{\prime \prime}} \circ \Phi_K \to 1$ the counit of the adjoint pair $\Phi_{K^{\prime \prime}}$ and $\Phi_K$. 
    As $K$ is relatively perfect over each $Y_i$, \Cref{lem:Ballard_relative_perf_implies_dual_is_such} implies $K^{\prime \prime}$ is relatively perfect over $Y_1$. 
    Hence, \Cref{lem:induced_dbcoh_perf_preservation_upon_pullback_with_t_affine} says $\Phi_{K^{\prime \prime}}$ restricts to a functor on $\operatorname{Perf}$.
    Since $\Phi_K$ restricts to a fully faithful functor on $\operatorname{Perf}$, \cite[\href{https://stacks.math.columbia.edu/tag/07RB}{Tag 07RB}]{stacks-project} implies $e_P$ is an isomorphism for all $P\in \operatorname{Perf}(Y_1)$. 
    Consider the strictly full subcategory of objects $B\in D_{\operatorname{qc}}(Y_1)$ such that $e_B$ is an isomorphism. 
    It is straightforward to check that is a localizing subcategory.
    The work above shows $\operatorname{Perf}(Y_1)\subseteq \mathcal{T}$. 
    Since $D_{\operatorname{qc}}(Y_1)$ is compactly generated, \Cref{lem:localizing_iff_cpt_gen} implies $D_{\operatorname{qc}}(Y_1) = \mathcal{T}$. 
    Then \cite[\href{https://stacks.math.columbia.edu/tag/07RB}{Tag 07RB}]{stacks-project} implies $\Phi_K$ is fully faithful on $D_{\operatorname{qc}}$.

    Next we prove the case for equivalence.
    By \Cref{lem:equivalence_or_fully_faithful_via_compacts}, $\eqref{lem:equivalences_induced1} \iff \eqref{lem:equivalences_induced3}$. 
    Note that $\Phi_{K^{\prime \prime}}$ and $\Phi_K$ restrict to adjoint pair on $\operatorname{Perf}$. 
    By \Cref{cor:relatively_perfect_y2_implies_right_adjoint_restrict_to_dbcoh}, $\Phi_K$ and $\Phi_{K^\prime}$ restrict to an adjoint pair on $D^b_{\operatorname{coh}}$. 

    Assume $\eqref{lem:equivalences_induced1}$ holds. 
    Since the unit and counit the adjoint pair $\Phi_K$ and $\Phi_{K^\prime}$ for objects in $D^b_{\operatorname{coh}}$ remain bounded and pseudocoherent, 
    \cite[\href{https://stacks.math.columbia.edu/tag/07RB}{Tag 07RB}]{stacks-project} shows that $\eqref{lem:equivalences_induced1} \implies \eqref{lem:equivalences_induced2}$. 
    
    Assume \eqref{lem:equivalences_induced2} holds. 
    Since $\Phi_K$ restricts to a functor on $\operatorname{Perf}$, it is a fully faithful functor on $\operatorname{Perf}$. 
    By the fully faithful case, we know that $\Phi_K$ is fully faithful on $D_{\operatorname{qc}}$.
    Set $\mathcal{S}$ to be the strictly full subcategory of $B\in D_{\operatorname{qc}}(Y_2)$ such the counit $\epsilon_B \colon (\Phi_K \circ \Phi_{K^\prime})(B) \to B$ is an isomorphism. 
    It can be checked that $\mathcal{S}$ is localizing. 
    Note that the morphism $\epsilon_B$ lies in $D^b_{\operatorname{coh}}(Y_2)$ for all $B\in D^b_{\operatorname{coh}}(Y_2)$.
    Hence, from the assumption, it follows that $D^b_{\operatorname{coh}}(Y_2) \subseteq \mathcal{S}$.
    Thus, $\operatorname{Perf}(Y_2)\subseteq \mathcal{S}$, and so \Cref{lem:localizing_iff_cpt_gen} implies $\mathcal{S}=D_{\operatorname{qc}}(Y_2)$. 
    This shows that $\eqref{lem:equivalences_induced2} \implies \eqref{lem:equivalences_induced1}$.

    We prove the last claim.
    Write $\eta:1\to\Theta\circ\Phi$ and $\varepsilon:\Phi\circ\Theta\to 1$ for the unit and counit, respectively.
    For $i=1,2$, let $G_i$ be a compact generator of $D_{\operatorname{qc}}(Y_i)$, see \Cref{ex:perfect_compacts_classical_generator}.
    Then $\Phi$ is fully faithful (resp.\ and equivalence) if and only if $\eta_{G_1}$ is an isomorphism (resp.\ both $\eta_{G_1}$ and $\varepsilon_{G_2}$ are isomorphisms).
    Indeed, let $\mathcal{T}^\prime\subseteq\mathcal{T}$ and \ $\mathcal{S}^\prime\subseteq\mathcal{S}$ be the strictly full subcategories where $\eta$ and $\varepsilon$ are isomorphisms.
    One may verify that $\mathcal{T}^\prime$ and $\mathcal{S}^\prime$ are closed under shifts, iterated cones, and small coproducts (hence, homotopy colimits).
    Therefore, \cite[\href{https://stacks.math.columbia.edu/tag/09SN}{Tag 09SN}]{stacks-project} tells us that the condition implies $\mathcal{T} \subseteq \mathcal{T}^\prime$
    (resp.\ $\mathcal{T} \subseteq \mathcal{T}^\prime$ and $\mathcal{S} \subseteq \mathcal{S}^\prime$).
    This completes the proof.
\end{proof}

\begin{proposition}
	\label[proposition]{prop:descent}
	Consider \Cref{setup:fm_cube_pullback}. Let $K\in D^b_{\operatorname{coh}}(Y_1\times_S Y_2)$ be relatively perfect over each $Y_i$. If $\Phi_{\mathbf{L}(t^\prime)^\ast K}$ is fully faithful (resp.\ an equivalence) on $D^b_{\operatorname{coh}}$ where $t$ is one of the following:
    \begin{enumerate}[label=(\arabic*), ref=\theremark(\arabic*)]
        \item \label[proposition]{prop:descent1} $t$ is affine and faithfully flat (e.g.\ $t=1_S$)
        \item \label[proposition]{prop:descent2} $t\colon \operatorname{Spec}(k)\to S$ is any morphism from a field (we require the condition to hold for all for all such morphisms)
        \item \label[proposition]{prop:descent2_closed} $t\colon \operatorname{Spec}(\kappa(p))\to S$ is any closed immersion associated to a closed point $p\in S$ (we require the condition to hold for all for all such morphisms),
    \end{enumerate}    
	%%NOTE: by same result, $\Phi_{\mathbf{L}(t^\prime)^\ast K^\prime }$ is right adjoint to $\Phi_{\mathbf{L}(t^\prime)^\ast K}$)
    then so is $\Phi_{K}$.
\end{proposition}

\begin{proof}
	Set $K^\prime$ to be the kernel of the integral transform obtained in \Cref{prop:right_adjoint_pullback} which is right adjoint to $\Phi_K$ on $D_{\operatorname{qc}}$. By \Cref{thm:bounded_pseudocoherence_perfectness_faithfully_flat_affine}, $\mathbf{L}(t^\prime)^\ast K$ is relatively perfect over $Y_2\times_S T$ for each $t\colon T \to S$ in the statement. Hence, \Cref{prop:right_adjoint_pullback} tells us that $\Phi_{\mathbf{L}(t^\prime)^\ast K^\prime}$ is right adjoint to $\Phi_{\mathbf{L}(t^\prime)^\ast K}$ on $D_{\operatorname{qc}}$ for each $t$ in the statement. So, from \Cref{cor:relatively_perfect_y2_implies_right_adjoint_restrict_to_dbcoh}, we know that $\Phi_K \dashv \Phi_{K^\prime}$ and $\Phi_{\mathbf{L}(t^\prime)^\ast K} \dashv \Phi_{\mathbf{L}(t^\prime)^\ast K^\prime}$ on $D^b_{\operatorname{coh}}$ where $t$ is any morphism in the statement.

    We start with the first case. Denote $\eta$ and $\varepsilon$ (resp.\ $\eta^\prime $, $\varepsilon^\prime $) for the unit and counit of the adjunction $\Phi_K\dashv\Phi_{K^\prime }$ (resp.\ of $\Phi_{(t^\prime)^\ast K}\dashv\Phi_{(t^\prime)^\ast K^\prime }$). Since $t$ is faithfully flat, so are each $t_i$; hence each $\mathbf{L}(t_i)^\ast $ reflects isomorphisms. Therefore, by commutativity of \eqref{eq:unit_of_integral_transforms_downstairs_and_upstairs_are_compatible} (resp.\ by \eqref{eq:counit_of_integral_transforms_downstairs_and_upstairs_are_compatible}),
	it follows that $\eta^\prime $ (resp.\ $\varepsilon^\prime $) being an isomorphism implies $\eta$ (resp.\ $\varepsilon$) is an isomorphism.
	The result then follows from \cite[\href{https://stacks.math.columbia.edu/tag/07RB}{Tag 07RB}]{stacks-project}.
    %%NOTE: See also \cite[\href{https://ncatlab.org/nlab/show/adjoint+functor\#FullyFaithfulAndInvertibleAdjoints}{Proposition 3.4}]{nlab}.
    Indeed, we check the case for fully faithful.
    Choose $E\in D^b_{\operatorname{coh}}(Y_1)$.
    Applying \eqref{eq:unit_of_integral_transforms_downstairs_and_upstairs_are_compatible}, we can form the morphism of distinguished triangles 
    \begin{displaymath}
        % https://q.uiver.app/#q=WzAsOCxbMCwwLCJcXG1hdGhiZntMfSB0XzFeXFxhc3QgRSJdLFsxLDEsIlxcbWF0aGJme0x9IHRfMV5cXGFzdFxcUGhpX3tLXlxccHJpbWV9XFxQaGlfSyAoRSkiXSxbMSwwLCJcXFBoaV97XFxtYXRoYmZ7TH0odF5cXHByaW1lKV5cXGFzdCBLXlxccHJpbWV9XFxQaGlfe1xcbWF0aGJme0x9KHReXFxwcmltZSleXFxhc3QgS30gXFxtYXRoYmZ7TH10XzFeXFxhc3QgRSJdLFswLDEsIlxcbWF0aGJme0x9IHRfMV5cXGFzdCBFIl0sWzIsMCwiXFxvcGVyYXRvcm5hbWV7Y29uZX0oXFxldGFeXFxwcmltZV97XFxtYXRoYmZ7TH10XzFeXFxhc3QgRX0pIl0sWzIsMSwiXFxtYXRoYmZ7TH0gdF8xXlxcYXN0IChcXG9wZXJhdG9ybmFtZXtjb25lfShcXGV0YV9FKSkiXSxbMywwLCJcXG1hdGhiZntMfSB0XzFeXFxhc3QgRSBbMV0iXSxbMywxLCJcXG1hdGhiZntMfSB0XzFeXFxhc3QgRSBbMV0uIl0sWzAsMiwiXFxldGFeXFxwcmltZV97XFxtYXRoYmZ7TH10XzFeXFxhc3QgRX0iXSxbMywxLCJcXG1hdGhiZntMfSB0XzFeXFxhc3QgKFxcZXRhX0UpIiwyXSxbMSwyLCJcXFBoaV97XFxtYXRoYmZ7TH0odF5cXHByaW1lKV5cXGFzdCBLXlxccHJpbWV9KFxcYmV0YV5LKSBcXGNpcmMgXFxiZXRhXntLXlxccHJpbWV9X3tcXFBoaV9LfSIsMl0sWzMsMCwiMSJdLFsyLDRdLFsxLDVdLFs0LDZdLFs1LDRdLFs1LDddLFs3LDZdXQ==
        \begin{tikzcd}
            {\mathbf{L} t_1^\ast E} & {\Phi_{\mathbf{L}(t^\prime)^\ast K^\prime}\Phi_{\mathbf{L}(t^\prime)^\ast K} \mathbf{L}t_1^\ast E} & {\operatorname{cone}(\eta^\prime_{\mathbf{L}t_1^\ast E})} & {\mathbf{L} t_1^\ast E [1]} \\
            {\mathbf{L} t_1^\ast E} & {\mathbf{L} t_1^\ast\Phi_{K^\prime}\Phi_K (E)} & {\mathbf{L} t_1^\ast (\operatorname{cone}(\eta_E))} & {\mathbf{L} t_1^\ast E [1].}
            \arrow["{\eta^\prime_{\mathbf{L}t_1^\ast E}}", from=1-1, to=1-2]
            \arrow[from=1-2, to=1-3]
            \arrow[from=1-3, to=1-4]
            \arrow["1", from=2-1, to=1-1]
            \arrow["{\mathbf{L} t_1^\ast (\eta_E)}"', from=2-1, to=2-2]
            \arrow["{\Phi_{\mathbf{L}(t^\prime)^\ast K^\prime}(\beta^K) \circ \beta^{K^\prime}_{\Phi_K}}"', from=2-2, to=1-2]
            \arrow[from=2-2, to=2-3]
            \arrow[from=2-3, to=1-3]
            \arrow[from=2-3, to=2-4]
            \arrow[from=2-4, to=1-4]
        \end{tikzcd}
    \end{displaymath}
    Since the left most and middle morphism is an isomorphism, we obtain 
    \begin{displaymath}
        \operatorname{cone}(\eta^\prime_{\mathbf{L}t_1^\ast E}) \cong \mathbf{L} t_1^\ast (\operatorname{cone}(\eta_E)).
    \end{displaymath}
    As $\Phi_{\mathbf{L}(t^\prime)^\ast K}$ is fully faithful on $D^b_{\operatorname{coh}}$, it follows that 
    \begin{displaymath}
        0 \cong \operatorname{cone}(\eta^\prime_{\mathbf{L}t_1^\ast E}) \cong \mathbf{L} t_1^\ast (\operatorname{cone}(\eta_E)).
    \end{displaymath}
    Since $t$ is faithfully flat, so is $t_1$. 
    Hence, $\operatorname{cone}(\eta_E)\cong 0$, which yields the claim.
    
    Now, we check the second and third case. 
    In both cases, $\Phi_{\mathbf{L}(t^\prime)^\ast K}$ is fully faithful (resp.\ an equivalence) on $D^b_{\operatorname{coh}}$ with $t\colon \operatorname{Spec}(\kappa(p))\to S$ is any closed immersion associated to a closed point $p\in S$. 
    We prove only fully faithfulness because equivalences is similar by argueing with the counit.
    In fact, the cone of the counit for the adjoint pair on $D^b_{\operatorname{coh}}$ is a bounded pseudocoherent complex because $K$ is relatively perfect over each $Y_i$. 
    This ingredient allows for the counit case to be proved similarly like below.

    To prove $\Phi_K$ is fully faithful, we have to show that the unit $\eta \colon 1 \to \Phi_{K^\prime} \circ \Phi_K$ is an isomorphism. 
    Pick $E\in D^b_{\operatorname{coh}}(X)$. 
    Consider the distinguished triangle 
    \begin{displaymath}
        E \xrightarrow{\eta_E} (\Phi_{K^\prime} \circ \Phi_K)(E) \to \operatorname{cone}(\eta_E) \to E[1].
    \end{displaymath}
    Choose a closed point $p\in S$. 
    Applying \eqref{eq:unit_of_integral_transforms_downstairs_and_upstairs_are_compatible}, we can form the morphism of distinguished triangles 
    \begin{displaymath}
        % https://q.uiver.app/#q=WzAsOCxbMCwwLCJcXG1hdGhiZntMfSB0XzFeXFxhc3QgRSJdLFsxLDEsIlxcbWF0aGJme0x9IHRfMV5cXGFzdFxcUGhpX3tLXlxccHJpbWV9XFxQaGlfSyAoRSkiXSxbMSwwLCJcXFBoaV97XFxtYXRoYmZ7TH0odF5cXHByaW1lKV5cXGFzdCBLXlxccHJpbWV9XFxQaGlfe1xcbWF0aGJme0x9KHReXFxwcmltZSleXFxhc3QgS30gXFxtYXRoYmZ7TH10XzFeXFxhc3QgRSJdLFswLDEsIlxcbWF0aGJme0x9IHRfMV5cXGFzdCBFIl0sWzIsMCwiXFxvcGVyYXRvcm5hbWV7Y29uZX0oXFxldGFeXFxwcmltZV97XFxtYXRoYmZ7TH10XzFeXFxhc3QgRX0pIl0sWzIsMSwiXFxtYXRoYmZ7TH0gdF8xXlxcYXN0IChcXG9wZXJhdG9ybmFtZXtjb25lfShcXGV0YV9FKSkiXSxbMywwLCJcXG1hdGhiZntMfSB0XzFeXFxhc3QgRSBbMV0iXSxbMywxLCJcXG1hdGhiZntMfSB0XzFeXFxhc3QgRSBbMV0uIl0sWzAsMiwiXFxldGFeXFxwcmltZV97XFxtYXRoYmZ7TH10XzFeXFxhc3QgRX0iXSxbMywxLCJcXG1hdGhiZntMfSB0XzFeXFxhc3QgKFxcZXRhX0UpIiwyXSxbMSwyLCJcXFBoaV97XFxtYXRoYmZ7TH0odF5cXHByaW1lKV5cXGFzdCBLXlxccHJpbWV9KFxcYmV0YV5LKSBcXGNpcmMgXFxiZXRhXntLXlxccHJpbWV9X3tcXFBoaV9LfSIsMl0sWzMsMCwiMSJdLFsyLDRdLFsxLDVdLFs0LDZdLFs1LDRdLFs1LDddLFs3LDZdXQ==
        \begin{tikzcd}
            {\mathbf{L} t_1^\ast E} & {\Phi_{\mathbf{L}(t^\prime)^\ast K^\prime}\Phi_{\mathbf{L}(t^\prime)^\ast K} \mathbf{L}t_1^\ast E} & {\operatorname{cone}(\eta^\prime_{\mathbf{L}t_1^\ast E})} & {\mathbf{L} t_1^\ast E [1]} \\
            {\mathbf{L} t_1^\ast E} & {\mathbf{L} t_1^\ast\Phi_{K^\prime}\Phi_K (E)} & {\mathbf{L} t_1^\ast (\operatorname{cone}(\eta_E))} & {\mathbf{L} t_1^\ast E [1].}
            \arrow["{\eta^\prime_{\mathbf{L}t_1^\ast E}}", from=1-1, to=1-2]
            \arrow[from=1-2, to=1-3]
            \arrow[from=1-3, to=1-4]
            \arrow["1", from=2-1, to=1-1]
            \arrow["{\mathbf{L} t_1^\ast (\eta_E)}"', from=2-1, to=2-2]
            \arrow["{\Phi_{\mathbf{L}(t^\prime)^\ast K^\prime}(\beta^K) \circ \beta^{K^\prime}_{\Phi_K}}"', from=2-2, to=1-2]
            \arrow[from=2-2, to=2-3]
            \arrow[from=2-3, to=1-3]
            \arrow[from=2-3, to=2-4]
            \arrow[from=2-4, to=1-4]
        \end{tikzcd}
    \end{displaymath}
    Since the left most and middle morphism is an isomorphism, we obtain 
    \begin{displaymath}
        \operatorname{cone}(\eta^\prime_{\mathbf{L}t_1^\ast E}) \cong \mathbf{L} t_1^\ast (\operatorname{cone}(\eta_E)).
    \end{displaymath}
    By \Cref{lem:equivalences_induced}, $\Phi_{\mathbf{L}(t^\prime)^\ast K}$ induces an equivalence on $D_{\operatorname{qc}}$, and so 
    \begin{displaymath}
        0 \cong \operatorname{cone}(\eta^\prime_{\mathbf{L}t_1^\ast E}) \cong \mathbf{L} t_1^\ast (\operatorname{cone}(\eta_E)).
    \end{displaymath}
    However, $\operatorname{cone}(\eta_E)\in D^b_{\operatorname{coh}}(Y_1)$, which implies $\mathbf{L} t_1^\ast (\operatorname{cone}(\eta_E))\in D^-_{\operatorname{coh}}(Y_1\times_S T)$.
    This shows that the derived pullback of $\operatorname{cone}(\eta_E)$ along $\operatorname{Spec}(\kappa(s))\to S$ for each closed point $s\in S$ is zero.
    Since $f_1$ is a closed map, by \Cref{lm:fiberwise_vanishing_implies_vanishing}, we conclude $\operatorname{cone}(\eta_E)\cong 0$, i.e.\ $\eta_E$ is an isomorphism.
    The proof is finished.
\end{proof}

While \Cref{prop:pullforwards_is_a_morphism_of_adjoint_integral_transforms} and \Cref{rmk:pullforwards_is_a_morphism_of_adjoint_integral_transforms} are not required, we record them for the sake of interest.
See also \Cref{rem:alt_proof_ascent}.

\begin{proposition}
    \label{prop:pullforwards_is_a_morphism_of_adjoint_integral_transforms}
    Consider \Cref{setup:fm_cube_pullback} where $t$ is affine. Let $K\in D^b_{\operatorname{coh}}(Y_1\times_S Y_2)$ be relatively perfect over $Y_2$. Denote by $K^\prime$ for the kernel of the integral transform obtained in \Cref{prop:right_adjoint_pullback} which is right adjoint to $\Phi_K$ on $D_{\operatorname{qc}}$.
    %%NOTE: by same result, $\Phi_{\mathbf{L}(t^\prime)^\ast K^\prime}$ is right adjoint to $\Phi_{\mathbf{L}(t^\prime)^\ast K}$).
    Then, after possibly replacing $\Phi_{K^\prime}$ by an isomorphic functor, we have a morphism of adjoint pairs as depicted in the diagram
    \begin{displaymath}
        % https://q.uiver.app/#q=WzAsNCxbMSwwLCJEX3tcXG9wZXJhdG9ybmFtZXtxY319KFlfMSkiXSxbMSwxLCJEX3tcXG9wZXJhdG9ybmFtZXtxY319KFlfMikiXSxbMCwwLCJEX3tcXG9wZXJhdG9ybmFtZXtxY319KFlfMVxcdGltZXNfU1QpIl0sWzAsMSwiRF97XFxvcGVyYXRvcm5hbWV7cWN9fShZXzJcXHRpbWVzX1NUKSJdLFsyLDAsIlxcbWF0aGJme1J9KHRfMSlfXFxhc3QiXSxbMywxLCJcXG1hdGhiZntSfSh0XzIpX1xcYXN0IiwyXSxbMSwwLCJcXFBoaV97S15cXHByaW1lfSIsMix7Im9mZnNldCI6Mn1dLFswLDEsIlxcUGhpX0siLDIseyJvZmZzZXQiOjJ9XSxbMywyLCJcXFBoaV97XFxtYXRoYmZ7TH0odF5cXHByaW1lKV5cXGFzdCBLXlxccHJpbWV9IiwyLHsib2Zmc2V0IjoyfV0sWzIsMywiXFxQaGlfe1xcbWF0aGJme0x9KHReXFxwcmltZSleXFxhc3QgS30iLDIseyJvZmZzZXQiOjJ9XSxbNyw2LCIiLDIseyJsZXZlbCI6MSwic3R5bGUiOnsibmFtZSI6ImFkanVuY3Rpb24ifX1dLFs5LDgsIiIsMix7ImxldmVsIjoxLCJzdHlsZSI6eyJuYW1lIjoiYWRqdW5jdGlvbiJ9fV1d
        \begin{tikzcd}[column sep=3em]
            {D_{\operatorname{qc}}(Y_1\times_ST)} & {D_{\operatorname{qc}}(Y_1)} \\
            {D_{\operatorname{qc}}(Y_2\times_ST)} & {D_{\operatorname{qc}}(Y_2)}
            \arrow["{\mathbf{R}(t_1)_\ast}", from=1-1, to=1-2]
            \arrow[""{name=0, anchor=center, inner sep=0}, "{\Phi_{\mathbf{L}(t^\prime)^\ast K}}"', shift right=2, from=1-1, to=2-1]
            \arrow[""{name=1, anchor=center, inner sep=0}, "{\Phi_K}"', shift right=2, from=1-2, to=2-2]
            \arrow[""{name=2, anchor=center, inner sep=0}, "{\Phi_{\mathbf{L}(t^\prime)^\ast K^\prime}}"', shift right=2, from=2-1, to=1-1]
            \arrow["{\mathbf{R}(t_2)_\ast}"', from=2-1, to=2-2]
            \arrow[""{name=3, anchor=center, inner sep=0}, "{\Phi_{K^\prime}}"', shift right=2, from=2-2, to=1-2]
            \arrow["\dashv"{anchor=center}, draw=none, from=0, to=2]
            \arrow["\dashv"{anchor=center}, draw=none, from=1, to=3]
        \end{tikzcd}
    \end{displaymath}
    with left and right comparison transformations given respectively by the isomorphisms
    \begin{displaymath}
        \begin{aligned}
            \alpha^K \colon \Phi_K\circ \mathbf{R}(t_1)_\ast \to
            \mathbf{R}(t_2)_\ast \circ\Phi_{\mathbf{L}(t^\prime)^\ast K},\\
            (\alpha^{K^\prime})^{-1}\colon \mathbf{R}(t_1)_\ast \circ\Phi_{\mathbf{L}(t^\prime)^\ast K^\prime}\to
            \Phi_{K^\prime}\circ \mathbf{R}(t_2)^\ast.
        \end{aligned}
    \end{displaymath}
\end{proposition}

\begin{proof}
    This argument is essentially dual to that of \Cref{prop:pullbacks_is_a_morphism_of_adjoint_integral_transforms}.
\end{proof}

\begin{remark}
    \label{rmk:pullforwards_is_a_morphism_of_adjoint_integral_transforms}
    Consider the situation as in \Cref{prop:pullforwards_is_a_morphism_of_adjoint_integral_transforms}.
	Denote $\varepsilon$ and $\varepsilon^\prime$ respectively the counits of the adjunctions $\Phi_K\dashv\Phi_{K^\prime}$ and $\Phi_{\mathbf{L}(t^\prime)^\ast K}\dashv\Phi_{\mathbf{L}(t^\prime)^\ast K^\prime}$.
	From \Cref{lem:morph_adj},
	it follows that (after possibly replacing $\Phi_{K^\prime}$ by an isomorphic functor) the following diagrams commute:
    \begin{equation}
        \label{eq:counit_of_integral_transforms_downstairs_and_upstairs_are_pushforward_compatible}
            % https://q.uiver.app/#q=WzAsNCxbMCwwLCJcXFBoaV9LXFxtYXRoYmZ7TH0odF8xKV9cXGFzdFxcUGhpX3tcXG1hdGhiZntMfSh0XlxccHJpbWUpXlxcYXN0IEteXFxwcmltZX0iXSxbMCwxLCJcXFBoaV9LXFxQaGlfe0teXFxwcmltZX1cXG1hdGhiZntSfSh0XzIpX1xcYXN0Il0sWzEsMSwiXFxtYXRoYmZ7Un0odF8yKV9cXGFzdCJdLFsxLDAsIlxcbWF0aGJme1J9KHRfMilfXFxhc3QgXFxQaGlfe1xcbWF0aGJme0x9KHReXFxwcmltZSleXFxhc3QgS31cXFBoaV97XFxtYXRoYmZ7TH0odF5cXHByaW1lKV5cXGFzdCBLXlxccHJpbWV9Il0sWzAsMywiXFxhbHBoYV5LX3tcXFBoaV97XFxtYXRoYmZ7TH0odF5cXHByaW1lKV5cXGFzdCBLXlxccHJpbWV9fSJdLFsxLDAsIlxcUGhpX0soXFxhbHBoYV57S15cXHByaW1lfSkiXSxbMSwyLCJcXHZhcmVwc2lsb25fe1xcbWF0aGJme1J9KHRfMilfXFxhc3R9IiwyXSxbMywyLCJcXG1hdGhiZntSfSh0XzIpX1xcYXN0IChcXHZhcmVwc2lsb25eXFxwcmltZSkiXV0=
            \begin{tikzcd}
                {\Phi_K\mathbf{L}(t_1)_\ast\Phi_{\mathbf{L}(t^\prime)^\ast K^\prime}} & {\mathbf{R}(t_2)_\ast \Phi_{\mathbf{L}(t^\prime)^\ast K}\Phi_{\mathbf{L}(t^\prime)^\ast K^\prime}} \\
                {\Phi_K\Phi_{K^\prime}\mathbf{R}(t_2)_\ast} & {\mathbf{R}(t_2)_\ast}
                \arrow["{\alpha^K_{\Phi_{\mathbf{L}(t^\prime)^\ast K^\prime}}}", from=1-1, to=1-2]
                \arrow["{\mathbf{R}(t_2)_\ast (\varepsilon^\prime)}", from=1-2, to=2-2]
                \arrow["{\Phi_K(\alpha^{K^\prime})}", from=2-1, to=1-1]
                \arrow["{\varepsilon_{\mathbf{R}(t_2)_\ast}}"', from=2-1, to=2-2]
            \end{tikzcd}
    \end{equation}
    \begin{equation}
    \label{eq:unit_of_integral_transforms_downstairs_and_upstairs_are_pushforward_compatible}
        % https://q.uiver.app/#q=WzAsNCxbMCwwLCJcXG1hdGhiZntSfSh0XzEpX1xcYXN0Il0sWzAsMSwiXFxtYXRoYmZ7Un0odF8xKV9cXGFzdFxcUGhpX3tcXG1hdGhiZntMfSh0XlxccHJpbWUpXlxcYXN0IEteXFxwcmltZX1cXFBoaV97XFxtYXRoYmZ7TH0odF5cXHByaW1lKV5cXGFzdCBLfSJdLFsxLDAsIlxcUGhpX3tLXlxccHJpbWV9XFxQaGlfS1xcbWF0aGJme1J9KHRfMSlfXFxhc3QiXSxbMSwxLCJcXFBoaV97S15cXHByaW1lfVxcbWF0aGJme1J9KHRfMilfXFxhc3QgXFxQaGlfe1xcbWF0aGJme0x9KHReXFxwcmltZSleXFxhc3QgS30iXSxbMCwyLCJcXGV0YV97XFxtYXRoYmZ7Un0odF8xKV9cXGFzdH0iXSxbMywxLCJcXGFscGhhXntLXlxccHJpbWV9X3tcXFBoaV97XFxtYXRoYmZ7TH0odF5cXHByaW1lKV5cXGFzdCBLfX0iXSxbMiwzLCJcXFBoaV97S15cXHByaW1lfShcXGFscGhhXkspIl0sWzAsMSwiXFxtYXRoYmZ7Un0odF8xKV9cXGFzdChcXGV0YV5cXHByaW1lKSIsMl1d
        \begin{tikzcd}
            {\mathbf{R}(t_1)_\ast} &[1em] {\Phi_{K^\prime}\Phi_K\mathbf{R}(t_1)_\ast} \\
        	{\mathbf{R}(t_1)_\ast\Phi_{\mathbf{L}(t^\prime)^\ast K^\prime}\Phi_{\mathbf{L}(t^\prime)^\ast K}} & {\Phi_{K^\prime}\mathbf{R}(t_2)_\ast \Phi_{\mathbf{L}(t^\prime)^\ast K}.}
        	\arrow["{\eta_{\mathbf{R}(t_1)_\ast}}", from=1-1, to=1-2]
        	\arrow["{\mathbf{R}(t_1)_\ast(\eta^\prime)}"', from=1-1, to=2-1]
        	\arrow["{\Phi_{K^\prime}(\alpha^K)}", from=1-2, to=2-2]
        	\arrow["{\alpha^{K^\prime}_{\Phi_{\mathbf{L}(t^\prime)^\ast K}}}", from=2-2, to=2-1]
        \end{tikzcd}
    \end{equation}
\end{remark}

\begin{proposition}
    \label{prop:ascending}
    Consider \Cref{setup:fm_cube_pullback} where $t$ is affine. Let $K\in D^b_{\operatorname{coh}}(Y_1\times_S Y_2)$ be relatively perfect over each  $Y_i$. If $\Phi_K$ is fully faithful (resp.\ an equivalence) on $D^b_{\operatorname{coh}}$, then so is $\Phi_{\mathbf{L}(t^\prime)^\ast K}$.
\end{proposition}

\begin{proof}
    Denote by $K^\prime$ for the kernel of the integral transform obtained in \Cref{prop:right_adjoint_pullback} which is right adjoint to $\Phi_K$ on $D_{\operatorname{qc}}$. By \Cref{thm:bounded_pseudocoherence_perfectness_faithfully_flat_affine}, $\mathbf{L}(t^\prime)^\ast K$ is relatively perfect over $Y_2\times_S T$. Hence, \Cref{prop:right_adjoint_pullback} tells us that $\Phi_{\mathbf{L}(t^\prime)^\ast K^\prime}$ is right adjoint to $\Phi_{\mathbf{L}(t^\prime)^\ast K}$ on $D_{\operatorname{qc}}$. So, from \Cref{cor:relatively_perfect_y2_implies_right_adjoint_restrict_to_dbcoh}, we know that $\Phi_K \dashv \Phi_{K^\prime}$ and $\Phi_{\mathbf{L}(t^\prime)^\ast K} \dashv \Phi_{\mathbf{L}(t^\prime)^\ast K^\prime}$ on $D^b_{\operatorname{coh}}$. Now, let $\eta$ and $\varepsilon$ (resp. $\eta^\prime$ and $\varepsilon^\prime$) be the unit and counit of the adjunction $\Phi_K\dashv\Phi_{K^\prime}$ (resp. of $\Phi_{\mathbf{L}(t^\prime)^\ast K}\dashv\Phi_{\mathbf{L}(t^\prime)^\ast K^\prime}$).
	Since $t$ is affine, base changes implies $t_1,t_2, t^\prime$ are affine and \cite[\href{https://stacks.math.columbia.edu/tag/08I8}{Tag 08I8}]{stacks-project} ensures their derived pullbacks preserves compact generators.

    We prove only fully faithfulness because equivalences is similar by argueing with the counit.
    In fact, the cone of the counit for the adjoint pairs on $D^b_{\operatorname{coh}}$ is a bounded pseudocoherent complex because $K$ is relatively perfect over each $Y_i$. 
    This ingredient allows for the counit case to be proved similarly like below.
    The case of the counit applies \eqref{eq:counit_of_integral_transforms_downstairs_and_upstairs_are_compatible}, which the argument belows detect vanishing of the cones.

	% Therefore, by commutativity of \eqref{eq:unit_of_integral_transforms_downstairs_and_upstairs_are_pushforward_compatible} (resp.\ by \eqref{eq:counit_of_integral_transforms_downstairs_and_upstairs_are_pushforward_compatible}),
    % from \cite[\href{https://stacks.math.columbia.edu/tag/0BQT}{Tag 0BQT}]{stacks-project}
	% we get that $\eta$ (resp.\ $\varepsilon$) being an isomorphism implies $\eta^\prime$ (resp.\ $\varepsilon^\prime$) is an isomorphism.
    % Applying this to a compact generator, we finish the proof by \Cref{lem:equivalences_induced}.
    % Indeed, we show the case for fully faithfulness.

    Let $G$ be a compact generator for $D_{\operatorname{qc}}(Y_1)$. 
    By \Cref{lem:equivalences_induced}, $\Phi_K$ is fully faithful on $D_{\operatorname{qc}}$.
    Now, \cite[\href{https://stacks.math.columbia.edu/tag/0BQT}{Tag 0BQT}]{stacks-project} says $\mathbf{L}t_1^\ast G$ is a compact generator for $D_{\operatorname{qc}}(Y_1 \times_S T)$.
    Applying \eqref{eq:unit_of_integral_transforms_downstairs_and_upstairs_are_compatible}, we can form the morphism of distinguished triangles 
    \begin{displaymath}
        % https://q.uiver.app/#q=WzAsOCxbMCwwLCJcXG1hdGhiZntMfSB0XzFeXFxhc3QgRSJdLFsxLDEsIlxcbWF0aGJme0x9IHRfMV5cXGFzdFxcUGhpX3tLXlxccHJpbWV9XFxQaGlfSyAoRSkiXSxbMSwwLCJcXFBoaV97XFxtYXRoYmZ7TH0odF5cXHByaW1lKV5cXGFzdCBLXlxccHJpbWV9XFxQaGlfe1xcbWF0aGJme0x9KHReXFxwcmltZSleXFxhc3QgS30gXFxtYXRoYmZ7TH10XzFeXFxhc3QgRSJdLFswLDEsIlxcbWF0aGJme0x9IHRfMV5cXGFzdCBFIl0sWzIsMCwiXFxvcGVyYXRvcm5hbWV7Y29uZX0oXFxldGFeXFxwcmltZV97XFxtYXRoYmZ7TH10XzFeXFxhc3QgRX0pIl0sWzIsMSwiXFxtYXRoYmZ7TH0gdF8xXlxcYXN0IChcXG9wZXJhdG9ybmFtZXtjb25lfShcXGV0YV9FKSkiXSxbMywwLCJcXG1hdGhiZntMfSB0XzFeXFxhc3QgRSBbMV0iXSxbMywxLCJcXG1hdGhiZntMfSB0XzFeXFxhc3QgRSBbMV0uIl0sWzAsMiwiXFxldGFeXFxwcmltZV97XFxtYXRoYmZ7TH10XzFeXFxhc3QgRX0iXSxbMywxLCJcXG1hdGhiZntMfSB0XzFeXFxhc3QgKFxcZXRhX0UpIiwyXSxbMSwyLCJcXFBoaV97XFxtYXRoYmZ7TH0odF5cXHByaW1lKV5cXGFzdCBLXlxccHJpbWV9KFxcYmV0YV5LKSBcXGNpcmMgXFxiZXRhXntLXlxccHJpbWV9X3tcXFBoaV9LfSIsMl0sWzMsMCwiMSJdLFsyLDRdLFsxLDVdLFs0LDZdLFs1LDRdLFs1LDddLFs3LDZdXQ==
        \begin{tikzcd}
            {\mathbf{L} t_1^\ast G} & {\Phi_{\mathbf{L}(t^\prime)^\ast K^\prime}\Phi_{\mathbf{L}(t^\prime)^\ast K} \mathbf{L}t_1^\ast G} & {\operatorname{cone}(\eta^\prime_{\mathbf{L}t_1^\ast G})} & {\mathbf{L} t_1^\ast G [1]} \\
            {\mathbf{L} t_1^\ast G} & {\mathbf{L} t_1^\ast\Phi_{K^\prime}\Phi_K (G)} & {\mathbf{L} t_1^\ast (\operatorname{cone}(\eta_G))} & {\mathbf{L} t_1^\ast G [1].}
            \arrow["{\eta^\prime_{\mathbf{L}t_1^\ast G}}", from=1-1, to=1-2]
            \arrow[from=1-2, to=1-3]
            \arrow[from=1-3, to=1-4]
            \arrow["1", from=2-1, to=1-1]
            \arrow["{\mathbf{L} t_1^\ast (\eta_G)}"', from=2-1, to=2-2]
            \arrow["{\Phi_{\mathbf{L}(t^\prime)^\ast K^\prime}(\beta^K) \circ \beta^{K^\prime}_{\Phi_K}}"', from=2-2, to=1-2]
            \arrow[from=2-2, to=2-3]
            \arrow[from=2-3, to=1-3]
            \arrow[from=2-3, to=2-4]
            \arrow[from=2-4, to=1-4]
        \end{tikzcd}
    \end{displaymath}
    Since the left most and middle morphism is an isomorphism, we obtain 
    \begin{displaymath}
        \operatorname{cone}(\eta^\prime_{\mathbf{L}t_1^\ast G}) \cong \mathbf{L} t_1^\ast (\operatorname{cone}(\eta_G)).
    \end{displaymath}
    As $\Phi_{K}$ is fully faithful on $D^b_{\operatorname{coh}}$, it follows that 
    \begin{displaymath}
        \operatorname{cone}(\eta^\prime_{\mathbf{L}t_1^\ast G}) \cong \mathbf{L} t_1^\ast (\operatorname{cone}(\eta_G))\cong 0.
    \end{displaymath}
    By \Cref{lem:equivalences_induced}, $\Phi_{\mathbf{L}(t^\prime)^\ast K}$ is fully faithful on $D_{\operatorname{qc}}$, and hence on $D^b_{\operatorname{coh}}$.
\end{proof}

\begin{remark}
    \label{rem:alt_proof_ascent}
    There is an alternative proof of \Cref{prop:ascending} in a similar spirit to the proof of the first case of \Cref{prop:descent},
    replacing in the argument \Cref{prop:pullbacks_is_a_morphism_of_adjoint_integral_transforms} by \Cref{prop:pullforwards_is_a_morphism_of_adjoint_integral_transforms}.
    Indeed, by \Cref{lem:equivalences_induced}, we might as well show a version of \Cref{prop:ascending} for full faithfulness (resp. equivalence) on $D_{\operatorname{qc}}$.
    Therefore, since $\mathbf{R}t_*$ is conservative \cite[\href{https://stacks.math.columbia.edu/tag/08I8}{Tag 08I8}]{stacks-project},
    it suffices to see $\mathbf{R}(t_1)_*(\eta^\prime)$ is an iso (resp. $\mathbf{R}(t_1)_*(\eta^\prime)$ and $\mathbf{R}(t_2)_*(\varepsilon^\prime)$ are isos).
    This follows from commutativity of \eqref{eq:counit_of_integral_transforms_downstairs_and_upstairs_are_pushforward_compatible} and \eqref{eq:unit_of_integral_transforms_downstairs_and_upstairs_are_pushforward_compatible}.
\end{remark}

\begin{theorem}
	\label{thm:descent_ascent}
	Consider \Cref{setup:fm_cube_pullback}. Let $K\in D^b_{\operatorname{coh}}(Y_1\times_S Y_2)$ be relatively perfect over each $Y_i$. Then the following are equivalent:
    \begin{enumerate}
        \item \label{thm:descent_ascent1} $\Phi_{K}$ is fully faithful (resp.\ an equivalence) on $D^b_{\operatorname{coh}}$
        \item \label{thm:descent_ascent2} $\Phi_{\mathbf{L}(t^\prime)^\ast K}$ is fully faithful (resp.\ an equivalence) on $D^b_{\operatorname{coh}}$ for any affine morphism $t$ (here, $T$ is Noetherian)
        \item \label{thm:descent_ascent3} $\Phi_{\mathbf{L}(t^\prime)^\ast K}$ is fully faithful (resp.\ an equivalence) on $D^b_{\operatorname{coh}}$ for every morphism $t\colon \operatorname{Spec}(k)\to S$ from a field
        \item \label{thm:descent_ascent3_closed} $\Phi_{\mathbf{L}(t^\prime)^\ast K}$ is fully faithful (resp.\ an equivalence) on $D^b_{\operatorname{coh}}$ for every closed point $p\in S$ with associated closed immersion $t\colon \operatorname{Spec}(\kappa(p))\to S$
        \item \label{thm:descent_ascent4} $\Phi_{\mathbf{L}(t^\prime)^\ast K}$ is fully faithful (resp.\ an equivalence) on $D^b_{\operatorname{coh}}$ for some affine surjection $t$ (here, $T$ is Noetherian).
    \end{enumerate}
\end{theorem}

\begin{proof}
    We know that $\eqref{thm:descent_ascent1}\implies \eqref{thm:descent_ascent2}$ follows from \Cref{prop:ascending}, whereas $\eqref{thm:descent_ascent2}\implies \eqref{thm:descent_ascent3}$ uses \cite[\href{https://stacks.math.columbia.edu/tag/01SI}{Tag 01SI}]{stacks-project}. Clearly, $\eqref{thm:descent_ascent3}\implies \eqref{thm:descent_ascent3_closed}$. Moreover, that $\eqref{thm:descent_ascent3_closed}\implies \eqref{thm:descent_ascent4}$ follows from \Cref{prop:descent2_closed} and taking $t=1_S$. Lastly, we check $\eqref{thm:descent_ascent4}\implies \eqref{thm:descent_ascent1}$.
    For each closed point $p\in S$, let $T_p := T\times_S \operatorname{Spec}(\kappa(p))$.
    By \Cref{prop:ascending}, we can ascend fully faithfulness or equivalences from $T$ to $T_p$.
    The morphism $T_p \to \operatorname{Spec}(\kappa(p))$ is affine and faithfully flat.
    Hence, \Cref{prop:descent1} allow us to descend fully faithfulness or equivalences from $T_p$ to $\operatorname{Spec}(\kappa(p))$.
    This shows every closed fiber has the desired property, and so \Cref{prop:descent2} implies $\Phi_K$ is fully faithful or an equivalence.
\end{proof}

%%%%%%%%%%%%%%%%%%%%%%%%%%%%%%%%%%%%%
\section{Applications}
\label{sec:applications}
%%%%%%%%%%%%%%%%%%%%%%%%%%%%%%%%%%%%%

In this section, we study the consequences of \Cref{thm:descent_ascent}. By genus of a curve, we mean arithmetic \cite[\href{https://stacks.math.columbia.edu/tag/0BY6}{Tag 0BY6}]{stacks-project}.

\begin{proof}
    [Proof of \Cref{prop:curves_isomorphic_after_base_change}]
    Let $\overline{k}$ be the algebraic closure of $k$. Suppose $K\in D^b_{\operatorname{coh}}(C\times_k C^\prime)$ has an associated integral transform $\Phi_K$ restricting to an equivalence $D^b_{\operatorname{coh}}(C) \to D^b_{\operatorname{coh}}(C^\prime)$. Denote by $\pi\colon (C\times_k C^\prime)\times_k \operatorname{Spec}(\overline{k}) \to C\times_k C^\prime$ the natural morphism. From \Cref{thm:descent_ascent}, we see that $\Phi_{\mathbf{L} \pi^\ast K}$ restricts to an equivalence $D^b_{\operatorname{coh}}(C\times_k \operatorname{Spec}(\overline{k})) \to D^b_{\operatorname{coh}}(C^\prime\times_k \operatorname{Spec}(\overline{k}))$. Then $C\times_k \operatorname{Spec}(\overline{k})$ and $C^\prime\times_k \operatorname{Spec}(\overline{k})$ are Fourier--Mukai partners over an algebraically closed field of characteristic zero. By \cite[Corollary 5.4]{Canonaco/Neeman/Stellari:2025}, we have that $\Phi_{\mathbf{L} \pi^\ast K}$ restricts to an equivalence $\operatorname{Perf}(C\times_k \operatorname{Spec}(\overline{k})) \to \operatorname{Perf}(C^\prime \times_k \operatorname{Spec}(\overline{k}))$. Consequently, we know that $C \times_k \operatorname{Spec}(\overline{k})\cong C^\prime \times_k \operatorname{Spec}(\overline{k})$ in the following cases:
    \begin{itemize}
        \item $C \times_k \operatorname{Spec}(\overline{k})$ is Gorenstein with ample or anti-ample canonical bundle \cite[Theorem 2]{Ballard:2011}
        \item $C \times_k \operatorname{Spec}(\overline{k})$ is not Gorenstein \cite[Theorem A]{Spence:2023}
        \item $C \times_k \operatorname{Spec}(\overline{k})$ is Gorenstein of genus one and trivial canonical line bundle \cite[Theorem 1.1]{LopezMartin:2014} ($C \times_k \operatorname{Spec}(\overline{k})$ is integral with genus one such that $\operatorname{deg}(\omega_X)=0$ implies $\omega_X \cong \mathcal{O}_X$ by \cite[Proposition 1.11]{Catanese:1982} where $\omega_X$ is the canonical sheaf).
        %%NOTE: Use \cite[\href{https://stacks.math.columbia.edu/tag/0H2V}{Tag 0H2V}]{stacks-project} if $\omega_X$ has positive degree as one can take tensor powers sufficiently high to make it very ample.
    \end{itemize}
    Observe these are the only possible cases for projective curves over an algebraically closed field in characteristic zero (see e.g.\ the two paragraphs above \cite[Theorem A]{Spence:2023} or use \cite[Theorem B, Proposition 1.8 and Proposition 1.11]{Catanese:1982}). 
    Hence, there is a finite field extension $L/k$ such that $C \times_k \operatorname{Spec}(L)\cong C^\prime \times_k \operatorname{Spec}(L)$ (see e.g.\ \cite[Proposition 3.2.5]{Poonen:2017}), which completes the proof. 
\end{proof}

\begin{corollary}
    \label{cor:special_fibers_genii_invariance}
    Let $S$ be a Noetherian $\mathbb{Q}$-scheme. Suppose $f_i\colon Y_i \to S$ are proper flat morphisms whose special fibers are geometrically integral curves. If $Y_1$ and $Y_2$ are Fourier--Mukai partners over $S$, then then genii of the special fibers of $Y_1$ coincide with that of $Y_2$.
\end{corollary}

\begin{proof}
    Let $s\in S$ be a closed point. Then the natural morphism
    \begin{displaymath}
        \operatorname{Spec}(\kappa(s)) \to S \to \operatorname{Spec}(\mathbb{Q})
    \end{displaymath}
    implies that $\mathbb{Q} \subseteq \kappa(s)$ (i.e.\ $\kappa(s)$ has characteristic zero). Note that $\operatorname{Spec}(\kappa(s)) \to S$ is affine (see e.g.\ \cite[\href{https://stacks.math.columbia.edu/tag/01SI}{Tag 01SI}]{stacks-project}). Denote by $Y_i^s$ for the base change of $Y_i$ along $\operatorname{Spec}(\kappa(s)) \to S$. Our hypothesis ensures that each $Y_i^s$ is a geometrically integral projective curve over $\kappa(s)$ (see e.g.\ \cite[\href{https://stacks.math.columbia.edu/tag/0A26}{Tag 0A26}]{stacks-project}). Then, by \Cref{thm:descent_ascent}, $Y^s_1$ and $Y_2^s$ are Fourier--Mukai partners over $\kappa(s)$. Furthermore, \Cref{prop:curves_isomorphic_after_base_change} guarantees that there is a finite field extension $K(s)/\kappa(s)$ such that $Y^s_1\times_{\kappa(s)} \operatorname{Spec}(K(s))$ is isomorphism to $Y^s_2\times_{\kappa(s)} \operatorname{Spec}(K(s))$. Hence, if coupled with \cite[\href{https://stacks.math.columbia.edu/tag/0BY9}{Tag 0BY9} \& \href{https://stacks.math.columbia.edu/tag/0FD2}{Tag 0FD2}]{stacks-project}, the genii of each $Y^s_i\times_{\kappa(s)} \operatorname{Spec}(K(s))$ coincides with $Y^s_j$ for $i\not=j$. In other words, the genii of $Y^s_1$ and $Y_2^s$ coincide, which completes the proof.
\end{proof}

\begin{lemma}
    \label{lem:smoothness_derived_invariance}
    Let $Y_1$ and $Y_2$ be proper schemes over a field $k$ (which need not be perfect). Suppose $Y_1$ and $Y_2$ are Fourier--Mukai partners. If $Y_1$ is geometrically regular, then so is $Y_2$. In other words, smoothness is an invariance for Fourier--Mukai partners.
\end{lemma}

\begin{proof}
    By \Cref{thm:descent_ascent}, $Y_1\times_k \operatorname{Spec}(L)$ and $Y_2 \times_k \operatorname{Spec}(L)$ are Fourier--Mukai partners for every finitely generated field extension $L/k$. So, $Y_1\times_k \operatorname{Spec}(L)$ is regular. Then \cite[Corollary 5.4]{Canonaco/Neeman/Stellari:2025} implies the following string of triangulated equivalences:
    \begin{displaymath}
        D^b_{\operatorname{coh}}(Y_1\times_k \operatorname{Spec}(L)) \cong \operatorname{Perf}(Y_1\times_k \operatorname{Spec}(L)) \cong \operatorname{Perf}(Y_2\times_k \operatorname{Spec}(L)).
    \end{displaymath}
    Hence, $\operatorname{Perf}(Y_2\times_k \operatorname{Spec}(L)) = D^b_{\operatorname{coh}}(Y_2\times_k \operatorname{Spec}(L))$, telling us that $Y_2\times_k \operatorname{Spec}(L)$ is regular. In other words, $Y_2$ must be geometrically regular.
\end{proof}

\begin{proof}
    [Proof of \Cref{prop:smooth_morphism_derived_invariance}]
    We prove only one direction as the other follows nearly verbatim. Assume $f_1$ is smooth. It suffices, by \cite[\href{https://stacks.math.columbia.edu/tag/01V8}{Tag 01V8}]{stacks-project}, to show that all the fibers $Y_2 \times_S \operatorname{Spec}(\kappa(s))$ are smooth over $\kappa(s)$ where $s\in S$ is any point. Fix a point $s\in S$ (which need not be closed). The natural morphism $t_s\colon \operatorname{Spec}(\mathcal{O}_{S,s}) \to S$ is affine (see e.g.\ \cite[\href{https://stacks.math.columbia.edu/tag/01SI}{Tag 01SI}]{stacks-project}). Hence, \Cref{thm:descent_ascent} tells us the base changes of $Y_1$ and $Y_2$ along $t_s$ are Fourier--Mukai partners. A further application \Cref{thm:descent_ascent} implies these base changes remain Fourier--Mukai partners if we base change further along the natural closed immersion $\operatorname{Spec}(\kappa(p)) \to \operatorname{Spec}(\mathcal{O}_{S,s})$. However, smoothness is stable under base change, so each base change of $Y_1$ along such morphisms remains smooth. Hence, the desired claim follows from \Cref{lem:smoothness_derived_invariance}.
\end{proof}

\begin{proof}
    [Proof of \Cref{prop:CM_or_Gorensteinness_morphism_derived_invariance}]
    It suffices to show $f_2$ is Cohen--Macaulay (resp.\ Gorenstein) if $f_1$ satisfies the same condition because the proof is analogous for the converse. In fact, we argue in a similar fashion to \Cref{prop:smooth_morphism_derived_invariance}. Moreover, by \cite[\href{https://stacks.math.columbia.edu/tag/045O}{Tag 045O}]{stacks-project} (resp.\ \cite[\href{https://stacks.math.columbia.edu/tag/0C03}{Tag 0C03}]{stacks-project}), we only need to show that all the fibers $Y_2 \times_S \operatorname{Spec}(\kappa(s))$ are Cohen--Macaulay (resp.\ Gorenstein) over an algebraic closure of each $\kappa(s)$ where $s\in S$ is any point. Fix $s\in S$ (perhaps, not closed). The natural morphism $t_s\colon \operatorname{Spec}(\mathcal{O}_{S,s}) \to S$ is affine (see e.g.\ \cite[\href{https://stacks.math.columbia.edu/tag/01SG}{Tag 01SG}]{stacks-project}). Hence, \Cref{thm:descent_ascent} tells us the base changes of $Y_1$ and $Y_2$ along $t_s$ are Fourier--Mukai partners. Then \Cref{thm:descent_ascent} implies these base changes remain Fourier--Mukai partners if we base change further along the natural closed immersion $\operatorname{Spec}(\kappa(p)) \to \operatorname{Spec}(\mathcal{O}_{S,s})$. Furthermore, if we base change further along $\overline{\kappa(s)}/\kappa(s)$ with $\overline{\kappa(s)}$ an algebraic closure of $\kappa(s)$, these base changes remain Fourier--Mukai partners via \Cref{thm:descent_ascent}. However, it follows (as fibers are geometrically integral) from \cite[Theorem 4.4]{Ruiperez/Hernandez/Martin/SanchodeSalas:2009}, that these integral projective schemes over $\overline{\kappa(s)}$ enjoy the property that  Cohen--Macaulayness (resp.\ Gorensteinness) is a derived invariance. This completes the proof.
\end{proof}

%%%%%%%%%%%%%%%%%%%%%%%%%%%%%%%%%%%%%
\begin{appendix}
\section{Gorenstein invariance}
%%%%%%%%%%%%%%%%%%%%%%%%%%%%%%%%%%%%%

Under additional assumptions, \Cref{prop:CM_or_Gorensteinness_morphism_derived_invariance} can be strengthened in the Gorenstein case. 
Its proof relies on techniques that differ from those in \Cref{prop:CM_or_Gorensteinness_morphism_derived_invariance}.
We record this for the sake of interest.

\begin{proposition}
    \label{prop:gorenstein_derived_equivalence_non_integral_case}
    Let $f_i \colon Y_i \to S$ be projective flat morphisms to a separated Noetherian scheme, where $Y_1$ and $Y_2$ are of finite Krull dimension, satisfying the resolution property (see e.g.\ \cite[\href{https://stacks.math.columbia.edu/tag/0F86}{Tag 0F86}]{stacks-project}). Assume that $Y_1$ and $Y_2$ are derived equivalent (e.g.\ the equivalence need not arise from an integral transform). Then $f_1$ is Gorenstein if, and only if, $f_2$ satisfies the same condition.
\end{proposition}

\begin{proof}
    We prove the case where $f_1$ is Gorenstein implies that of $f_2$ because the proof for the converse holds verbatim. Denote by $D^b_{\operatorname{fid}}(Y_2)$ the full triangulated subcategory in $D^b_{\operatorname{coh}}(Y_2)$ of complexes with bounded cohomology and with a bounded injective resolution (i.e.\ can be represented by a complex with finitely many nonzero components where each is injective $\mathcal{O}_X$-module). By \cite[Subchapter V.2]{Hartshorne1966}, the dualizing sheaf $f_2^!\mathcal{O}_S$ is contained in $D^b_{\operatorname{fid}}(Y_2)$. Furthermore,
    \cite[\href{https://stacks.math.columbia.edu/tag/0A89}{Tag 0A89}]{stacks-project} tells us that there is an equivalence
    \begin{equation}
        \label{eq:duality}
        D:=\mathbf{R}\operatorname{\mathcal{H}\! \mathit{om}}(-, f_2^!\mathcal{O}_S ) \colon D^b_{\operatorname{coh}} (Y_2) \xrightarrow{\sim} D^b_{\operatorname{coh}} (Y_2)^{\operatorname{op}}
    \end{equation}
    satisfying $D\circ D = \operatorname{id}$.
    Moreover, this equivalence restricts to an equivalence \cite[Proposition V.2.6 (a)]{Hartshorne1966},
    \begin{equation}
        \label{eq:duality-perf}
        D  \colon \operatorname{Perf} (Y_2) \xrightarrow{\sim} D^b_{\operatorname{fid}} (Y_2)^{\operatorname{op}}
    \end{equation}
    So, we have the following intrinsic characterization of $D^b_{\operatorname{fid}}(Y_1)$ in $D^b_{\operatorname{coh}}(Y_1)$,
    \begin{equation}
        \label{eq:charact-fid}
            D^b_{\operatorname{fid}}(Y_1) = \{ B \in D^b_{\operatorname{coh}}(Y_1) \mid  \operatorname{Hom} ( A , B [p]  ) = 0 , \text{for}\ 0\ll |p|\text{ and all } A\in D^b_{\operatorname{coh}}(Y_1)\} .
    \end{equation}
    Indeed, by \cite[Proposition 1.11]{Orlov_2006}, we have a dual characterization for perfect complexes in $D^b_{\operatorname{coh}}(Y_1)$,
    \begin{displaymath}
        \operatorname{Perf}(Y_1) = \{ B \in D^b_{\operatorname{coh}}(Y_1) \mid \operatorname{Hom} ( B , A [p]  ) = 0 , \text{for}\ 0\ll |p|\text{ and all } A\in D^b_{\operatorname{coh}}(Y_1) \} ,
    \end{displaymath}
    because $Y_1$ is separated Noetherian of finite Krull dimension and satisfies the resolution property (Orlov calls this (ELF) scheme - note that the proof in loc. cit. does not use that the scheme is defined over a field). This characterization of $\operatorname{Perf}$, together with the equivalences \eqref{eq:duality} and \eqref{eq:duality-perf}, implies \eqref{eq:charact-fid}.

    Recall that $Y_1$ over $S$ is Gorenstein if, and only if, the dualizing sheaf $f_1^!\mathcal{O}_S$ is invertible (see e.g.\ \cite[Exercise V.9.7]{Hartshorne1966}). In particular, $f_1^!\mathcal{O}_S$ is a perfect complex. If now $\Phi \colon D_{\operatorname{qc}}(Y_1) \to D_{\operatorname{qc}}(Y_2) $ is an equivalence, then $\Phi (f_1^!\mathcal{O}_S)$ is a perfect complex in $D_{\operatorname{qc}}(Y_2)$. 
    Furthermore, as $\Phi$ restricts to an equivalence on $D^b_{\operatorname{coh}}$, we see by the characterization \eqref{eq:charact-fid} that $\Phi (f_1^!\mathcal{O}_S)$ is contained in $D^b_{\operatorname{fid}} (Y_2 )$.
    Hence, by \eqref{eq:duality-perf}, we see that $\Phi (f_1^!\mathcal{O}_S)$ is of the form
    \begin{displaymath}
        \Phi (f_1^!\mathcal{O}_S) \cong D ( P ) = \mathbf{R}\operatorname{\mathcal{H}\! \mathit{om}} ( P , f_2^!\mathcal{O}_S) \cong P^\vee \otimes^{\mathbf{L}} f_2^!\mathcal{O}_S \in D^b_{\operatorname{fid}} (Y_2 )
    \end{displaymath}
    for some $P$ in $\operatorname{Perf}(Y_2)$. 
    %However, $\Phi (f_1^!\mathcal{O}_S)$ is perfect, and so, $P^\vee \otimes^{\mathbf{L}} f_2^!\mathcal{O}_S$ is in $\operatorname{Perf}(Y_2)$. 
    
    By the description of \eqref{eq:charact-fid}, the image of perfects under $\Phi$ belongs to $D^b_{\operatorname{fid}} (Y_2)$. 
    However, since $D^b_{\operatorname{fid}} (Y_1)$ coincides with $\operatorname{Perf}(Y_1)$ (because $f_1^! \mathcal{O}_S$ is invertible), and since $\Phi$ also preserves $\operatorname{Perf}$, we have that $D^b_{\operatorname{fid}} (Y_2)$ coincides with $\operatorname{Perf}(Y_2)$.
    In particular, $f_2^! \mathcal{O}_S$, which lives naturally in $D^b_{\operatorname{fid}} (Y_2)$, is perfect.
    Moreover, since $f_2^!\mathcal{O}_S^\vee \otimes^{\mathbf{L}} f_2^!\mathcal{O}_S \cong \mathbf{R}\operatorname{\mathcal{H}\! \mathit{om}}(f_2^!\mathcal{O}_S , f_2^!\mathcal{O}_S ) \cong \mathcal{O}_{Y_2} $ by the definition of a dualizing complex, it is easy to deduce that $f_2^!\mathcal{O}_S$ is a shift of a line bundle of $Y_2$ (see e.g.\ \cite[Lemma V.3.3]{Hartshorne1966}). Again, by \cite[Exercise V.9.7]{Hartshorne1966}, we conclude that $Y_2$ over $S$ is Gorenstein, which completes the proof.
\end{proof}

%%%%%%%%%%%%%%%%%%%%%%%%%%%%%%%%%%%%%
\end{appendix}
%%%%%%%%%%%%%%%%%%%%%%%%%%%%%%%%%%%%%

\bibliographystyle{alpha}
\bibliography{mainbib}

\end{document}
